% Lesson 25 — Writing: Compositions about developing and sharing interests in a trade/job/profession — Anglais 1ère A — Cours signé OnBuch+
% Programme officiel MINESEC. Compiler avec : tectonic ENA25-writing-compositions-about-developing-and-sharing-inter.tex
\newcommand{\DOCMATIERE}{Anglais}
\newcommand{\DOCNIVEAU}{1ère A}
\newcommand{\DOCMODULE}{Chapter 2 : Economic Life and Occupations}
\newcommand{\DOCLECON}{ENA25}
\newcommand{\DOCTITRE}{Lesson 25 — Writing: Compositions about developing and sharing interests in a trade/job/profession}
% OnBuch+ — préambule « cours signés » ENRICHI (compilé avec tectonic / XeLaTeX)
% Système visuel « Pop » d'OnBuch+ : fond crème (jamais blanc), bordures encre
% épaisses, ombres « dures » décalées, titres Archivo Black en capitales.
% Version MATHÉMATIQUES : graphes (pgfplots/tikz), boîtes pédagogiques
% (définition, propriété, méthode, attention, le savais-tu, exercice résolu,
% exercice, TP, bilan), page de garde et signature OnBuch+ ancrée en bas.
%
% Macros attendues dans main.tex AVANT \input{preamble} :
%   \DOCMATIERE  \DOCNIVEAU  \DOCMODULE  \DOCLECON (numéro)  \DOCTITRE
\documentclass[11pt]{article}

\usepackage{fontspec}
\usepackage[english,french]{babel} % français = langue principale ; l'anglais via \eng{..} / textebox
\usepackage[a4paper,margin=2cm,top=3.4cm,bottom=2.8cm,headheight=1.4cm,footskip=1.3cm]{geometry}
\usepackage{amsmath,amssymb,amsfonts}
\usepackage{mathtools} % \xmapsto, \coloneqq... souvent utilisés par les modèles
\usepackage{cancel} % simplifications barrées (bilans de forces)
% \abs{z} (module, valeur absolue) et \norm{u} (norme), utilisés par les modèles
\providecommand{\abs}{}\let\abs\relax\DeclarePairedDelimiter{\abs}{\lvert}{\rvert}
\providecommand{\norm}{}\let\norm\relax\DeclarePairedDelimiter{\norm}{\lVert}{\rVert}
\usepackage[table]{xcolor} % [table] : fournit aussi \rowcolors (alternance de lignes)
\usepackage{pagecolor}
\usepackage{tikz}
\usetikzlibrary{arrows.meta,positioning,calc,shapes.geometric,shapes.misc,shapes.arrows,decorations.pathmorphing,decorations.pathreplacing,decorations.markings,patterns,shadows,mindmap,trees,fit,backgrounds,matrix,intersections,angles,quotes}
\usepackage{pgfplots}
\usepackage[version=4]{mhchem} % équations nucléaires \ce{^{235}_{92}U}
\usepackage[european,siunitx]{circuitikz} % schémas électriques
\pgfplotsset{compat=1.18}
\usepgfplotslibrary{fillbetween,groupplots} % aires entre courbes (calcul intégral)
\usepackage{enumitem}
\usepackage{fancyhdr}
% [most] charge toutes les bibliothèques tcolorbox (listings, minted, raster,
% documentation...) et gonfle inutilement la mémoire TeX sur un document long
% (~300 boîtes) : on ne charge que ce qui est réellement utilisé.
\usepackage[skins,breakable]{tcolorbox}
\usepackage{graphicx}
\usepackage{colortbl}
\usepackage{booktabs}
\usepackage{tabularx}
\usepackage{diagbox} % case d'en-tête diagonale des tableaux à double entrée
% Colonnes extensibles centrée / gauche / droite (C, L, R), souvent utilisées par les modèles
\newcolumntype{C}{>{\centering\arraybackslash}X}
\newcolumntype{L}{>{\raggedright\arraybackslash}X}
\newcolumntype{R}{>{\raggedleft\arraybackslash}X}
\usepackage{array}
\usepackage{multicol}
\usepackage{multirow}
\usepackage{siunitx}
% parse-numbers=false : accepte aussi \SI{2,5 \times 10^{-3}}{...}
\sisetup{locale=FR,output-decimal-marker={,},group-minimum-digits=4,parse-numbers=false}
\DeclareSIUnit\ohm{\ensuremath{\symup{\Omega}}} % U+2126 absent de la police
\DeclareSIUnit\division{div} % réglages d'oscilloscope (ms/div, V/div)
\DeclareSIUnit\tr{tr} % tours (vitesse de rotation, ex. \SI{1500}{\tr\per\minute})
\DeclareSIUnit\molar{M} % concentration molaire (ex. \SI{3}{\molar}, \SI{1}{\micro\molar})
\DeclareSIUnit\an{an} % année (le modèle confond parfois avec \year, un compteur TeX) — ex. \SI{5}{\kilo\meter\per\an}
\DeclareSIUnit\FCFA{FCFA} % franc CFA (ex. \SI{500}{\FCFA\per\kilo\gram})
\DeclareSIUnit\degreeBrix{\textdegree Brix} % teneur en sucre d'un sirop/jus (réfractométrie)
\DeclareSIUnit\month{mois} % le modèle utilise parfois \month, un compteur TeX, comme unité de durée
\usepackage{qrcode}
% \degree : commande fréquemment utilisée par les modèles pour un angle ou une
% température en dehors de \SI{}{} (non fournie par les paquets chargés).
\providecommand{\degree}{^{\circ}}
% \qed (fin de démonstration), utilisé par les modèles sans amsthm
\providecommand{\qed}{\hfill$\square$}
% \barre : double barre des valeurs interdites dans un tableau de variations (tkz-tab)
\providecommand{\barre}{\ensuremath{\|}}
% environnement proof (amsthm non chargé) utilisé par les modèles
\ifdefined\proof\else\newenvironment{proof}[1][Démonstration]{\par\noindent\textit{#1.}\ }{\qed\par}\fi
\usepackage{lastpage}
\usepackage{hyperref}
\hypersetup{hidelinks}

% ============================================================
% Palette « Pop » OnBuch+ — mêmes hex que la classe P (plus_ui.dart)
% ============================================================
\definecolor{poporange}{HTML}{F59321}
\definecolor{popdark}{HTML}{E07A0C}
\definecolor{popcream}{HTML}{FFF6E8}
\definecolor{popink}{HTML}{1C1714}
\definecolor{poppurple}{HTML}{7A5AE0}
\definecolor{popgreen}{HTML}{2E9E6A}
\definecolor{poppink}{HTML}{E0457B}
\definecolor{popblue}{HTML}{2F7DE1}
\definecolor{popgold}{HTML}{C98A12}
\definecolor{popmuted}{HTML}{8A7F76}
\definecolor{popline}{HTML}{EADFD2}
\definecolor{popgray}{HTML}{8A7F76}
\colorlet{popgrey}{popgray}
% Alias tolérants (le modèle nomme parfois les couleurs différemment)
\colorlet{popwhite}{white}
\colorlet{popblack}{popink}
\colorlet{popred}{poppink}
\colorlet{popyellow}{popgold}
% Teintes claires (fonds de tableaux/encadrés) — le modèle nomme parfois
% les couleurs avec un suffixe L (ex. popblueL) sans les avoir définies.
\colorlet{poporangeL}{poporange!20}
\colorlet{popdarkL}{popdark!20}
\colorlet{poppurpleL}{poppurple!20}
\colorlet{popgreenL}{popgreen!20}
\colorlet{poppinkL}{poppink!20}
\colorlet{popblueL}{popblue!20}
\colorlet{popgoldL}{popgold!20}
\colorlet{popredL}{popred!20}
\colorlet{popgrayL}{popgray!20}
% Teintes claires pour fonds de tableaux / zones
\colorlet{popblueL}{popblue!10!white}
\colorlet{poporangeL}{poporange!14!white}
\colorlet{popgreenL}{popgreen!12!white}
\colorlet{poppurpleL}{poppurple!12!white}
\colorlet{poppinkL}{poppink!12!white}
\colorlet{popgoldL}{popgold!14!white}

\pagecolor{popcream}

% ============================================================
% Typographie
% ============================================================
\defaultfontfeatures{Ligatures=TeX}
\setmainfont{PlusJakartaSans-Medium.ttf}[Path=../../../fonts/,BoldFont=PlusJakartaSans-Bold.ttf]
% Symboles, API et emojis absents de la police principale : repli sur DejaVu Sans (caractères actifs)
\newfontface\symfont{DejaVuSans.ttf}[Path=../../../fonts/]
\makeatletter
\newcommand\symactive[1]{\catcode#1=13 \begingroup\lccode`\~=#1 \lowercase{\endgroup\def~}{{\symfont\char#1}}}
\newcommand\symblank[1]{\catcode#1=13 \begingroup\lccode`\~=#1 \lowercase{\endgroup\def~}{}}
\newcommand\symhyph[1]{\catcode#1=13 \begingroup\lccode`\~=#1 \lowercase{\endgroup\def~}{\mbox{-}}}
\newcount\symc
\newcommand\symrange[2]{\symc=#1 \@whilenum\symc<#2 \do{\expandafter\symactive\expandafter{\the\symc}\advance\symc by 1 }}
\newcommand\symblankrange[2]{\symc=#1 \@whilenum\symc<#2 \do{\expandafter\symblank\expandafter{\the\symc}\advance\symc by 1 }}
\makeatother
\symrange{"0250}{"02FF}\symactive{"2713}\symactive{"2714}\symactive{"2717}\symactive{"2718}\symactive{"2610}\symactive{"2611}\symactive{"2612}
\symhyph{"2011}\symhyph{"2E1A}\symblankrange{"1F300}{"1FAFF}\symblank{"FE0F}
% Maths en sans-serif (Fira Math) avec chiffres et lettres droites de Plus
% Jakarta, pour que les formules se fondent dans le texte.
\usepackage[math-style=ISO,bold-style=ISO]{unicode-math}
\setmathfont{FiraMath-Regular.otf}[Path=../../../fonts/]
\setmathfont{PlusJakartaSans-Medium.ttf}[Path=../../../fonts/,range={up/num,up/{Latin,latin},\mathup}]
\usepackage{icomma} % virgule décimale sans espace en mode math
% Caractères mathématiques Unicode absents des polices de texte (Plus Jakarta,
% Archivo Black) : rendus par leur équivalent mathématique, en texte comme en
% formule — sinon ils s'affichent en carré vide (ex. « Arithmétique dans ℤ »).
\usepackage{newunicodechar}
\newunicodechar{ℝ}{\ensuremath{\mathbb{R}}}
\newunicodechar{ℤ}{\ensuremath{\mathbb{Z}}}
\newunicodechar{ℂ}{\ensuremath{\mathbb{C}}}
\newunicodechar{ℕ}{\ensuremath{\mathbb{N}}}
\newunicodechar{ℚ}{\ensuremath{\mathbb{Q}}}
\newunicodechar{𝔻}{\ensuremath{\mathbb{D}}}
\newunicodechar{α}{\ensuremath{\alpha}}
\newunicodechar{β}{\ensuremath{\beta}}
\newunicodechar{γ}{\ensuremath{\gamma}}
\newunicodechar{δ}{\ensuremath{\delta}}
\newunicodechar{ε}{\ensuremath{\varepsilon}}
\newunicodechar{θ}{\ensuremath{\theta}}
\newunicodechar{λ}{\ensuremath{\lambda}}
\newunicodechar{μ}{\ensuremath{\mu}}
\newunicodechar{σ}{\ensuremath{\sigma}}
\newunicodechar{τ}{\ensuremath{\tau}}
\newunicodechar{φ}{\ensuremath{\varphi}}
\newunicodechar{ψ}{\ensuremath{\psi}}
\newunicodechar{ω}{\ensuremath{\omega}}
\newunicodechar{Σ}{\ensuremath{\Sigma}}
% majuscules produites par \MakeUppercase dans les titres (α-aminés -> Α)
\newunicodechar{Α}{\ensuremath{\alpha}}
\newunicodechar{Β}{\ensuremath{\beta}}
\newunicodechar{Γ}{\ensuremath{\Gamma}}
\newunicodechar{Δ}{\ensuremath{\Delta}}
\newunicodechar{Ε}{\ensuremath{\varepsilon}}
\newunicodechar{Θ}{\ensuremath{\Theta}}
\newunicodechar{Λ}{\ensuremath{\Lambda}}
\newunicodechar{Μ}{\ensuremath{\mu}}
\newunicodechar{Τ}{\ensuremath{\tau}}
\newunicodechar{Φ}{\ensuremath{\Phi}}
\newunicodechar{Ψ}{\ensuremath{\Psi}}
\newunicodechar{Ω}{\ensuremath{\Omega}}
\newunicodechar{↦}{\ensuremath{\mapsto}}
\newunicodechar{∈}{\ensuremath{\in}}
\newunicodechar{∉}{\ensuremath{\notin}}
\newunicodechar{∧}{\ensuremath{\wedge}}
\newunicodechar{⇔}{\ensuremath{\Leftrightarrow}}
\newunicodechar{⟺}{\ensuremath{\Longleftrightarrow}}
\newunicodechar{⇒}{\ensuremath{\Rightarrow}}
\newunicodechar{⟹}{\ensuremath{\Longrightarrow}}
\newunicodechar{∘}{\ensuremath{\circ}}
\newunicodechar{∪}{\ensuremath{\cup}}
\newunicodechar{∩}{\ensuremath{\cap}}
\newunicodechar{≡}{\ensuremath{\equiv}}
\newunicodechar{⊂}{\ensuremath{\subset}}
\newunicodechar{⊥}{\ensuremath{\perp}}
\newunicodechar{∃}{\ensuremath{\exists}}
\newunicodechar{∀}{\ensuremath{\forall}}
\newunicodechar{∛}{\ensuremath{\sqrt[3]{\,}}}
\newunicodechar{∜}{\ensuremath{\sqrt[4]{\,}}}
\newunicodechar{⃗}{\ensuremath{{}^{\to}}}
\newunicodechar{ⁿ}{\ifmmode^{n}\else\textsuperscript{\ensuremath{n}}\fi}
\newunicodechar{ˣ}{\ifmmode^{x}\else\textsuperscript{\ensuremath{x}}\fi}
\newunicodechar{ᵇ}{\ifmmode^{b}\else\textsuperscript{\ensuremath{b}}\fi}
\newunicodechar{ʳ}{\ifmmode^{r}\else\textsuperscript{\ensuremath{r}}\fi}
\newunicodechar{ᵘ}{\ifmmode^{u}\else\textsuperscript{\ensuremath{u}}\fi}
\newunicodechar{ʸ}{\ifmmode^{y}\else\textsuperscript{\ensuremath{y}}\fi}
\newunicodechar{ᵃ}{\ifmmode^{a}\else\textsuperscript{\ensuremath{a}}\fi}
\newunicodechar{ᵉ}{\ifmmode^{e}\else\textsuperscript{\ensuremath{e}}\fi}
\newunicodechar{ᵏ}{\ifmmode^{k}\else\textsuperscript{\ensuremath{k}}\fi}
\newunicodechar{ᵖ}{\ifmmode^{p}\else\textsuperscript{\ensuremath{p}}\fi}
\newunicodechar{ᴺ}{\ifmmode^{N}\else\textsuperscript{\ensuremath{N}}\fi}
\newunicodechar{ᶜ}{\ifmmode^{c}\else\textsuperscript{\ensuremath{c}}\fi}
\newunicodechar{ʰ}{\ifmmode^{h}\else\textsuperscript{\ensuremath{h}}\fi}
\newunicodechar{ᵠ}{\ifmmode^{q}\else\textsuperscript{\ensuremath{q}}\fi}
\newunicodechar{ₙ}{\ifmmode_{n}\else\textsubscript{\ensuremath{n}}\fi}
\newunicodechar{ₐ}{\ifmmode_{a}\else\textsubscript{\ensuremath{a}}\fi}
\newunicodechar{ᵢ}{\ifmmode_{i}\else\textsubscript{\ensuremath{i}}\fi}
\newunicodechar{ᵣ}{\ifmmode_{r}\else\textsubscript{\ensuremath{r}}\fi}
\newunicodechar{ₖ}{\ifmmode_{k}\else\textsubscript{\ensuremath{k}}\fi}
\newunicodechar{ᵦ}{\ifmmode_{\beta}\else\textsubscript{\ensuremath{\beta}}\fi}
\newunicodechar{ₓ}{\ifmmode_{x}\else\textsubscript{\ensuremath{x}}\fi}
\newfontfamily\popdisplay{ArchivoBlack-Regular.ttf}[Path=../../../fonts/]
\newfontfamily\poplogo{SpaceGrotesk-Bold.ttf}[Path=../../../fonts/]
\newfontfamily\popsemi{PlusJakartaSans-SemiBold.ttf}[Path=../../../fonts/]
\newfontfamily\popxbold{PlusJakartaSans-ExtraBold.ttf}[Path=../../../fonts/]
\newfontfamily\popxboldls{PlusJakartaSans-ExtraBold.ttf}[Path=../../../fonts/,LetterSpace=3.0]

\setlist[enumerate]{leftmargin=1.7em,itemsep=5pt,topsep=4pt}
\setlist[itemize]{leftmargin=1.7em,itemsep=4pt,topsep=4pt,label={\color{poporange}\textbullet}}
\setlength{\parindent}{0pt}
\setlength{\parskip}{5pt}
\renewcommand{\arraystretch}{1.25}

% pgfplots : style Pop par défaut
\pgfplotsset{
  every axis/.append style={axis line style={popink,line width=1pt},tick style={popink},
    grid style={popline,line width=0.5pt},label style={font=\small},tick label style={font=\footnotesize}},
  popcurve/.style={poporange,line width=1.8pt,no markers,smooth},
}
% Même style disponible dans un \draw TikZ simple (hors pgfplots)
\tikzset{popcurve/.style={poporange,line width=1.8pt}}
\tikzset{popgrid/.style={popline,very thin}}
\colorlet{popcurve}{poporange} % le nom du style est parfois utilisé comme couleur

% ============================================================
% Pill / squiggle
% ============================================================
\newcommand{\poppill}[3]{%
  \tikz[baseline=-0.55ex]\node[draw=popink,line width=1.2pt,rounded corners=2.6pt,%
    fill=#2,inner xsep=6.5pt,inner ysep=3pt]{%
    \popxboldls\fontsize{7}{7}\selectfont\color{#3}\MakeUppercase{#1}};%
}
\newcommand{\popsquiggle}[1]{%
  \tikz[baseline=-0.4ex]{
    \draw[#1,line width=2.6pt,line cap=round]
      plot [smooth] coordinates {(0,0.09) (0.34,0.01) (0.68,0.21) (1.02,0.01) (1.36,0.10)};
  }%
}
% Mot-clé surligné (marqueur orange sous le texte)
\newcommand{\cle}[1]{{\popxbold\color{popdark}#1}}

% ============================================================
% En-tête / pied
% ============================================================
\pagestyle{fancy}
\fancyhf{}
\renewcommand{\headrulewidth}{0pt}
\renewcommand{\footrulewidth}{0pt}
\lhead{\raisebox{-0.15ex}{\poplogo\fontsize{13}{13}\selectfont\color{popink}OnBuch\color{poporange}+}}
\rhead{\poppill{\DOCMATIERE\ \textbullet\ \DOCNIVEAU\ \textbullet\ Leçon \DOCLECON}{white}{popink}}
\lfoot{\popxbold\fontsize{7.6}{7.6}\selectfont\color{popmuted}\MakeUppercase{Cours signé OnBuch+}\ \textbullet\ {\popsemi\DOCTITRE}}
\rfoot{\tikz[baseline=-0.6ex]\node[rounded corners=8.5pt,fill=popink,minimum height=17pt,minimum width=17pt,inner xsep=5pt,inner sep=0]{\popxbold\fontsize{8}{8}\selectfont\color{popcream}\thepage\,/\,\pageref*{LastPage}};}

% ============================================================
% Titres
% ============================================================
\newcommand{\chaptitre}[2]{%
  \begin{tcolorbox}[enhanced,colback=poporange,colframe=popink,boxrule=2.6pt,arc=11pt,
      drop shadow=popink,left=15pt,right=15pt,top=12pt,bottom=11pt,before skip=3pt,after skip=18pt]
    \poppill{#2}{popink}{popcream}\par\vspace{9pt}
    {\raggedright\hyphenpenalty=10000\exhyphenpenalty=10000\popdisplay\fontsize{22}{24}\selectfont\color{popcream}\MakeUppercase{#1}\par}
  \end{tcolorbox}
}
% Section numérotée : pastille orange avec le numéro + titre Archivo
\newcounter{popsec}
\newcommand{\coursec}[1]{%
  \stepcounter{popsec}\setcounter{popsub}{0}%
  \par\vspace{16pt}\noindent
  \begin{minipage}{\linewidth}
  \tikz[baseline=-0.7ex]\node[draw=popink,line width=1.6pt,fill=poporange,rounded corners=4pt,minimum size=22pt,inner sep=0]{\popdisplay\fontsize{12}{12}\selectfont\color{popcream}\thepopsec};%
  \hspace{8pt}{\popdisplay\fontsize{15}{17}\selectfont\color{popink}\MakeUppercase{#1}}\par
  \vspace{4pt}\noindent{\color{poporange}\rule{36pt}{3.6pt}}
  \end{minipage}\par\nopagebreak\vspace{8pt}
}
\newcounter{popsub}
\newcommand{\courssub}[1]{%
  \stepcounter{popsub}%
  \par\vspace{9pt}\noindent{\popxbold\fontsize{11.5}{13}\selectfont\color{popink}{\color{poporange}\thepopsec.\thepopsub}\hspace{6pt}#1}\par\nopagebreak\vspace{4pt}
}

% ============================================================
% Boîtes pédagogiques — même recette : fond blanc, bordure épaisse colorée,
% ombre dure encre, badge plein. Titre optionnel [#1].
% ============================================================
\tcbset{popbase/.style={breakable,enhanced,colback=white,boxrule=2.3pt,arc=9pt,
  shadow={3pt}{-3pt}{0pt}{popink},left=14pt,right=14pt,top=11pt,bottom=11pt,before skip=11pt,after skip=15pt,
  fontupper=\popsemi\fontsize{10.6}{14.5}\selectfont\color{popink},
  fontlower=\popsemi\fontsize{10.6}{14.5}\selectfont\color{popink}}}
% Coche dessinée (le glyphe ✓ n'existe pas dans les polices du document)
\newcommand{\popcheck}[1]{\tikz[baseline=-0.5ex]\draw[#1,line width=1.6pt,line cap=round,line join=round](-0.13,0.0)--(-0.04,-0.09)--(0.13,0.1);}
\providecommand{\coche}{\popcheck{popgreen}}
\providecommand{\resultat}[1]{\par\smallskip\noindent\fcolorbox{popgreen}{popgreenL}{\parbox{\dimexpr\linewidth-2\fboxsep-2\fboxrule\relax}{\textbf{Résultat.} #1}}\par\smallskip}
\newcommand{\popboxhead}[3]{\poppill{#1}{#2}{white}\ifx\relax#3\relax\else\hspace{6pt}{\popxbold\fontsize{10.6}{12}\selectfont\color{popink}#3}\fi\par\vspace{7pt}}

\newtcolorbox{aretenir}[1][]{popbase,colframe=popgreen,before upper={\popboxhead{À retenir}{popgreen}{#1}}}
% Texte en anglais (document de lecture, dialogue, extrait) : cadre
% d'étude. Un titre optionnel (source, auteur) s'affiche dans l'en-tête.
\newtcolorbox{textebox}[1][]{popbase,colframe=popblue,before upper={\popboxhead{Text}{popblue}{#1}\begin{otherlanguage}{english}},after upper={\end{otherlanguage}}}
% Marques ✓ / ✗ (forme correcte / fautive) souvent utilisées par les modèles
\providecommand{\cross}{\textcolor{poppink}{\ensuremath{\times}}}
\providecommand{\xmark}{\cross}\providecommand{\crossmark}{\cross}\providecommand{\cmark}{\ensuremath{\checkmark}}
% Ligne de réponse / justification à compléter (inventée par les modèles)
\providecommand{\justification}[1]{\par\noindent\textit{Justification~:} \dotfill\par}
% \textipa (package tipa non chargé) : la police ne couvre pas l'API -> simple passage
\providecommand{\textipa}[1]{#1}
% Soulignement (ulem non chargé) : \uline, \ul
\providecommand{\uline}[1]{\underline{#1}}\providecommand{\ul}[1]{\underline{#1}}
% \glossaire{\item ...} inventée par les modèles : liste de glossaire
\providecommand{\glossaire}[1]{\begin{itemize}#1\end{itemize}}
% \alert (beamer) inventée par les modèles : texte mis en évidence
\providecommand{\alert}[1]{\textbf{\textcolor{poppink}{#1}}}
% variantes ASCII d'ordinaux écrites en macro
\providecommand{\iere}{\textsuperscript{re}}\providecommand{\ieme}{\textsuperscript{e}}\providecommand{\ere}{\textsuperscript{re}}\providecommand{\eme}{\textsuperscript{e}}\providecommand{\er}{\textsuperscript{er}}\providecommand{\ie}{\textsuperscript{e}}
% Ordinaux français écrits en macro par les modèles : 1\ière, 3\ième
\newcommand{\ière}{\textsuperscript{re}}\newcommand{\ième}{\textsuperscript{e}}
% \exemple (inventée par les modèles) : étiquette « Exemple : »
\providecommand{\exemple}{\textbf{Exemple~:}\ }
% \square utilisable aussi hors mode math (cases à cocher)
\AtBeginDocument{\let\squareMath\square\renewcommand{\square}{\ensuremath{\squareMath}}}
% Mot ou phrase en anglais à l'intérieur d'un paragraphe français
\newcommand{\eng}[1]{\ifmmode\text{\textit{#1}}\else\begingroup\selectlanguage{english}\itshape #1\endgroup\fi}
\let\esp\eng % alias toléré
% flèches d'intonation inventées par les modèles
\providecommand{\textsearrow}{}\renewcommand{\textsearrow}{\ensuremath{\searrow}}\providecommand{\textnearrow}{}\renewcommand{\textnearrow}{\ensuremath{\nearrow}}
% \fr{..} : mot français (inventé par les modèles)
\providecommand{\fr}[1]{\textit{#1}}
\newtcolorbox{exemplebox}[1][]{popbase,colframe=poporange,before upper={\popboxhead{Exemple}{poporange}{#1}}}
\newtcolorbox{definition}[1][]{popbase,colframe=popblue,before upper={\popboxhead{Définition}{popblue}{#1}}}
\newtcolorbox{propriete}[1][]{popbase,colframe=poppurple,before upper={\popboxhead{Propriété}{poppurple}{#1}}}
\newtcolorbox{methode}[1][]{popbase,colframe=popgold,before upper={\popboxhead{Méthode}{popgold}{#1}}}
\newtcolorbox{attention}[1][]{popbase,colframe=poppink,before upper={\popboxhead{Attention}{poppink}{#1}}}
\newtcolorbox{experience}[1][]{popbase,colframe=popblue,colback=popblueL,before upper={\popboxhead{Expérience}{popblue}{#1}}}
% « Le savais-tu ? » : carte violette pleine (contexte, culture, Cameroun)
\newtcolorbox{savaistu}[1][]{popbase,colback=poppurple,colframe=popink,
  fontupper=\popsemi\fontsize{10.4}{14.2}\selectfont\color{white},
  before upper={\poppill{Le savais-tu\,?}{white}{poppurple}\ifx\relax#1\relax\else\hspace{6pt}{\popxbold\color{white}#1}\fi\par\vspace{7pt}}}
% Exercice résolu : énoncé en haut, \tcblower puis la solution sur fond crème
\newcounter{popexo}\newcounter{popexr}
\newtcolorbox{exoresolu}[1][]{popbase,colframe=poporange,colbacklower=popcream,
  segmentation style={popink,line width=1.2pt,dashed},
  before upper={\stepcounter{popexr}\popboxhead{Exercice résolu \thepopexr}{poporange}{#1}},
  before lower={\poppill{Solution}{popink}{popcream}\par\vspace{6pt}}}
% Exercice (non résolu en place) — la correction est dans la partie Corrigés
\newtcolorbox{exercice}[1][]{popbase,colframe=popink,
  before upper={\stepcounter{popexo}\popboxhead{Exercice \thepopexo}{popink}{#1}}}
% Corrigé
\newtcolorbox{corrige}[1][]{popbase,colframe=popgreen,colback=popgreenL,
  before upper={\popboxhead{Corrigé}{popgreen}{#1}}}
% Objectifs / prérequis / bilan
\newtcolorbox{objectifs}{popbase,colframe=popink,colback=poporangeL,
  before upper={\popboxhead{Objectifs de la leçon}{popink}{}}}
\newtcolorbox{prerequis}{popbase,colframe=popink,colback=popblueL,
  before upper={\popboxhead{Prérequis}{popblue}{}}}
\newtcolorbox{bilan}{popbase,colframe=popink,colback=white,boxrule=2.8pt,
  before upper={\popboxhead{Fiche bilan}{poporange}{L'essentiel à connaître pour l'examen}}}
% Figure encadrée avec légende
\newtcolorbox{popfigure}[1][]{enhanced,breakable=false,colback=white,colframe=popline,boxrule=1.4pt,arc=8pt,
  left=10pt,right=10pt,top=10pt,bottom=8pt,before skip=10pt,after skip=12pt,halign=center,
  lower separated=false,halign lower=center,
  fontlower=\popsemi\fontsize{9}{11.5}\selectfont\color{popmuted},
  #1}
\newcommand{\legende}[1]{\par\vspace{6pt}{\popsemi\fontsize{9}{11.5}\selectfont\color{popmuted}#1}}

% ============================================================
% Signature OnBuch+
% ============================================================
\newcommand{\poplogoinline}[1]{{\poplogo\fontsize{#1}{#1}\selectfont\color{popink}OnBuch\color{poporange}+}}
\newcommand{\signatureonbuch}{%
  \begin{tcolorbox}[enhanced,colback=popink,colframe=popink,boxrule=2.2pt,arc=10pt,
      drop shadow=poporange,left=14pt,right=14pt,top=10pt,bottom=10pt,before skip=0pt,after skip=0pt]
    \tikz[baseline=-0.7ex]\node[circle,fill=popgreen,draw=popcream,line width=1.3pt,minimum size=22pt,inner sep=0]{\popcheck{white}};%
    \hspace{9pt}\begin{minipage}[c]{0.62\linewidth}
      {\popxbold\fontsize{12}{13}\selectfont\color{popcream}Signé\ \textbullet\ {\poplogo OnBuch\color{poporange}+}}\par\vspace{2pt}
      {\raggedright\popsemi\fontsize{8}{10.5}\selectfont\color{popline}\MakeUppercase{Équipe pédagogique OnBuch+}\\\MakeUppercase{Conforme au programme officiel MINESEC}\par}
    \end{minipage}\hfill
    {\poplogo\fontsize{22}{22}\selectfont\color{popcream}OnBuch\color{poporange}+}
  \end{tcolorbox}%
}

% ============================================================
% Page de garde
% #1 titre · #2 matière · #3 niveau · #4 module · #5 n° leçon · #6 durée ·
% #7 description / objectifs (texte ou itemize)
% ============================================================
\newcommand{\pagegarde}[7]{%
  \thispagestyle{empty}
  \newgeometry{a4paper,margin=1.7cm,top=1.6cm,bottom=1.4cm}%
  \begin{tikzpicture}[remember picture,overlay]
    \fill[poporange!18] ([xshift=-3.2cm,yshift=-3.6cm]current page.north east) circle (4.6cm);
    \fill[poppurple!14] ([xshift=2.4cm,yshift=7.6cm]current page.south west) circle (3.8cm);
  \end{tikzpicture}%
  {\poplogo\fontsize{30}{30}\selectfont\color{popink}OnBuch\color{poporange}+}\hfill
  \raisebox{0.6ex}{\poppill{Cours signé \textbullet\ Premium}{popink}{popcream}}\par
  \vspace{1pt}\popsquiggle{poppurple}\par
  \vspace{8pt}
  \begin{tcolorbox}[enhanced,colback=poporange,colframe=popink,boxrule=2.8pt,arc=13pt,
      drop shadow=popink,left=18pt,right=18pt,top=12pt,bottom=13pt,after skip=10pt]
    \poppill{#2 \textbullet\ #3}{popink}{popcream}\hspace{5pt}\poppill{#4}{white}{popink}\par\vspace{8pt}
    {\popdisplay\fontsize{14}{14}\selectfont\color{popink}LEÇON #5}\par\vspace{4pt}
    {\raggedright\hyphenpenalty=10000\exhyphenpenalty=10000\popdisplay\fontsize{25}{27}\selectfont\color{popcream}\MakeUppercase{#1}\par}
    \vspace{8pt}
    {\popxbold\fontsize{9.5}{9.5}\selectfont\color{popink}\MakeUppercase{Durée conseillée : #6}}
  \end{tcolorbox}
  \begin{tcolorbox}[enhanced,colback=white,colframe=popink,boxrule=2.2pt,arc=10pt,
      drop shadow=popink,left=16pt,right=16pt,top=10pt,bottom=10pt,after skip=8pt]
    \poppill{Dans cette leçon}{poppurple}{white}\par\vspace{5pt}
    \popsemi\fontsize{9.3}{12.6}\selectfont\color{popink}#7
  \end{tcolorbox}
  \noindent\begin{minipage}[t]{0.487\linewidth}
    \begin{tcolorbox}[enhanced,colback=popgreen,colframe=popink,boxrule=2.2pt,arc=10pt,
        drop shadow=popink,left=11pt,right=11pt,top=10pt,bottom=10pt,height=3.1cm,valign=center]
      \tcbox[colback=white,colframe=popink,boxrule=1.3pt,arc=5pt,boxsep=0pt,nobeforeafter,
        left=3pt,right=3pt,top=3pt,bottom=3pt,valign=center]{\qrcode[height=1.9cm]{https://play.google.com/store/apps/details?id=com.luvvix.onbuch&hl=fr}}%
      \hspace{8pt}\begin{minipage}[c]{0.5\linewidth}
        {\popdisplay\fontsize{11}{12.5}\selectfont\color{white}\MakeUppercase{Télécharge OnBuch}}\par\vspace{4pt}
        {\raggedright\popsemi\fontsize{8.4}{11}\selectfont\color{white}Gratuit \textbullet\ Google Play\\Tuteur IA, annales, concours, résultats\par}
      \end{minipage}
    \end{tcolorbox}
  \end{minipage}\hfill
  \begin{minipage}[t]{0.487\linewidth}
    \begin{tcolorbox}[enhanced,colback=poppurple,colframe=popink,boxrule=2.2pt,arc=10pt,
        drop shadow=popink,left=11pt,right=11pt,top=10pt,bottom=10pt,height=3.1cm,valign=center]
      \tikz[baseline=-0.7ex]\node[circle,fill=white,draw=popink,line width=1.3pt,minimum size=1.6cm,inner sep=0]{\popdisplay\fontsize{24}{24}\selectfont\color{poppurple}+};%
      \hspace{8pt}\begin{minipage}[c]{0.6\linewidth}
        {\popdisplay\fontsize{11}{12.5}\selectfont\color{white}\MakeUppercase{Passe à OnBuch+}}\par\vspace{4pt}
        {\raggedright\popsemi\fontsize{8.4}{11}\selectfont\color{white}Tous les cours signés, Tuteur IA illimité, fiches PDF — onglet OnBuch+ de l'app\par}
      \end{minipage}
    \end{tcolorbox}
  \end{minipage}
  \vspace{8pt}
  \popcontenu
  \vfill
  \signatureonbuch
  \restoregeometry
  \newpage
}

% Rangée « Ce cours contient » (page de garde)
\newcommand{\poptile}[3]{% #1 couleur #2 gros texte #3 légende
  \begin{tcolorbox}[enhanced,colback=#1,colframe=popink,boxrule=2pt,arc=9pt,shadow={2.5pt}{-2.5pt}{0pt}{popink},
      left=6pt,right=6pt,top=9pt,bottom=9pt,halign=center,valign=center,height=1.75cm,width=\linewidth,nobeforeafter]
    {\popdisplay\fontsize{12.5}{13}\selectfont\color{white}\MakeUppercase{#2}}\par\vspace{3pt}
    {\centering\popsemi\fontsize{7.8}{9.5}\selectfont\color{white}#3\par}
  \end{tcolorbox}}
\newcommand{\popcontenu}{%
  \par\noindent{\popxboldls\fontsize{7.5}{7.5}\selectfont\color{popmuted}CE COURS SIGNÉ CONTIENT}\par\vspace{7pt}
  \noindent\begin{minipage}[t]{0.235\linewidth}\poptile{popblue}{Cours}{complet, illustré, pas à pas}\end{minipage}\hfill
  \begin{minipage}[t]{0.235\linewidth}\poptile{poporange}{Méthodes}{et exercices résolus}\end{minipage}\hfill
  \begin{minipage}[t]{0.235\linewidth}\poptile{poppink}{Exercices}{type examen + corrigés}\end{minipage}\hfill
  \begin{minipage}[t]{0.235\linewidth}\poptile{popgreen}{Fiche bilan}{l'essentiel à retenir}\end{minipage}\par}

% Sommaire
\newtcolorbox{sommaire}{enhanced,colback=white,colframe=popink,boxrule=2.2pt,arc=10pt,
  drop shadow=popink,left=18pt,right=18pt,top=13pt,bottom=13pt,before skip=6pt,after skip=20pt,
  before upper={\poppill{Sommaire}{poppurple}{white}\par\vspace{9pt}\popsemi\fontsize{10.6}{14.5}\selectfont\color{popink}}}

% Signature de fin de cours — poussée tout en bas de la dernière page
\newcommand{\findecours}{%
  \par\vfill\nopagebreak
  \begin{center}\popsquiggle{poporange}\end{center}\vspace{4pt}
  \signatureonbuch
}
\AtBeginDocument{\def\llbracket{\mathopen{[\mkern-3.5mu[}}\def\rrbracket{\mathclose{]\mkern-3.5mu]}}} % intervalles d'entiers (glyphe ⟦ absent de la police)
% Glyphes absents de Fira Math : redéfinis à partir de glyphes disponibles
\AtBeginDocument{%
  \def\setminus{\mathbin{\backslash}}\let\smallsetminus\setminus
  \def\vdots{\mathord{\vbox{\baselineskip3.2pt\lineskiplimit0pt\kern4pt\hbox{.}\hbox{.}\hbox{.}}}}%
  \def\ddots{\mathinner{\mkern1mu\raise6.5pt\hbox{.}\mkern2mu\raise3.6pt\hbox{.}\mkern2mu\raise0.7pt\hbox{.}\mkern1mu}}%
  \def\triangle{\mathord{\scalebox{0.9}{$\symup{\Delta}$}}}%
  \def\nabla{\mathord{\raisebox{\depth}{\scalebox{1}[-1]{$\symup{\Delta}$}}}}%
  \def\checkmark{\popcheck{popdark}}%
  \def\star{\mathbin{\ast}}\def\bigstar{\mathord{\ast}}%
  \def\diamond{\mathbin{\scalebox{0.6}{$\Diamond$}}}%
  \def\bigcup{\mathop{\mathchoice{\raisebox{-0.3em}{\scalebox{1.7}{$\cup$}}}{\scalebox{1.25}{$\cup$}}{\cup}{\cup}}\displaylimits}%
  \def\bigcap{\mathop{\mathchoice{\raisebox{-0.3em}{\scalebox{1.7}{$\cap$}}}{\scalebox{1.25}{$\cap$}}{\cap}{\cap}}\displaylimits}%
  \def\lesseqqgtr{\mathrel{\lessgtr}}\def\bigoplus{\mathop{\oplus}\displaylimits}%
}
\providecommand{\ssi}{si et seulement si} % abréviation fréquente
\AtBeginDocument{\providecommand{\wideparen}{\overparen}} % arc de cercle

%% FIGFIX-BEGIN v5 — correctifs globaux des figures (docs/onbuch-generation/tools/figfix)
\usepackage{adjustbox}\usepackage{etoolbox}\tcbuselibrary{raster}
% toute figure TikZ/pgfplots plus large que la ligne est réduite à la largeur disponible
\BeforeBeginEnvironment{tikzpicture}{\begin{adjustbox}{max width=\linewidth}}
\AfterEndEnvironment{tikzpicture}{\end{adjustbox}}
% légendes pgfplots lisibles par-dessus les courbes
\pgfplotsset{every axis/.append style={legend style={fill=white,fill opacity=0.92,text opacity=1,draw=black!15,inner sep=3pt}}}
% étiquettes d'arêtes (syntaxe edge["..."]) avec fond blanc
\tikzset{every edge quotes/.append style={fill=white,inner sep=1.2pt}}
%% FIGFIX-END
\begin{document}

\pagegarde{Lesson 25 — Writing: Compositions about developing and sharing interests in a trade/job/profession}{Anglais}{1ère A}{Chapter 2 : Economic Life and Occupations}{ENA25}{2 h}{Cette leçon d'expression écrite apprend aux élèves de 1ère A à planifier et rédiger une composition (introduction, développement, conclusion) expliquant comment ils développent et partagent leur intérêt pour un métier, un trade ou une profession. Elle fournit connecteurs logiques, temps verbaux, vocabulaire des métiers, un modèle commenté, les erreurs à éviter et une grille d'évaluation.
\vspace{4pt}
\begin{multicols}{2}\small
\begin{itemize}
\item À la fin de cette leçon, je sais identifier la structure d'une composition : introduction, développement, conclusion.
\item À la fin de cette leçon, je sais construire un plan de rédaction (brainstorming, plan détaillé) avant d'écrire.
\item À la fin de cette leçon, je sais utiliser le vocabulaire des métiers, trades et professions en contexte.
\item À la fin de cette leçon, je sais employer les connecteurs logiques (first, moreover, however, therefore, in conclusion) pour organiser mes idées.
\item À la fin de cette leçon, je sais choisir les temps verbaux corrects pour parler de mon intérêt, de mes expériences et de mes projets professionnels.
\item À la fin de cette leçon, je sais rédiger une introduction accrocheuse et une conclusion efficace.
\end{itemize}\end{multicols}}

\begin{sommaire}
\begin{multicols}{2}
\begin{enumerate}
\item Découverte : lire un modèle de composition
\item Le vocabulaire des métiers et de la formation
\item Connecteurs logiques et temps verbaux
\item La structure de la composition et le plan de rédaction
\item Modèle commenté et erreurs à éviter
\item Grille d'évaluation et production guidée
\item Activité d'intégration
\item Méthodes et exercices résolus
\item Exercices
\item Corrigés des exercices
\item Fiche bilan
\end{enumerate}
\end{multicols}
\end{sommaire}

\begin{objectifs}
\begin{itemize}
\item À la fin de cette leçon, je sais identifier la structure d'une composition : introduction, développement, conclusion.
\item À la fin de cette leçon, je sais construire un plan de rédaction (brainstorming, plan détaillé) avant d'écrire.
\item À la fin de cette leçon, je sais utiliser le vocabulaire des métiers, trades et professions en contexte.
\item À la fin de cette leçon, je sais employer les connecteurs logiques (first, moreover, however, therefore, in conclusion) pour organiser mes idées.
\item À la fin de cette leçon, je sais choisir les temps verbaux corrects pour parler de mon intérêt, de mes expériences et de mes projets professionnels.
\item À la fin de cette leçon, je sais rédiger une introduction accrocheuse et une conclusion efficace.
\item À la fin de cette leçon, je sais relire et corriger ma production (grammaire, orthographe, ponctuation, cohérence).
\item À la fin de cette leçon, je sais évaluer une composition à l'aide d'une grille critériée.
\end{itemize}
\end{objectifs}

\begin{prerequis}
\begin{itemize}
\item Vocabulaire de base des métiers et professions (Seconde : jobs and occupations)
\item Les temps principaux : present simple, present continuous, past simple, future with 'going to' et 'will'
\item La rédaction d'un paragraphe simple (topic sentence + supporting sentences)
\item Les connecteurs de base : and, but, because, so
\item La structure d'une phrase anglaise (sujet + verbe + complément)
\end{itemize}
\end{prerequis}

\newpage

\coursec{Découverte : lire un modèle de composition}

\courssub{Situation d'accroche et problématique}

\begin{textebox}[At the school magazine — Buea]
At the Government Technical High School in Buea, students are asked to write a short composition for the school magazine: “Why I want to become a mechanic / a nurse / a computer engineer.” In Cameroon, choosing a trade early — carpentry, tailoring, farming, ICT — can change a family's future. In the UK and the USA, teenagers also write “career essays” to explain their dream job. Writing about your interest in a trade or profession is a skill you will need for the Bac, for scholarship applications and for job motivation letters. Today, you will learn how to plan, organise and write a clear, convincing composition on this topic.
\end{textebox}

Avant d'écrire, il faut savoir \emph{lire} un bon texte. C'est la première étape de tout travail d'expression écrite : observer comment un modèle réussi est construit, quels mots il utilise, comment les idées s'enchaînent. La problématique de cette leçon est la suivante : \textbf{qu'est-ce qui rend une composition sur un métier claire, bien organisée et convaincante ?} Pour y répondre, nous allons lire ensemble une composition modèle, la décortiquer, puis repérer ses outils de langue.

\begin{definition}[What is a composition?]
A \cle{composition} is an organised piece of writing with an \cle{introduction}, a \cle{body} and a \cle{conclusion}, expressing and developing ideas on a given topic.\\
En français : une composition est un texte organisé comportant une introduction, un développement et une conclusion, qui exprime et développe des idées sur un sujet donné. Au Bac camerounais, elle fait généralement entre 150 et 250 mots.
\end{definition}

\courssub{Lecture du texte modèle}

\begin{methode}[How to read a model composition — 4 steps]
\begin{enumerate}
\item \textbf{First reading for gist} : lire une première fois en silence, sans dictionnaire, pour comprendre le sens général (\eng{Who writes? About what? For whom?}).
\item \textbf{Identify the three parts} : repérer l'introduction, le développement (2 ou 3 paragraphes) et la conclusion.
\item \textbf{Highlight the language tools} : souligner les connecteurs (\eng{first, then, moreover, however, in conclusion}), le vocabulaire des métiers et les temps verbaux utilisés.
\item \textbf{Steal good ideas} : noter dans son cahier les expressions et les bonnes idées à réutiliser dans sa propre production.
\end{enumerate}
\end{methode}

Lisez maintenant le texte modèle une première fois en silence, puis à voix haute, en appliquant les étapes 1 et 2.

\begin{textebox}[Developing my interest in tailoring — a model composition]
Ever since I was a child, I have admired my aunt who sews beautiful dresses in Bamenda. Whenever I visited her workshop, I watched her transform a simple piece of fabric into an elegant outfit. That is how my interest in tailoring was born.

First, she taught me how to use a needle and thread; then I learnt to cut fabric and to take measurements. Little by little, I started making simple skirts and shirts for my younger brothers and sisters. Moreover, tailoring allows me to be creative: I can choose colours, patterns and styles that please my customers. It also allows me to earn money. Last holidays, I sewed five dresses for my neighbours and earned enough to buy my school books.

However, I know I need more training. Tailoring is not only about sewing; a good tailor must also understand fashion, manage money and communicate with customers. That is why I plan to attend a technical school after my A-Levels. I would like to obtain a professional certificate and, one day, open my own workshop in Bamenda.

In conclusion, tailoring is not just a hobby for me: it is the profession I want to share with my community. I believe that a skilled tailor can create jobs, dress people with dignity and contribute to the economy of our country.
\end{textebox}

\textbf{Glossaire} (mots difficiles) :

\begin{center}
\begin{tabularx}{\linewidth}{|l|l|X|}
\hline
\rowcolor{popblueL} \textbf{English} & \textbf{Nature} & \textbf{Français} \\
\hline
to sew (sewed, sewn) & verbe irrégulier & coudre \\
\hline
a needle & nom & une aiguille \\
\hline
thread & nom & du fil (à coudre) \\
\hline
fabric & nom & le tissu, l'étoffe \\
\hline
to take measurements & expression & prendre les mesures \\
\hline
a workshop & nom & un atelier \\
\hline
to earn money & expression & gagner de l'argent \\
\hline
training & nom & la formation \\
\hline
skilled & adjectif & qualifié, compétent \\
\hline
a hobby & nom & un passe-temps \\
\hline
\end{tabularx}
\end{center}

\courssub{Repérage global : qui écrit ? de quoi ? pour qui ?}

Après une première lecture, répondez mentalement à ces trois questions essentielles :

\begin{itemize}
\item \textbf{Who writes?} A student (probably a girl or a boy in a technical stream) whose aunt is a tailor in Bamenda. — \emph{Un élève dont la tante est couturière à Bamenda.}
\item \textbf{What is it about?} The writer's growing interest in tailoring and his or her plans for the future. — \emph{L'intérêt croissant de l'auteur pour la couture et ses projets d'avenir.}
\item \textbf{For whom?} For a teacher, an examiner or the readers of a school magazine. — \emph{Pour un enseignant, un examinateur ou les lecteurs d'un journal scolaire.}
\end{itemize}

Ce repérage global est fondamental : avant d'écrire votre propre composition, posez-vous toujours ces trois questions. Elles déterminent le \cle{ton} (ici : personnel mais sérieux), le \cle{registre} (semi-formel) et le \cle{point de vue} (première personne, \eng{I}).

\courssub{Repérage des trois grandes parties}

Une composition bien construite ressemble à un immeuble : l'introduction pose les fondations, le développement élève les étages, la conclusion pose le toit.

\begin{popfigure}
\begin{center}
\begin{tikzpicture}[node distance=0.5cm]
\node[draw=popblue, fill=popblueL, rounded corners, minimum width=8.5cm, minimum height=1.1cm, align=center, font=\small\bfseries] (intro) {INTRODUCTION\\[-0.1cm]{\footnotesize\mdseries hook + topic sentence — “Ever since I was a child...”}};
\node[draw=poporange, fill=poporangeL, rounded corners, minimum width=8.5cm, minimum height=2.6cm, align=center, font=\small\bfseries, below=of intro] (body) {BODY (development)\\[0.05cm]{\footnotesize\mdseries P1 : how I learnt the trade (first, then)}\\{\footnotesize\mdseries P2 : why I love it (moreover) — creativity + money}\\{\footnotesize\mdseries P3 : my future plans (however, that is why)}};
\node[draw=popgreen, fill=popgreenL, rounded corners, minimum width=8.5cm, minimum height=1.1cm, align=center, font=\small\bfseries, below=of body] (concl) {CONCLUSION\\[-0.1cm]{\footnotesize\mdseries restate the main idea + open up — “In conclusion...”}};
\draw[-{Stealth}, very thick, popdark] (intro) -- (body);
\draw[-{Stealth}, very thick, popdark] (body) -- (concl);
\end{tikzpicture}
\end{center}
\legende{La structure en trois parties de la composition modèle}
\end{popfigure}

\begin{propriete}[Rôle de chaque partie]
\begin{itemize}
\item \textbf{Introduction} (2–3 phrases) : elle \emph{accroche} le lecteur (\eng{hook}) et annonce le sujet (\eng{topic sentence}). Ici : \eng{“Ever since I was a child, I have admired my aunt...”}
\item \textbf{Body / development} (2–3 paragraphes) : chaque paragraphe développe \emph{une seule idée principale}, introduite par un connecteur et illustrée par un exemple concret.
\item \textbf{Conclusion} (2–3 phrases) : elle reformule l'idée principale sans copier l'introduction et ouvre sur une perspective (\eng{“...contribute to the economy of our country.”}).
\end{itemize}
\end{propriete}

\courssub{Chasse aux outils de langue dans le texte}

Troisième étape de la méthode : soulignons maintenant dans le texte trois familles d'outils.

\textbf{1. Le vocabulaire des métiers et de la formation.} Relevez : \eng{tailoring, tailor, workshop, needle, thread, fabric, measurements, customers, training, technical school, certificate, profession, skilled, to earn money}. Ce champ lexical précis montre à l'examinateur que l'élève maîtrise le vocabulaire du thème.

\textbf{2. Les connecteurs logiques.} Ce sont les « charnières » du texte :

\begin{center}
\begin{tabularx}{\linewidth}{|l|X|X|}
\hline
\rowcolor{popblueL} \textbf{Connecteur} & \textbf{Fonction} & \textbf{Exemple dans le texte} \\
\hline
\eng{first / then} & ordre chronologique & \eng{First, she taught me...; then I learnt to cut fabric.} \\
\hline
\eng{moreover} & ajout d'un argument & \eng{Moreover, tailoring allows me to be creative.} \\
\hline
\eng{however} & opposition, nuance & \eng{However, I know I need more training.} \\
\hline
\eng{that is why / so} & conséquence & \eng{That is why I plan to attend a technical school.} \\
\hline
\eng{in conclusion} & annonce de la fin & \eng{In conclusion, tailoring is not just a hobby...} \\
\hline
\end{tabularx}
\end{center}

\textbf{3. Les temps verbaux.} Observez la logique des temps :

\begin{center}
\begin{tabularx}{\linewidth}{|X|X|l|}
\hline
\rowcolor{popblueL} \textbf{Temps} & \textbf{Exemple du texte} & \textbf{Emploi} \\
\hline
Present perfect & \eng{I have admired my aunt} & bilan d'une expérience qui dure \\
\hline
Simple past & \eng{she taught me, I learnt, I sewed} & événements passés datés \\
\hline
Simple present & \eng{tailoring allows me to be creative} & vérités générales, goûts \\
\hline
Futur / projets & \eng{I plan to attend..., I would like to obtain...} & projets d'avenir \\
\hline
\end{tabularx}
\end{center}

\begin{attention}[Pièges fréquents des francophones]
\begin{itemize}
\item Ne dites pas \eng{*I have learnt to sew last year} : avec un moment passé précis (\eng{last year}), utilisez le simple past : \eng{I learnt to sew last year.} « J'ai appris à coudre l'année dernière. »
\item \eng{Actually} ne veut pas dire « actuellement » mais « en fait ». Pour « actuellement », dites \eng{at present} ou \eng{currently}.
\item \eng{A formation} n'existe pas : « une formation » se dit \eng{training} (invariable). Pas de \eng{*trainings}.
\item Attention à l'orthographe et à la prononciation : \eng{to sew} se prononce « so » (comme \eng{so}) ; le \emph{w} ne se prononce pas. Ne confondez pas avec \eng{to sow} (semer), qui se prononce pareil !
\end{itemize}
\end{attention}

\courssub{Carte mentale : les idées à développer sur “My dream job”}

Quand le sujet porte sur le métier de vos rêves, cinq axes permettent de trouver des idées rapidement. Mémorisez cette carte : elle vous servira de plan de brainstorming avant toute rédaction.

\begin{popfigure}
\begin{center}
\begin{center}\tcbox[colback=popblue,colframe=popblue,coltext=white,arc=9pt,boxsep=4pt,left=6pt,right=6pt]{\bfseries\small My dream job}\end{center}
\vspace{-2pt}
\begin{tcbraster}[raster columns=3,raster equal height=rows,raster column skip=6pt,raster row skip=6pt,enhanced,blankest]
\begin{tcolorbox}[enhanced,colback=white,colframe=poporange,boxrule=0.9pt,arc=3pt,title={reasons {\footnotesize why?}},fonttitle=\bfseries\scriptsize,coltitle=white,colbacktitle=poporange,left=4pt,right=4pt,top=2pt,bottom=2pt,boxsep=1pt,toptitle=1.5pt,bottomtitle=1.5pt,halign=left]
\begin{itemize}[nosep,leftmargin=*,itemsep=1pt,font=\scriptsize]
\item \footnotesize passion
\item \footnotesize money
\end{itemize}
\end{tcolorbox}
\begin{tcolorbox}[enhanced,colback=white,colframe=popgreen,boxrule=0.9pt,arc=3pt,title={skills {\footnotesize what I can do}},fonttitle=\bfseries\scriptsize,coltitle=white,colbacktitle=popgreen,left=4pt,right=4pt,top=2pt,bottom=2pt,boxsep=1pt,toptitle=1.5pt,bottomtitle=1.5pt,halign=left]
\begin{itemize}[nosep,leftmargin=*,itemsep=1pt,font=\scriptsize]
\item \footnotesize creativity
\item \footnotesize patience
\end{itemize}
\end{tcolorbox}
\begin{tcolorbox}[enhanced,colback=white,colframe=poppurple,boxrule=0.9pt,arc=3pt,title={training {\footnotesize how to learn}},fonttitle=\bfseries\scriptsize,coltitle=white,colbacktitle=poppurple,left=4pt,right=4pt,top=2pt,bottom=2pt,boxsep=1pt,toptitle=1.5pt,bottomtitle=1.5pt,halign=left]
\begin{itemize}[nosep,leftmargin=*,itemsep=1pt,font=\scriptsize]
\item \footnotesize technical school
\item \footnotesize certificate
\end{itemize}
\end{tcolorbox}
\begin{tcolorbox}[enhanced,colback=white,colframe=poppink,boxrule=0.9pt,arc=3pt,title={people who inspire me},fonttitle=\bfseries\scriptsize,coltitle=white,colbacktitle=poppink,left=4pt,right=4pt,top=2pt,bottom=2pt,boxsep=1pt,toptitle=1.5pt,bottomtitle=1.5pt,halign=left]
\begin{itemize}[nosep,leftmargin=*,itemsep=1pt,font=\scriptsize]
\item \footnotesize my aunt
\item \footnotesize a famous tailor
\end{itemize}
\end{tcolorbox}
\begin{tcolorbox}[enhanced,colback=white,colframe=popgold,boxrule=0.9pt,arc=3pt,title={future plans},fonttitle=\bfseries\scriptsize,coltitle=white,colbacktitle=popgold,left=4pt,right=4pt,top=2pt,bottom=2pt,boxsep=1pt,toptitle=1.5pt,bottomtitle=1.5pt,halign=left]
\begin{itemize}[nosep,leftmargin=*,itemsep=1pt,font=\scriptsize]
\item \footnotesize open a workshop
\item \footnotesize create jobs
\end{itemize}
\end{tcolorbox}
\end{tcbraster}

\end{center}
\legende{Carte mentale de brainstorming autour de “My dream job”}
\end{popfigure}

\courssub{Questions de compréhension sur le texte modèle}

\begin{exoresolu}[Reading comprehension — solved questions]
\textbf{A. True or False?} Justify with a quotation from the text.
\begin{enumerate}
\item The writer's aunt works in Douala.
\item The writer has already earned money thanks to tailoring.
\item The writer thinks he or she knows everything about tailoring.
\end{enumerate}
\textbf{B. Multiple choice.} The writer plans to attend a technical school : (a) before the A-Levels (b) after the A-Levels (c) instead of the A-Levels.
\textbf{C. Short answer.} Give two reasons why the writer loves tailoring.
\tcblower
\textbf{A.} 1. \textbf{False} — \eng{“...my aunt who sews beautiful dresses in Bamenda.”} 2. \textbf{True} — \eng{“I sewed five dresses for my neighbours and earned enough to buy my school books.”} 3. \textbf{False} — \eng{“However, I know I need more training.”}
\textbf{B.} (b) — \eng{“I plan to attend a technical school after my A-Levels.”}
\textbf{C.} Two possible reasons : \eng{“Tailoring allows me to be creative”} and \eng{“It also allows me to earn money.”}
\end{exoresolu}

\begin{exercice}[Your turn — find the tools in the text]
\begin{enumerate}
\item Underline three connectors of addition or opposition in the model composition and give their French equivalents.
\item Find one verb in the present perfect and two verbs in the simple past. Copy them with their subject.
\item Identify the sentence that announces the topic in the introduction (the \eng{topic sentence}).
\item In one sentence, say what makes this composition convincing.
\end{enumerate}
\end{exercice}

\begin{corrige}[Exercice — find the tools in the text]
\begin{enumerate}
\item Par exemple : \eng{moreover} « de plus », \eng{however} « cependant », \eng{that is why} « c'est pourquoi » (on accepte aussi \eng{first, then, in conclusion}).
\item Present perfect : \eng{I have admired}. Simple past : \eng{she taught me}, \eng{I learnt}, \eng{I sewed}, \eng{I earned}.
\item \eng{“That is how my interest in tailoring was born.”} (elle annonce le sujet : la couture).
\item Exemple de réponse : \eng{The composition is convincing because it is well organised, uses clear connectors and gives concrete examples from the writer's real life.}
\end{enumerate}
\end{corrige}

\courssub{Discussion : qu'est-ce qui rend cette composition convaincante ?}

\begin{experience}[Class discussion — What makes it good?]
En groupes de quatre, discutez en anglais pendant cinq minutes : \eng{“What makes this composition clear and convincing?”} Utilisez les amorces suivantes, puis un rapporteur présente vos conclusions à la classe :
\begin{itemize}
\item \eng{In my opinion, the introduction is effective because...}
\item \eng{I think the examples are convincing because...}
\item \eng{The connectors help the reader to...}
\item \eng{What I would like to imitate in my own composition is...}
\end{itemize}
\end{experience}

Éléments de réponse attendus : le texte est \cle{personnel} (exemples vécus : la tante, les voisins), \cle{organisé} (trois parties nettes), \cle{nuancé} (\eng{however} montre que l'élève sait qu'il lui reste à apprendre) et \cle{tourné vers l'avenir} (projets concrets). Ce sont exactement les critères de la grille d'évaluation du Bac.

\begin{savaistu}[Career essays in the UK and the USA]
Au Royaume-Uni et aux États-Unis, les lycéens rédigent des \eng{personal statements} et des \eng{career essays} pour entrer à l'université ou obtenir une bourse : ils doivent y expliquer leur parcours, leurs motivations et leurs projets professionnels — exactement le même type d'écrit qu'au Bac camerounais ! Savoir écrire une composition sur son intérêt pour un métier est donc une compétence internationale, utile bien au-delà de l'examen.
\end{savaistu}

\begin{aretenir}[À retenir de cette section]
\begin{itemize}
\item Une composition = \textbf{introduction} (accroche + sujet) + \textbf{body} (2–3 paragraphes, une idée chacun) + \textbf{conclusion} (bilan + ouverture).
\item Lire un modèle en 4 étapes : sens général, trois parties, outils de langue, idées à réutiliser.
\item Les connecteurs (\eng{first, then, moreover, however, that is why, in conclusion}) guident le lecteur.
\item Les temps obéissent à une logique : present perfect pour l'expérience, simple past pour les faits datés, present pour les vérités, \eng{plan to / would like to} pour les projets.
\end{itemize}
\end{aretenir}

\coursec{Le vocabulaire des métiers et de la formation}

Dans la section précédente, nous avons lu et observé un modèle de composition intitulé \eng{“Why I want to become a mechanic”}. Vous avez remarqué que ce texte utilisait des mots comme \eng{trade}, \eng{apprenticeship}, \eng{skills} ou encore \eng{to follow in my father's footsteps}. Avant de rédiger votre propre composition, il faut donc d'abord maîtriser ce vocabulaire. C'est l'objectif de cette section : construire, champ lexical par champ lexical, tout le lexique nécessaire pour parler des métiers, de la formation et de l'intérêt que l'on porte à un métier.

\courssub{Observer avant d'apprendre : un court texte de référence}

Lisons d'abord ce petit texte original. Soulignez mentalement tous les mots qui se rapportent aux métiers et à la formation.

\begin{textebox}[At the vocational training centre in Bamenda]
Every morning, twenty young people arrive at the vocational training centre in Bamenda. Some of them want to learn a trade: Grace is training to become a tailor, and Peter is doing an apprenticeship with a local carpenter. Others dream of a profession that requires long studies. “I want to become a nurse,” says Lilian, “because I am passionate about helping sick people.” Her friend Samuel wants to be an electrician first, and later an engineer. “My uncle is a plumber,” he explains. “He followed in his grandfather's footsteps, and he earns a good living. He taught me that a trade is not a small thing: it is a real career.”

The centre offers workshops in carpentry, plumbing, mechanics and hairdressing. At the end of the training, every student receives a certificate. With this qualification, they can find a job or even set up their own business. Mrs Ndi, the director, is proud of her students. “Some people think that only doctors, lawyers and engineers are important,” she says. “But who repairs your car? Who builds your house? A skilled mechanic or a good baker is just as useful to society.”

Many students also want to share their knowledge. “When I finish my training,” says Grace, “I will teach other girls in my village. My teacher inspired me, and now I want to inspire others.”
\end{textebox}

\textbf{Glossaire :}
\begin{itemize}
\item \cle{vocational training centre} (n.) : centre de formation professionnelle
\item \cle{apprenticeship} (n.) : apprentissage (formation pratique auprès d'un patron)
\item \cle{skilled} (adj.) : qualifié, compétent
\item \cle{to earn a living} (exp.) : gagner sa vie
\item \cle{to set up a business} (exp.) : créer une entreprise
\item \cle{to inspire} (v.) : inspirer, donner envie
\end{itemize}

\textbf{Questions de compréhension :}
\begin{enumerate}
\item What trade is Grace learning?
\item Why does Lilian want to become a nurse?
\item What do students receive at the end of their training?
\item What does Grace want to do after her training?
\end{enumerate}

\courssub{Champ lexical 1 : les trades (métiers manuels et artisanaux)}

\begin{definition}[Trade]
A \cle{trade} is a job that requires special manual or practical skills, e.g. carpentry, plumbing. (Un métier artisanal, manuel, qui s'apprend surtout par la pratique.)
\end{definition}

Observez les phrases suivantes, tirées de situations camerounaises quotidiennes :

\eng{My father is a carpenter; he makes furniture in his workshop in Buea.} « Mon père est menuisier ; il fabrique des meubles dans son atelier à Buea. »

\eng{The tailor sewed my school uniform in two days.} « Le tailleur a cousu mon uniforme scolaire en deux jours. »

\eng{We called a plumber because the kitchen tap was leaking.} « Nous avons appelé un plombier parce que le robinet de la cuisine fuyait. »

\begin{center}
\begin{tabularx}{\linewidth}{|X|X|X|}\hline
\rowcolor{popblueL} English & Français & Remarque \\ \hline
a tailor & un tailleur & attention : \eng{tailor}, pas \eng{taylor} \\ \hline
a carpenter & un menuisier & la menuiserie : \eng{carpentry} \\ \hline
a mechanic & un mécanicien & il répare les voitures et les bendskins \\ \hline
an electrician & un électricien & article \eng{an} devant voyelle \\ \hline
a plumber & un plombier & le \eng{b} final ne se prononce pas : « pleu-meur » \\ \hline
a hairdresser & un coiffeur / une coiffeuse & un seul mot \\ \hline
a farmer & un agriculteur, un cultivateur & très courant dans les villages du Cameroun \\ \hline
a baker & un boulanger & le pain : \eng{bread} \\ \hline
\end{tabularx}
\end{center}

\begin{attention}[Prononciation et orthographe]
\eng{Plumber} se prononce « pleu-meur » (le \eng{b} est muet). \eng{Carpenter} se prononce « car-peun-teur ». N'écrivez jamais \eng{taylor} pour \eng{tailor} : \eng{Taylor} est un nom propre anglais !
\end{attention}

\courssub{Champ lexical 2 : les professions (professions qualifiées)}

\begin{definition}[Profession]
A \cle{profession} is a job that requires advanced education and training, e.g. medicine, law. (Une profession intellectuelle qui exige de longues études, souvent à l'université.)
\end{definition}

\eng{She is passionate about nursing.} « Elle est passionnée par le métier d'infirmière. »

\eng{My elder sister studied law in Yaoundé and now she is a lawyer.} « Ma grande sœur a étudié le droit à Yaoundé et elle est maintenant avocate. »

\begin{center}
\begin{tabularx}{\linewidth}{|X|X|X|}\hline
\rowcolor{popblueL} English & Français & Remarque \\ \hline
a doctor & un médecin & \eng{to see a doctor} = consulter un médecin \\ \hline
a nurse & un(e) infirmier(-ière) & le métier : \eng{nursing} \\ \hline
a teacher & un enseignant, un professeur & faux ami : voir plus bas \\ \hline
a lawyer & un avocat, un juriste & le droit : \eng{law} \\ \hline
an engineer & un ingénieur & prononcer « èn-dji-nieur » \\ \hline
an accountant & un comptable & la comptabilité : \eng{accountancy} \\ \hline
a journalist & un journaliste & le journalisme : \eng{journalism} \\ \hline
a computer scientist & un informaticien & l'informatique : \eng{computer science} \\ \hline
\end{tabularx}
\end{center}

\begin{attention}[Faux ami : teacher $\neq$ professeur ?]
\eng{A teacher} se traduit bien par « un enseignant / un professeur ». Mais attention au faux ami inverse : le mot anglais \eng{a professor} désigne un \emph{professeur d'université} (le grade le plus élevé). Ne dites pas \eng{My father is a professor} si votre père enseigne au lycée : dites \eng{My father is a teacher}.
\end{attention}

\courssub{Champ lexical 3 : la formation et l'apprentissage}

Pour raconter comment on se prépare à un métier, il faut le vocabulaire de la formation.

\eng{He did a three-year apprenticeship with a mechanic in Douala.} « Il a fait un apprentissage de trois ans chez un mécanicien à Douala. »

\eng{After two years at a technical college, she received her certificate.} « Après deux ans dans un collège technique, elle a reçu son certificat. »

\begin{center}
\begin{tabularx}{\linewidth}{|X|X|X|}\hline
\rowcolor{popblueL} English & Français & Remarque \\ \hline
training & la formation & mot incontournable \\ \hline
an apprenticeship & un apprentissage & \eng{to do an apprenticeship} \\ \hline
a vocational school & une école de formation professionnelle & \eng{vocational} = professionnel \\ \hline
a technical college & un collège technique, un lycée technique & comme le CETIC ou un GTHS au Cameroun \\ \hline
a workshop & un atelier & aussi : un stage pratique \\ \hline
skills (pl.) & les compétences, le savoir-faire & presque toujours au pluriel \\ \hline
a qualification & un diplôme, une qualification & \eng{to have the right qualifications} \\ \hline
a certificate & un certificat, un diplôme & \eng{to receive a certificate} \\ \hline
\end{tabularx}
\end{center}

\begin{attention}[Faux ami célèbre : formation]
Le français « formation » se dit \cle{training} en anglais. Le mot anglais \eng{formation} existe, mais il signifie « formation / constitution » au sens de \emph{création, structure} (the formation of a group, the formation of a rock). Écrivez donc \eng{I received good training} et jamais \eng{I received a good formation}.
\end{attention}

\courssub{Champ lexical 4 : développer et partager un intérêt}

C'est le cœur du sujet de composition : expliquer \emph{pourquoi} un métier vous attire et comment vous voulez le partager.

\begin{center}
\begin{tabularx}{\linewidth}{|X|X|X|}\hline
\rowcolor{popblueL} English & Français & Construction \\ \hline
to be interested in + n./V-ing & être intéressé par & \eng{I am interested in mechanics.} \\ \hline
to be passionate about + n./V-ing & être passionné par & \eng{She is passionate about teaching.} \\ \hline
to develop a taste for + n. & prendre goût à & \eng{He developed a taste for carpentry.} \\ \hline
to learn from + someone & apprendre auprès de quelqu'un & \eng{I learnt a lot from my uncle.} \\ \hline
to teach others & enseigner aux autres & \eng{I want to teach others.} \\ \hline
to inspire + someone & inspirer quelqu'un & \eng{My teacher inspired me.} \\ \hline
to share one's knowledge & partager son savoir & \eng{She shares her knowledge with the girls.} \\ \hline
\end{tabularx}
\end{center}

\begin{propriete}[Prépositions obligatoires]
Après \eng{interested} on utilise toujours \cle{in} ; après \eng{passionate} toujours \cle{about}. Si un verbe suit, il prend la forme en \cle{-ing} : \eng{I am interested in becoming a nurse.} « Je suis intéressé par le fait de devenir infirmière. » Jamais \eng{interested to become} ni \eng{passionate of}.
\end{propriete}

\courssub{Collocations utiles pour la composition}

Une \cle{collocation} est une association de mots qui vont naturellement ensemble. Les employer rend votre copie plus riche et plus naturelle.

\begin{exemplebox}[Cinq collocations à réutiliser]
\begin{itemize}
\item \eng{to learn a trade} : apprendre un métier — \eng{I want to learn a trade after my BEPC.} « Je veux apprendre un métier après mon BEPC. »
\item \eng{to acquire skills} : acquérir des compétences — \eng{At the workshop, students acquire practical skills.} « À l'atelier, les élèves acquièrent des compétences pratiques. »
\item \eng{to follow in someone's footsteps} : suivre les traces de quelqu'un — \eng{He followed in his grandfather's footsteps.} « Il a suivi les traces de son grand-père. »
\item \eng{to earn a living} : gagner sa vie — \eng{She earns a living as a hairdresser.} « Elle gagne sa vie comme coiffeuse. »
\item \eng{to set up a business} : créer une entreprise — \eng{One day, I will set up my own business.} « Un jour, je créerai ma propre entreprise. »
\end{itemize}
\end{exemplebox}

\courssub{Job, trade, profession, career : ne plus les confondre}

\begin{attention}[Quatre mots, quatre sens]
\begin{itemize}
\item \cle{job} : le mot général, « un emploi » (\eng{He is looking for a job.}).
\item \cle{trade} : un métier manuel ou artisanal (carpenter, plumber).
\item \cle{profession} : une profession qualifiée exigeant des études (doctor, lawyer).
\item \cle{career} : la carrière, c'est-à-dire \emph{tout le parcours} professionnel d'une vie (\eng{She had a long career as a nurse.}).
\end{itemize}
Ne dites jamais \eng{I want to do the profession of doctor} : c'est un calque du français. Dites simplement \eng{I want to become a doctor.} « Je veux devenir médecin. »
\end{attention}

\begin{center}
\begin{tabularx}{\linewidth}{|X|X|X|}\hline
\rowcolor{popblueL} Français & English & Exemple \\ \hline
un emploi (général) & a job & \eng{He found a job in Douala.} \\ \hline
un métier manuel & a trade & \eng{Plumbing is a useful trade.} \\ \hline
une profession qualifiée & a profession & \eng{Medicine is a noble profession.} \\ \hline
une carrière & a career & \eng{She started her career in 2010.} \\ \hline
\end{tabularx}
\end{center}

\courssub{Carte mentale récapitulative}

\begin{popfigure}
\begin{center}
\begin{center}\tcbox[colback=popblue,colframe=popblue,coltext=white,arc=9pt,boxsep=4pt,left=6pt,right=6pt]{\bfseries\small JOB}\end{center}
\vspace{-2pt}
\begin{tcbraster}[raster columns=2,raster equal height=rows,raster column skip=6pt,raster row skip=6pt,enhanced,blankest]
\begin{tcolorbox}[enhanced,colback=white,colframe=poporange,boxrule=0.9pt,arc=3pt,title={trades},fonttitle=\bfseries\scriptsize,coltitle=white,colbacktitle=poporange,left=4pt,right=4pt,top=2pt,bottom=2pt,boxsep=1pt,toptitle=1.5pt,bottomtitle=1.5pt,halign=left]
\begin{itemize}[nosep,leftmargin=*,itemsep=1pt,font=\scriptsize]
\item tailor carpenter
\item mechanic plumber
\end{itemize}
\end{tcolorbox}
\begin{tcolorbox}[enhanced,colback=white,colframe=popgreen,boxrule=0.9pt,arc=3pt,title={professions},fonttitle=\bfseries\scriptsize,coltitle=white,colbacktitle=popgreen,left=4pt,right=4pt,top=2pt,bottom=2pt,boxsep=1pt,toptitle=1.5pt,bottomtitle=1.5pt,halign=left]
\begin{itemize}[nosep,leftmargin=*,itemsep=1pt,font=\scriptsize]
\item doctor nurse
\item lawyer engineer
\end{itemize}
\end{tcolorbox}
\begin{tcolorbox}[enhanced,colback=white,colframe=poppurple,boxrule=0.9pt,arc=3pt,title={training},fonttitle=\bfseries\scriptsize,coltitle=white,colbacktitle=poppurple,left=4pt,right=4pt,top=2pt,bottom=2pt,boxsep=1pt,toptitle=1.5pt,bottomtitle=1.5pt,halign=left]
\begin{itemize}[nosep,leftmargin=*,itemsep=1pt,font=\scriptsize]
\item apprenticeship
\item certificate skills
\end{itemize}
\end{tcolorbox}
\begin{tcolorbox}[enhanced,colback=white,colframe=poppink,boxrule=0.9pt,arc=3pt,title={sharing interest},fonttitle=\bfseries\scriptsize,coltitle=white,colbacktitle=poppink,left=4pt,right=4pt,top=2pt,bottom=2pt,boxsep=1pt,toptitle=1.5pt,bottomtitle=1.5pt,halign=left]
\begin{itemize}[nosep,leftmargin=*,itemsep=1pt,font=\scriptsize]
\item passionate about
\item inspire teach others
\end{itemize}
\end{tcolorbox}
\end{tcbraster}

\end{center}
\legende{Carte mentale du vocabulaire : quatre branches autour du mot \eng{job}.}
\end{popfigure}

\courssub{Exercice résolu et entraînement}

\begin{exoresolu}[Compléter avec le bon mot]
Complétez chaque phrase avec un mot de la liste : \eng{trade — training — career — passionate — footsteps}.
\begin{enumerate}
\item My brother is \ldots\ about computers; he wants to become a computer scientist.
\item Carpentry is a \ldots\ that you learn with your hands.
\item She received excellent \ldots\ at the technical college in Bafoussam.
\item I want to follow in my mother's \ldots\ and become a nurse.
\item He had a long \ldots\ as a teacher before he retired.
\end{enumerate}
\tcblower
\textbf{Solution détaillée :}
\begin{enumerate}
\item \cle{passionate} — construction : \eng{to be passionate about} + nom. « Mon frère est passionné par l'informatique. »
\item \cle{trade} — la menuiserie est un métier manuel, donc \eng{trade}, pas \eng{profession}. « La menuiserie est un métier qui s'apprend avec les mains. »
\item \cle{training} — « formation » en français = \eng{training} (faux ami : pas \eng{formation} !).
\item \cle{footsteps} — collocation : \eng{to follow in someone's footsteps}. « Je veux suivre les traces de ma mère. »
\item \cle{career} — tout le parcours professionnel = \eng{career}. « Il a fait une longue carrière comme enseignant. »
\end{enumerate}
\end{exoresolu}

\begin{exercice}[À vous de jouer]
Traduisez en anglais :
\begin{enumerate}
\item Mon oncle est électricien et ma tante est coiffeuse.
\item Elle a fait un apprentissage chez une tailleuse à Bamenda.
\item Je suis intéressé par le métier de journaliste.
\item Après ma formation, je veux créer ma propre entreprise.
\item Mon professeur m'a inspiré et je veux partager mon savoir avec les autres.
\end{enumerate}
\end{exercice}

\begin{corrige}[Exercice « À vous de jouer »]
\begin{enumerate}
\item \eng{My uncle is an electrician and my aunt is a hairdresser.} (article \eng{an} devant voyelle)
\item \eng{She did an apprenticeship with a tailor in Bamenda.}
\item \eng{I am interested in journalism.} (ou \eng{in becoming a journalist})
\item \eng{After my training, I want to set up my own business.}
\item \eng{My teacher inspired me and I want to share my knowledge with others.}
\end{enumerate}
\end{corrige}

\begin{aretenir}[L'essentiel de la section]
\begin{itemize}
\item \eng{trade} = métier manuel ; \eng{profession} = profession qualifiée ; \eng{job} = emploi (général) ; \eng{career} = carrière.
\item « Formation » = \eng{training} ; \eng{formation} est un faux ami.
\item \eng{interested in} et \eng{passionate about} + nom ou verbe en \eng{-ing}.
\item Collocations clés : \eng{learn a trade}, \eng{acquire skills}, \eng{follow in someone's footsteps}, \eng{earn a living}, \eng{set up a business}.
\item Pour dire « je veux devenir médecin » : \eng{I want to become a doctor.}
\end{itemize}
\end{aretenir}

Ce vocabulaire est maintenant votre boîte à outils. Dans la section suivante, nous verrons comment relier ces mots entre eux grâce aux connecteurs logiques et aux temps verbaux adaptés, afin de construire des phrases fluides pour votre composition.

\coursec{Connecteurs logiques et temps verbaux}

Dans la section précédente, nous avons enrichi notre vocabulaire des métiers et de la formation : \eng{a carpenter}, \eng{a tailor}, \eng{an apprenticeship}, \eng{a workshop}… Mais une bonne composition ne se résume pas à une liste de mots. Pour obtenir une excellente note, il faut \cle{relier les idées} entre elles grâce aux \cle{connecteurs logiques} et choisir le \cle{temps verbal} adapté à chaque moment du récit. C'est précisément l'objet de cette section : donner à votre texte du \emph{liant} et une \emph{chronologie claire}.

\courssub{Observer d'abord : un paragraphe bien relié}

\begin{textebox}[A well-connected paragraph --- observe the linking words]
My name is Larissa and I want to become a hairdresser. \textbf{First of all}, I love this trade \textbf{because} it makes people happy and confident. \textbf{When I was a child}, I watched my aunt braiding hair in her small salon in Bafoussam, \textbf{and} I \textbf{have learnt} many techniques from her \textbf{since then}. \textbf{However}, this profession is not easy: it requires patience, creativity and long hours of standing. \textbf{Moreover}, the training is expensive. \textbf{Although} these difficulties are real, I am determined. \textbf{For example}, I practise every weekend on my sisters' hair. \textbf{Therefore}, I work hard at school, \textbf{because} good marks will open the doors of a good training centre. \textbf{In conclusion}, hairdressing is my passion, \textbf{and} I \textbf{am going to open} my own salon one day.
\end{textebox}

\textbf{Glossaire :} \emph{to braid} (v.) : tresser ; \emph{training} (n.) : formation ; \emph{determined} (adj.) : déterminé(e) ; \emph{marks} (n.) : notes ; \emph{centre} (n.) : centre (orthographe britannique).

\textbf{Questions d'observation :}
\begin{enumerate}
\item Soulignez tous les mots qui relient les phrases ou les idées. Que remarquez-vous sur leur position dans la phrase ?
\item Quels temps verbaux Larissa utilise-t-elle pour parler de son enfance ? De son expérience actuelle ? De son projet ?
\end{enumerate}

On observe que les connecteurs se placent généralement \cle{en début de phrase}, suivis d'une virgule, et que chaque étape du récit (souvenir, expérience, passion présente, projet) appelle un temps verbal différent. Étudions cela en détail.

\courssub{Les connecteurs logiques par fonction}

\begin{definition}[Qu'est-ce qu'un connecteur logique ?]
Un \cle{connecteur logique} (\eng{linking word} ou \eng{connector}) est un mot ou une expression qui relie deux phrases ou deux idées et qui indique leur relation : addition, opposition, cause, conséquence, exemple ou conclusion. Sans connecteurs, une composition ressemble à une liste télégraphique ; avec eux, elle devient un texte fluide et argumenté.
\end{definition}

\begin{propriete}[Les cinq familles de connecteurs]
\begin{itemize}
\item \textbf{Addition} (ajouter une idée) : \eng{first}, \eng{first of all} (d'abord), \eng{then} (ensuite), \eng{also} (aussi), \eng{moreover} (de plus), \eng{furthermore} (en outre), \eng{in addition} (en plus).
\item \textbf{Opposition} (contraster) : \eng{but} (mais), \eng{however} (cependant), \eng{although} (bien que), \eng{whereas} (tandis que), \eng{on the other hand} (d'un autre côté).
\item \textbf{Cause et conséquence} : \eng{because} (parce que), \eng{since} / \eng{as} (puisque, comme), \eng{so} (donc), \eng{therefore} (c'est pourquoi, par conséquent), \eng{that is why} (c'est pourquoi).
\item \textbf{Exemple et illustration} : \eng{for example}, \eng{for instance} (par exemple), \eng{such as} (tel que, comme).
\item \textbf{Conclusion} : \eng{in conclusion} (en conclusion), \eng{to sum up} (pour résumer), \eng{all in all} (tout compte fait), \eng{finally} (finalement).
\end{itemize}
\end{propriete}

\begin{center}
\begin{tabularx}{\linewidth}{|l|X|X|}\hline
\rowcolor{popblueL} \textbf{Fonction} & \textbf{Connecteurs} & \textbf{Français} \\ \hline
Addition & first of all, then, moreover, furthermore, in addition, also & d'abord, ensuite, de plus, en outre, aussi \\ \hline
Opposition & but, however, although, whereas, on the other hand & mais, cependant, bien que, tandis que \\ \hline
Cause & because, since, as & parce que, puisque, comme \\ \hline
Conséquence & so, therefore, that is why & donc, c'est pourquoi \\ \hline
Exemple & for example, for instance, such as & par exemple, tel que \\ \hline
Conclusion & in conclusion, to sum up, all in all, finally & en conclusion, pour résumer, finalement \\ \hline
\end{tabularx}
\end{center}

\begin{exemplebox}[Exemples traduits]
\eng{First, I observed my mother; then I tried myself.} « D'abord, j'observais ma mère ; ensuite j'ai essayé moi-même. »

\eng{However, this job requires patience.} « Cependant, ce métier exige de la patience. »

\eng{Therefore, I will work hard at school.} « C'est pourquoi je vais travailler dur à l'école. »

\eng{Although the work is tiring, I enjoy it.} « Bien que le travail soit fatigant, je l'aime. »

\eng{My brother likes office jobs, whereas I prefer manual work.} « Mon frère aime les métiers de bureau, tandis que je préfère le travail manuel. »

\eng{I want to learn a trade such as plumbing or welding.} « Je veux apprendre un métier tel que la plomberie ou la soudure. »

\eng{In conclusion, carpentry is the job of my dreams.} « En conclusion, la menuiserie est le métier de mes rêves. »
\end{exemplebox}

\begin{attention}[Pièges de traduction --- erreurs fréquentes des francophones]
\begin{itemize}
\item « C'est pourquoi » ne se traduit \textbf{jamais} par \eng{it is why}. Dites \eng{that is why} ou \eng{therefore} : \eng{I love this trade. That is why I practise every day.}
\item « Par exemple » se dit \eng{for example}, \textbf{jamais} \eng{by example}.
\item Après \eng{although}, on ne met \textbf{pas} \eng{but} dans la même phrase : \eng{Although the job is difficult, I like it.} (et non \emph{Although the job is difficult, but I like it}). En anglais, un seul connecteur suffit, contrairement au français « bien que… mais… ».
\item \eng{Because} et \eng{so} ne s'emploient pas ensemble : \eng{Because I was ill, I stayed at home} ou \eng{I was ill, so I stayed at home} --- mais pas les deux à la fois.
\item Attention à la virgule : \eng{however}, \eng{therefore}, \eng{moreover}, \eng{for example} en tête de phrase sont suivis d'une virgule.
\end{itemize}
\end{attention}

\begin{methode}[Varier ses connecteurs comme un bon rédacteur]
\begin{enumerate}
\item Avant de rédiger, listez 6 à 8 connecteurs différents que vous comptez utiliser (au moins un par famille).
\item Ne répétez \textbf{jamais} le même connecteur deux fois de suite : après un \eng{first}, utilisez \eng{then} ou \eng{moreover}, pas un autre \eng{first}.
\item Réservez les connecteurs de conclusion (\eng{in conclusion}, \eng{to sum up}) au tout dernier paragraphe.
\item À la relecture, comptez vos connecteurs : visez \textbf{au moins 5 connecteurs différents} dans la composition. Soulignez-les pour vérifier.
\end{enumerate}
\end{methode}

\begin{popfigure}
\begin{tikzpicture}[node distance=0.35cm and 0.5cm]
\node[draw=popblue, fill=popblueL, rounded corners, align=center, font=\bfseries] (add) {Addition\\ first, then, moreover,\\ furthermore, in addition};
\node[draw=poporange, fill=poporangeL, rounded corners, align=center, font=\bfseries, right=of add] (opp) {Opposition\\ but, however, although,\\ whereas, on the other hand};
\node[draw=popgreen, fill=popgreenL, rounded corners, align=center, font=\bfseries, below=0.6cm of add] (cause) {Cause / Conséquence\\ because, since, as,\\ so, therefore, that is why};
\node[draw=poppurple, fill=poppurpleL, rounded corners, align=center, font=\bfseries, right=of cause] (ex) {Exemple\\ for example, for instance,\\ such as};
\node[draw=popgold, fill=popgoldL, rounded corners, align=center, font=\bfseries, below=0.6cm of $(cause)!0.5!(ex)$] (concl) {Conclusion\\ in conclusion, to sum up, all in all, finally};
\draw[-{Stealth}, thick, popmuted] (add) -- (cause);
\draw[-{Stealth}, thick, popmuted] (opp) -- (ex);
\draw[-{Stealth}, thick, popmuted] (cause) -- (concl);
\draw[-{Stealth}, thick, popmuted] (ex) -- (concl);
\end{tikzpicture}
\legende{Carte mentale des connecteurs classés par fonction logique.}
\end{popfigure}

\courssub{Les temps verbaux du récit personnel}

Une composition sur son intérêt pour un métier raconte un parcours : un souvenir d'enfance, une expérience accumulée, une passion présente et un projet d'avenir. Chaque étape appelle un temps précis.

\begin{propriete}[La règle des quatre temps]
\begin{itemize}
\item \textbf{Past simple} : pour un souvenir \cle{daté et terminé} (souvent avec \eng{when I was a child}, \eng{last year}, \eng{in 2020}). \eng{My uncle taught me how to sew.} « Mon oncle m'a appris à coudre. » \eng{Last year, I repaired my first radio.} « L'année dernière, j'ai réparé ma première radio. »
\item \textbf{Present perfect} : pour l'\cle{expérience acquise} jusqu'à aujourd'hui, sans date précise (souvent avec \eng{already}, \eng{since}, \eng{never}, \eng{many}). \eng{I have learnt a lot from my father.} « J'ai beaucoup appris de mon père. » \eng{I have repaired many radios.} « J'ai réparé de nombreuses radios. »
\item \textbf{Present simple} : pour les \cle{goûts permanents} et les habitudes. \eng{I love mechanics.} « J'adore la mécanique. » \eng{I enjoy repairing radios.} « J'aime réparer les radios. »
\item \textbf{Going to / will} : pour les \cle{projets}. \eng{Going to} exprime un projet décidé : \eng{I am going to open a workshop.} « Je vais ouvrir un atelier. » \eng{Will} exprime une promesse ou une prédiction : \eng{I will work hard at school.} « Je vais travailler dur à l'école. »
\end{itemize}
\end{propriete}

\begin{center}
\begin{tabular}{|l|l|l|}\hline
\rowcolor{popblueL} \textbf{Temps} & \textbf{Emploi} & \textbf{Marqueurs typiques} \\ \hline
Past simple & souvenir daté, terminé & when I was a child, last year \\ \hline
Present perfect & expérience acquise & already, since then, many times \\ \hline
Present simple & goût, habitude, passion & always, every day, usually \\ \hline
Going to / will & projet, avenir & one day, in the future, next year \\ \hline
\end{tabular}
\end{center}

\begin{popfigure}
\begin{tikzpicture}[font=\small]
\draw[-{Stealth}, very thick, popdark] (0,0) -- (12.5,0);
\foreach \x/\t/\d/\c in {1.5/{Past simple}/{childhood memory}/poporange, 5/{Present perfect}/{experience so far}/popgreen, 8.5/{Present simple}/{passion now}/popblue, 11.5/{Going to}/{future plans}/poppurple}{
  \fill[\c] (\x,0) circle (0.12);
  \node[above=0.15cm, align=center, text=\c, font=\bfseries] at (\x,0) {\t};
  \node[below=0.15cm, align=center, text=popmuted] at (\x,0) {\d};
}
\node[align=center, text=popmuted, font=\itshape] at (6.25,-1.1) {My uncle taught me $\rightarrow$ I have learnt a lot $\rightarrow$ I love this trade $\rightarrow$ I am going to succeed};
\end{tikzpicture}
\legende{Frise chronologique des temps dans une composition sur un métier.}
\end{popfigure}

\begin{attention}[Erreurs fréquentes des francophones avec les temps]
\begin{itemize}
\item Ne dites pas \emph{I am loving mechanics} : les verbes de goût (\eng{love}, \eng{like}, \eng{enjoy}, \eng{prefer}, \eng{want}) n'ont pas de forme en \emph{-ing}. Dites \eng{I love mechanics}.
\item Le present perfect exige l'auxiliaire \eng{have/has} + participe passé : \eng{I have learnt} (et non \emph{I have learn} ni \emph{I am learn}).
\item Avec une date passée précise (\eng{last year}, \eng{in 2019}), utilisez le past simple, jamais le present perfect : \eng{Last year, I started my apprenticeship} (et non \emph{I have started}).
\item Après \eng{going to}, mettez la base verbale : \eng{I am going to become a technician} (et non \emph{to becoming}).
\end{itemize}
\end{attention}

\courssub{Exercice résolu : relier des phrases simples}

\begin{exoresolu}[Relier avec le connecteur approprié]
Énoncé --- Reliez chaque paire de phrases avec le connecteur qui convient parmi : \eng{however}, \eng{because}, \eng{therefore}, \eng{for example}, \eng{moreover}, \eng{although}.

1. I want to become a mechanic. I love engines.
2. The training is long. I am ready to work hard.
3. My father is a welder. He earns a good salary.
4. Many trades are useful. Plumbing and carpentry are always needed.
5. I have no money for tools. I will start with small repairs.
\tcblower
Solution détaillée :

1. \eng{I want to become a mechanic because I love engines.} (cause : la seconde phrase explique la première.)
2. \eng{The training is long; however, I am ready to work hard.} ou \eng{Although the training is long, I am ready to work hard.} (opposition entre la difficulté et la motivation.)
3. \eng{My father is a welder; moreover, he earns a good salary.} (addition d'un argument supplémentaire.)
4. \eng{Many trades are useful; for example, plumbing and carpentry are always needed.} (la seconde phrase illustre la première.)
5. \eng{I have no money for tools; therefore, I will start with small repairs.} (conséquence logique.)

Remarque : pour la phrase 2, on ne peut pas écrire \emph{Although the training is long, but I am ready} : \eng{although} et \eng{but} ne se combinent jamais.
\end{exoresolu}

\begin{exercice}[À vous de jouer]
Reliez les phrases suivantes avec un connecteur différent à chaque fois, puis indiquez le temps verbal à utiliser :

1. When I was ten, I (watch) my grandmother cooking. Since then, I (learn) many recipes.
2. I (love) cooking now. One day, I (open) a restaurant in Douala.
3. This job is demanding. It is rewarding.
\end{exercice}

\begin{corrige}[Exercice --- Relier et conjuguer]
1. \eng{When I was ten, I watched my grandmother cooking. Since then, I have learnt many recipes.} (past simple pour le souvenir daté ; present perfect pour l'expérience acquise, marqueur \eng{since then}.)
2. \eng{I love cooking now. Moreover, one day, I am going to open a restaurant in Douala.} (present simple pour le goût ; \eng{going to} pour le projet décidé.)
3. \eng{This job is demanding; however, it is rewarding.} ou \eng{Although this job is demanding, it is rewarding.} (opposition.)
\end{corrige}

\begin{aretenir}[L'essentiel de la section]
Une composition réussie combine \cle{connecteurs variés} (au moins 5 différents, un par famille) et \cle{temps verbaux cohérents} : past simple pour le souvenir, present perfect pour l'expérience, present simple pour la passion, \eng{going to} pour le projet. Retenez les pièges : \eng{that is why} (et non \emph{it is why}), \eng{for example} (et non \emph{by example}), et jamais \eng{but} après \eng{although}. Vous disposez maintenant de tous les outils linguistiques : la section suivante vous montrera comment les organiser dans la structure complète de la composition --- introduction, développement, conclusion.
\end{aretenir}

\coursec{La structure de la composition et le plan de rédaction}

Dans la section précédente, nous avons appris à relier nos idées grâce aux connecteurs logiques (\eng{first, moreover, as a result, finally}) et à choisir les bons temps verbaux. Mais des connecteurs ne suffisent pas : encore faut-il savoir \emph{où} placer chaque idée. Une bonne composition anglaise n'est pas un flux de pensées comme en français parlé ; c'est un \cle{édifice} construit selon un plan rigide et visible. C'est exactement ce que les correcteurs du Bac attendent. Nous allons maintenant apprendre à construire cet édifice, étage par étage.

\courssub{Observer d'abord : deux débuts de composition}

Lisez ces deux ouvertures sur le même sujet : \eng{Write about a trade or profession that interests you.}

\begin{textebox}[Two openings — compare them]
\textbf{Opening A:} “I like farming. Farming is good. My grandmother is a farmer. I will write about farming.”

\textbf{Opening B:} “Have you ever watched a tiny maize seed become a tall green plant heavy with corn? I have, and that miracle changed my life. Farming, the oldest profession in the world, is the trade I dream of practising. In this composition, I will explain where my interest comes from, how I am developing it, and how I hope to share it with other young people.”
\end{textebox}

\textbf{Observation.} Les deux textes annoncent le même sujet, mais l'ouverture B est infiniment meilleure. Pourquoi ? Elle contient les \cle{trois ingrédients obligatoires} d'une introduction anglaise : une \cle{accroche} (\emph{hook}) qui capte l'attention, la \cle{présentation} du métier choisi, et l'\cle{annonce du plan} (\emph{outline}). L'ouverture A, elle, répète le même mot, commence chaque phrase par \eng{I} et n'annonce rien. Retenez ce contraste : il résume toute cette section.

\courssub{Vue d'ensemble : les trois parties de la composition}

\begin{definition}[La structure en trois parties]
Toute composition anglaise au lycée (et au Bac) comporte trois parties :
\begin{itemize}
\item the \cle{introduction} : accroche + présentation du sujet + annonce du plan (3 à 4 phrases) ;
\item the \cle{body} (développement) : 2 ou 3 paragraphes, chacun consacré à \emph{une seule} idée ;
\item the \cle{conclusion} : résumé de l'idée principale + phrase d'ouverture (2 à 3 phrases).
\end{itemize}
Au Bac, on attend environ \cle{150 à 200 mots}, répartis en \cle{4 à 5 paragraphes}.
\end{definition}

\begin{attention}[Différence majeure avec le français]
En français, on apprend parfois que l'annonce du plan est « scolaire » ou facultative. En anglais, c'est l'\emph{inverse} : une introduction qui n'annonce pas le plan est jugée \emph{incomplète}. La phrase d'annonce (\eng{In this composition, I will explain...}) est un marqueur d'organisation que le correcteur cherche explicitement. Ne la supprimez jamais.
\end{attention}

\courssub{L'introduction : accroche, présentation, annonce du plan}

\begin{methode}[Rédiger l'introduction en trois étapes]
\begin{enumerate}
\item \textbf{Une phrase d'accroche} (\emph{hook}) : une question, un souvenir ou une citation. Exemple : \eng{Have you ever watched a mechanic bring a dead engine back to life?} « Avez-vous déjà vu un mécanicien ramener un moteur mort à la vie ? »
\item \textbf{Présenter le métier choisi} en une phrase générale : \eng{Motor mechanics is a trade that keeps our towns moving.} « La mécanique automobile est un métier qui fait avancer nos villes. »
\item \textbf{Annoncer ce que vous allez expliquer} : \eng{I will explain how this interest started, how I am developing it, and what my plans are.} « Je vais expliquer comment cet intérêt est né, comment je le développe et quels sont mes projets. »
\end{enumerate}
\end{methode}

\begin{exemplebox}[Trois types d'accroches]
\begin{itemize}
\item \textbf{Question :} \eng{Have you ever wondered who builds the houses we live in?} « Vous êtes-vous déjà demandé qui construit les maisons où nous vivons ? »
\item \textbf{Souvenir :} \eng{When I was ten, I watched my uncle repair a radio in his small workshop in Bamenda.} « Quand j'avais dix ans, j'ai regardé mon oncle réparer une radio dans son petit atelier de Bamenda. »
\item \textbf{Citation ou proverbe :} \eng{As the saying goes, “Give a man a fish and you feed him for a day; teach him to fish and you feed him for a lifetime.”} « Comme le dit le proverbe : donne un poisson à un homme et tu le nourris un jour ; apprends-lui à pêcher et tu le nourris toute sa vie. »
\end{itemize}
\end{exemplebox}

\courssub{Le développement : trois paragraphes, trois missions}

Pour notre thème (\eng{developing and sharing an interest in a trade}), le développement suit un plan en trois paragraphes très clair :

\begin{center}
\begin{tabularx}{\linewidth}{|l|X|X|}
\hline
\rowcolor{popblueL} \textbf{Paragraphe} & \textbf{Mission} & \textbf{Questions directrices} \\
\hline
Body 1 & \cle{Origine de l'intérêt} & Who inspired me? What memory do I have? When did it start? \\
\hline
Body 2 & \cle{Développement de l'intérêt} & What training do I follow? How do I practise? What do I read? \\
\hline
Body 3 & \cle{Partage et projets} & How will I help others? Will I open a workshop? What is my dream? \\
\hline
\end{tabularx}
\end{center}

\begin{propriete}[La loi d'or du paragraphe anglais]
Chaque paragraphe anglais obéit à une architecture fixe :
\begin{itemize}
\item \textbf{une \cle{topic sentence}} (phrase d'idée principale), placée \emph{en premier} : elle annonce le contenu du paragraphe ;
\item \textbf{deux ou trois \cle{supporting sentences}} (phrases de soutien) : détails, exemples, souvenirs, chiffres ;
\item éventuellement \textbf{une \cle{linking sentence}} (phrase de transition) qui prépare le paragraphe suivant.
\end{itemize}
Règle absolue : \emph{un paragraphe = une seule idée}. Si vous changez d'idée, changez de paragraphe.
\end{propriete}

\begin{exemplebox}[Topic sentence et supporting sentences — exemples traduits]
\textbf{Topic sentence :} \eng{My interest in farming began in my grandmother's village.} « Mon intérêt pour l'agriculture est né dans le village de ma grand-mère. »

\textbf{Supporting sentences :} \eng{Every holiday, I helped her plant maize and harvest cassava.} « Chaque vacances, je l'aidais à planter le maïs et à récolter le manioc. » \eng{Moreover, she taught me how to read the sky before the rains.} « De plus, elle m'a appris à lire le ciel avant les pluies. »

\textbf{Linking sentence :} \eng{These precious memories pushed me to take farming seriously.} « Ces précieux souvenirs m'ont poussé à prendre l'agriculture au sérieux. »
\end{exemplebox}

\begin{popfigure}
\begin{center}
\begin{tikzpicture}[node distance=0.5cm,
box/.style={draw=popblue, fill=popblueL, rounded corners=2pt, align=center, text width=5.2cm, inner sep=5pt, font=\small},
arr/.style={-{Stealth[length=2.5mm]}, thick, popdark}]
\node[box, fill=poporangeL, draw=poporange, align=center] (t){\textbf{TOPIC SENTENCE}\\ \eng{My interest in farming began in my grandmother's village.}};
\node[box, below=of t, align=center] (s1){\textbf{Supporting sentence 1}\\ \eng{Every holiday, I helped her plant maize.}};
\node[box, below=of s1, align=center] (s2){\textbf{Supporting sentence 2}\\ \eng{Moreover, she taught me to read the sky.}};
\node[box, below=of s2, fill=popgreenL, draw=popgreen, align=center] (l){\textbf{Linking sentence}\\ \eng{These memories pushed me to take farming seriously.}};
\draw[arr] (t) -- (s1);
\draw[arr] (s1) -- (s2);
\draw[arr] (s2) -- (l);
\end{tikzpicture}
\end{center}
\legende{Anatomie d'un paragraphe anglais : une idée principale, des preuves, une transition.}
\end{popfigure}

\courssub{La conclusion : résumé et ouverture}

\begin{methode}[Rédiger la conclusion en deux étapes]
\begin{enumerate}
\item \textbf{Reformuler l'idée principale avec d'autres mots} (ne copiez jamais l'introduction) : \eng{In short, farming is not just a job for me; it is a passion born in my family.} « En bref, l'agriculture n'est pas seulement un métier pour moi ; c'est une passion née dans ma famille. »
\item \textbf{Ajouter une phrase d'ouverture} : un conseil aux jeunes ou un espoir pour l'avenir : \eng{I hope every young Cameroonian will find a trade they love.} « J'espère que chaque jeune Camerounais trouvera un métier qu'il aime. »
\end{enumerate}
\end{methode}

\begin{attention}[Erreurs fréquentes des francophones]
\begin{itemize}
\item \textbf{Ne commencez pas tous les paragraphes par \eng{I}.} Variez les sujets : \eng{This trade requires patience.} / \eng{My training at the technical college is demanding.} / \eng{Young people often underestimate farming.}
\item Pas de \textbf{nouvelle idée} dans la conclusion : elle résume, elle n'ajoute rien.
\item Attention aux calques du français : on écrit \eng{in conclusion} ou \eng{to conclude}, jamais \eng{in the conclusion}.
\end{itemize}
\end{attention}

\courssub{La technique du brainstorming : de l'idée au plan}

Avant de rédiger, le bon élève passe par quatre étapes de préparation. C'est ce chemin que récompense le correcteur.

\begin{enumerate}
\item \textbf{Liste d'idées} (\emph{brainstorming}) : notez en vrac, en 3 minutes, tout ce qui concerne le sujet : souvenirs, personnes, lieux, projets, vocabulaire.
\item \textbf{Regroupement} : classez les idées par familles (origine / formation / projets).
\item \textbf{Classement} : numérotez les familles dans l'ordre logique (chronologique ou du passé vers l'avenir).
\item \textbf{Plan numéroté} : transformez chaque famille en paragraphe avec sa topic sentence.
\end{enumerate}

\begin{popfigure}
\begin{center}
\begin{tikzpicture}[node distance=0.55cm and 0.7cm,
stage/.style={draw=popblue, fill=popblueL, rounded corners=3pt, align=center, text width=3.1cm, inner sep=5pt, font=\small},
mini/.style={draw=poppurple, fill=poppurpleL, rounded corners=2pt, align=center, font=\footnotesize, inner sep=2pt},
arr/.style={-{Stealth[length=2.5mm]}, thick, popdark}]
\node[stage, fill=popgoldL, draw=popgold, align=center] (a){\textbf{SUJET}\\ \eng{My interest in farming}};
\node[stage, right=of a, align=center] (b){\textbf{BRAINSTORMING}\\ liste d'idées en vrac};
\node[stage, right=of b, align=center] (c){\textbf{PLAN}\\ Intro / P1 / P2 / P3 / Concl.};
\node[stage, right=of c, align=center] (d){\textbf{RÉDACTION}\\ 150--200 mots};
\node[stage, right=of d, fill=popgreenL, draw=popgreen, align=center] (e){\textbf{RELECTURE}\\ correction};
\draw[arr] (a) -- (b);
\draw[arr] (b) -- (c);
\draw[arr] (c) -- (d);
\draw[arr] (d) -- (e);
\node[mini, below=0.35cm of c] (p1) {Intro : hook + plan};
\node[mini, below=0.15cm of p1] (p2) {P1 : origine};
\node[mini, below=0.15cm of p2] (p3) {P2 : formation};
\node[mini, below=0.15cm of p3] (p4) {P3 : partage};
\node[mini, below=0.15cm of p4] (p5) {Concl. : résumé};
\end{tikzpicture}
\end{center}
\legende{La chaîne de production d'une composition : du sujet à la relecture.}
\end{popfigure}

\courssub{Exemple de plan détaillé commenté : “My interest in farming”}

\begin{exemplebox}[Un plan complet, prêt à rédiger]
\textbf{Introduction} — Hook : question sur le miracle d'une graine de maïs. Présentation : farming, the oldest profession. Annonce : origin, development, plans.

\textbf{Paragraph 1 — Origin.} Topic sentence : \eng{My interest in farming began in my grandmother's village.} Supports : planting maize, harvesting cassava, her advice.

\textbf{Paragraph 2 — Development.} Topic sentence : \eng{Since then, I have worked hard to develop this passion.} Supports : agricultural science at school, a small garden behind our house in Yaoundé, reading about modern techniques.

\textbf{Paragraph 3 — Sharing and plans.} Topic sentence : \eng{In the future, I want to share this interest with my community.} Supports : teaching young people, opening a modern farm, creating jobs.

\textbf{Conclusion} — Reformulation : farming is a passion, not just a job. Ouverture : hope for young Cameroonians.
\end{exemplebox}

Remarquez la logique : \emph{passé} (origine) $\rightarrow$ \emph{présent} (formation) $\rightarrow$ \emph{futur} (projets). Cette progression chronologique est la plus sûre pour ce type de sujet.

\begin{savaistu}[Le plan, votre meilleur investissement au Bac]
Au Baccalauréat, les correcteurs attribuent des points \emph{séparés} pour le respect du sujet, l'organisation, le vocabulaire et la grammaire. Conséquence précieuse : un plan clair avec introduction annoncée, paragraphes distincts et conclusion garantit déjà une grande partie des points d'organisation et de respect du sujet — avant même que la grammaire soit jugée ! Cinq minutes de planification peuvent donc valoir la moitié de la note.
\end{savaistu}

\begin{exoresolu}[Construire un plan à partir d'un brainstorming]
\textbf{Énoncé.} Voici les idées notées en vrac par une élève sur le sujet \eng{My interest in tailoring} : \eng{my mother's sewing machine / open a fashion shop in Douala / watching her sew dresses / lessons at a training centre / teach my little sister / practise every Saturday / a famous tailor in our quarter / buy beautiful fabrics}. Regroupez-les en trois familles et proposez une topic sentence pour chaque paragraphe.
\tcblower
\textbf{Solution détaillée.}
\begin{itemize}
\item \textbf{Famille 1 — Origine :} my mother's sewing machine ; watching her sew dresses ; a famous tailor in our quarter. Topic sentence : \eng{My passion for tailoring was born at home, beside my mother's sewing machine.}
\item \textbf{Famille 2 — Développement :} lessons at a training centre ; practise every Saturday ; buy beautiful fabrics. Topic sentence : \eng{To turn this passion into a skill, I am following serious training.}
\item \textbf{Famille 3 — Partage et projets :} open a fashion shop in Douala ; teach my little sister. Topic sentence : \eng{One day, I hope to share my skills and open my own fashion shop in Douala.}
\end{itemize}
L'ordre passé $\rightarrow$ présent $\rightarrow$ futur donne un plan logique et facile à rédiger.
\end{exoresolu}

\begin{exercice}[À vous de jouer]
Choisissez un métier (\eng{carpentry, nursing, computing, hairdressing...}). Faites un brainstorming de huit idées, regroupez-les en trois familles, puis rédigez les trois topic sentences et la phrase d'annonce du plan de l'introduction. Vous rédigerez la composition complète dans la section 6.
\end{exercice}

\begin{aretenir}[L'essentiel de la section]
\begin{itemize}
\item Une composition = introduction (accroche + présentation + annonce du plan) + développement (3 paragraphes : origine, développement, partage) + conclusion (résumé + ouverture).
\item Chaque paragraphe = une topic sentence + 2 ou 3 supporting sentences (+ linking sentence).
\item Toujours passer par : brainstorming $\rightarrow$ regroupement $\rightarrow$ plan numéroté $\rightarrow$ rédaction $\rightarrow$ relecture.
\item Au Bac : 150 à 200 mots, 4 à 5 paragraphes ; le plan rapporte des points à lui tout seul.
\end{itemize}
\end{aretenir}

\coursec{Modèle commenté et erreurs à éviter}

Dans la section précédente, nous avons construit le \cle{plan} de notre composition : une introduction qui accroche le lecteur et annonce le sujet, un développement en deux ou trois paragraphes reliés par des connecteurs, et une conclusion qui répond à la question posée. Il est temps maintenant de voir ce plan \emph{en action}. Nous allons lire ensemble une composition complète, la décortiquer phrase par phrase, puis apprendre à repérer et à corriger les cinq erreurs les plus fréquentes des élèves francophones. Objectif : que votre propre texte ressemble au modèle, et non au contre-modèle !

\courssub{Un modèle de composition : analyse guidée}

\begin{textebox}[Model composition — Why I want to become a nurse]
“Have you ever seen a nurse save a life? I have, and since that day, my dream has been to wear the white uniform.

First of all, I discovered this profession when my little brother was hospitalised in Douala. The nurses were kind and skilful. They comforted him and explained everything to our family.

Moreover, I have already learnt first aid at school, and I often help at the health centre during the holidays.

However, I know that nursing requires long studies; therefore, I am working hard in science subjects.

In conclusion, nursing is more than a job for me: it is a way of serving my community, and I am ready to make sacrifices for it.”
\end{textebox}

\textbf{Glossaire}

\begin{center}
\begin{tabularx}{\linewidth}{|l|l|X|}
\hline
\rowcolor{popblueL} Mot anglais & Nature & Traduction française \\
\hline
to save a life & verbe & sauver une vie \\
\hline
since that day & locution & depuis ce jour-là \\
\hline
skilful & adjectif & habile, compétent \\
\hline
to comfort & verbe & réconforter \\
\hline
first aid & nom & les premiers secours \\
\hline
health centre & nom & centre de santé \\
\hline
to require & verbe & exiger, nécessiter \\
\hline
a way of serving & nom + gérondif & une manière de servir \\
\hline
sacrifice & nom & sacrifice \\
\hline
\end{tabularx}
\end{center}

\courssub{Lecture commentée : les marges du texte}

Observons maintenant chaque phrase et demandons-nous : \emph{que fait l'élève à cet endroit précis ?} C'est exactement ce que l'examinateur regarde.

\begin{exemplebox}[Analyse ligne par ligne]
\begin{itemize}
\item \eng{“Have you ever seen a nurse save a life?”} → \cle{l'accroche} : une question directe au lecteur, au \emph{present perfect}, qui crée immédiatement l'intérêt. « Avez-vous déjà vu une infirmière sauver une vie ? »
\item \eng{“I have, and since that day, my dream has been to wear the white uniform.”} → \cle{l'annonce du sujet} : réponse courte (\eng{I have}) puis présentation du thème au \emph{present perfect} (\eng{has been}) : un rêve qui a commencé dans le passé et continue aujourd'hui.
\item \eng{“First of all, I discovered this profession when my little brother was hospitalised in Douala.”} → \cle{past simple} pour raconter \emph{l'origine} de l'intérêt + connecteur d'ordre \eng{first of all}. « J'ai découvert ce métier quand mon petit frère a été hospitalisé à Douala. »
\item \eng{“The nurses were kind and skilful.”} → description au past simple : deux adjectifs précis, pas de phrase vide.
\item \eng{“Moreover, I have already learnt first aid at school, and I often help at the health centre during the holidays.”} → \cle{connecteur d'addition} \eng{moreover} + \emph{present perfect} (\eng{have learnt}) pour une expérience acquise + \emph{present simple} (\eng{help}) pour une habitude. « De plus, j'ai déjà appris les premiers secours à l'école, et j'aide souvent au centre de santé pendant les vacances. »
\item \eng{“However, I know that nursing requires long studies; therefore, I am working hard in science subjects.”} → \cle{connecteur d'opposition} \eng{however} + \cle{connecteur de conséquence} \eng{therefore} : l'élève montre qu'il est lucide et qu'il a un \emph{projet} (\emph{present continuous} : \eng{I am working}). « Cependant, je sais que le métier d'infirmier exige de longues études ; c'est pourquoi je travaille dur dans les matières scientifiques. »
\item \eng{“In conclusion, nursing is more than a job for me: it is a way of serving my community.”} → \cle{la conclusion} : elle reprend l'idée principale sans introduire d'idée nouvelle, et s'achève sur une \emph{ouverture} (servir la communauté). « En conclusion, le métier d'infirmier est plus qu'un emploi pour moi : c'est une manière de servir ma communauté. »
\end{itemize}
\end{exemplebox}

\begin{exemplebox}[Exemples traduits à retenir]
\eng{Nursing is more than a job: it is a vocation.} « Le métier d'infirmier est plus qu'un emploi : c'est une vocation. »

\eng{I often help at the health centre.} « J'aide souvent au centre de santé. »

\eng{Since that day, my dream has been to become a mechanic.} « Depuis ce jour-là, mon rêve est de devenir mécanicien. »

\eng{My uncle, who is a carpenter in Bamenda, taught me to love wood.} « Mon oncle, qui est menuisier à Bamenda, m'a appris à aimer le bois. »
\end{exemplebox}

\begin{aretenir}[La recette du modèle]
Le modèle réussit parce qu'il respecte quatre règles simples :
\begin{enumerate}
\item chaque paragraphe a \textbf{une seule fonction} (accroche / origine / arguments / projet / conclusion) ;
\item chaque changement d'idée est annoncé par un \textbf{connecteur} (\eng{first of all, moreover, however, therefore, in conclusion}) ;
\item les \textbf{temps verbaux} sont choisis logiquement : past simple pour le souvenir, present perfect pour l'expérience, present simple pour les habitudes et les vérités, present continuous pour le projet en cours ;
\item l'élève parle de \textbf{lui-même} (\eng{I, my, me}) : il ne décrit pas le métier en général, il explique \emph{son} intérêt personnel.
\end{enumerate}
\end{aretenir}

\courssub{Les cinq erreurs à éviter absolument}

\textbf{Erreur 1 : le hors-sujet}

Le sujet n'est pas \eng{Describe the nursing profession} mais \eng{Explain why YOU are interested in it}. Une composition qui énumère les tâches d'un infirmier sans jamais dire \eng{I, my, me} est un hors-sujet partiel. Chaque paragraphe doit contenir au moins une phrase à la première personne.

\begin{exemplebox}[Hors-sujet ou pas ?]
\eng{Nurses work in hospitals. They wear uniforms and give injections.} → description générale : \textbf{hors-sujet}.

\eng{I admire nurses because they cared for my grandmother; that is why I want to join them.} → intérêt personnel : \textbf{dans le sujet}.
\end{exemplebox}

\textbf{Erreur 2 : la traduction mot à mot du français}

\begin{attention}[Les calques du français]
\eng{I have 17 years} est un calque de « j'ai 17 ans ». En anglais, l'âge se dit avec \emph{to be} : \eng{I am 17 years old.}

Autres calques fréquents :
\begin{itemize}
\item « je veux faire médecine » → \eng{I want to make medicine} est faux ; dites \eng{I want to study medicine.}
\item « j'ai appris ce métier » → \eng{I have learned this trade} (ou \eng{I trained for this job}) ; attention, le participe passé de \emph{teach} est \emph{taught}, on n'écrit jamais *\eng{I have teached}.
\item « je suis intéressé par » → \eng{I am interested by} est un calque ; dites \eng{I am interested \textbf{in} nursing.}
\end{itemize}
\end{attention}

\textbf{Erreur 3 : l'absence de connecteurs (ou la répétition de « and »)}

Un texte du type \eng{I like nursing and nurses are kind and I helped my brother and I want to study science and...} fatigue le lecteur et perd des points. Remplacez les \eng{and} en chaîne par des connecteurs variés :

\begin{center}
\begin{tabularx}{\linewidth}{|l|X|X|}
\hline
\rowcolor{popblueL} Fonction & Connecteur anglais & Équivalent français \\
\hline
ordre & first of all, then, finally & d'abord, ensuite, enfin \\
\hline
addition & moreover, furthermore, also & de plus, en outre, aussi \\
\hline
opposition & however, whereas, although & cependant, tandis que, bien que \\
\hline
cause & because, since, as & parce que, puisque \\
\hline
conséquence & therefore, so, that is why & c'est pourquoi, donc \\
\hline
conclusion & in conclusion, to sum up & en conclusion, pour résumer \\
\hline
\end{tabularx}
\end{center}

\textbf{Erreur 4 : le mélange des temps sans raison et l'oubli du -s}

\begin{attention}[La 3e personne du singulier]
N'oubliez pas le \cle{-s} au present simple : \eng{My sister work\textbf{s} in a garage} (et non \eng{my sister work}). De même : \eng{He want\textbf{s} to be a pilot}, \eng{This job require\textbf{s} patience}.

Ne changez pas de temps sans raison : si vous racontez un souvenir au past simple, restez-y dans tout le paragraphe : \eng{I visited the workshop, I watched the mechanics and I asked questions} — et non \eng{I visited... I watch... I am asking...}.
\end{attention}

\textbf{Erreur 5 : pas de conclusion, ou une conclusion qui introduit une idée nouvelle}

La conclusion doit \emph{reformuler} l'idée principale, pas ouvrir un nouveau débat. Terminer par \eng{In conclusion, I also want to become a footballer} détruit tout le texte. Une bonne formule de clôture : \eng{In conclusion, + reprise de l'idée + ouverture personnelle}.

\begin{popfigure}
\begin{center}
\begin{tabular}{|p{5.6cm}|p{5.6cm}|}
\hline
\rowcolor{poporangeL} \textbf{Erreur fréquente} & \textbf{Correction} \\
\hline
\eng{I have 17 years.} & \eng{I am 17 years old.} \\
\hline
\eng{I want to make medicine.} & \eng{I want to study medicine.} \\
\hline
\eng{My sister work in a garage.} & \eng{My sister works in a garage.} \\
\hline
\eng{I like nursing and nurses are kind and I helped my brother.} & \eng{I like nursing because nurses are kind; moreover, I once helped my brother.} \\
\hline
\eng{In conclusion, I also love football.} & \eng{In conclusion, nursing is my dream because I want to serve my community.} \\
\hline
\end{tabular}
\end{center}
\legende{Les cinq erreurs classiques et leurs corrections}
\end{popfigure}

\courssub{Correction collective d'un texte d'élève}

\begin{exoresolu}[Trouvez et corrigez les 8 erreurs]
Énoncé. Voici le brouillon d'un élève de Première. Il contient \textbf{8 erreurs} (temps, calques, connecteurs, 3e personne, conclusion). Soulignez-les et corrigez-les.

\eng{I have 16 years and I want to become a mechanic. Last year I visit my uncle's garage in Douala and I am very impressed. The mechanics was kind and they explain everything to me. Moreover I often repair bicycles for my friends. My father work as a driver and he say that mechanics is a good job. I want to make mechanics at the technical school. In conclusion, I also like cooking very much.}

\tcblower

Solution détaillée — les 8 erreurs corrigées :

\begin{enumerate}
\item \eng{I have 16 years} → \eng{I \textbf{am} 16 years \textbf{old}} (calque du français « j'ai 16 ans »).
\item \eng{Last year I visit} → \eng{Last year I \textbf{visited}} (souvenir daté : past simple).
\item \eng{I am very impressed} → \eng{I \textbf{was} very impressed} (on reste au past simple dans le récit).
\item \eng{The mechanics was kind} → \eng{The mechanics \textbf{were} kind} (sujet pluriel).
\item \eng{they explain} → \eng{they \textbf{explained}} (cohérence des temps dans le récit).
\item \eng{My father work} → \eng{My father \textbf{works}} (-s de la 3e personne du singulier).
\item \eng{he say} → \eng{he \textbf{says}} (-s obligatoire ; attention, \eng{says} se prononce « sèz »).
\item \eng{I want to make mechanics} → \eng{I want to \textbf{study} mechanics} (calque de « faire mécanique »).
\end{enumerate}

Et le problème de fond : la phrase finale \eng{In conclusion, I also like cooking very much} introduit une \textbf{idée nouvelle} dans la conclusion. Il faut la remplacer, par exemple : \eng{In conclusion, mechanics is more than a job for me: it is a passion I want to turn into my future profession.}

Texte corrigé complet :

\eng{I am 16 years old and I want to become a mechanic. Last year I visited my uncle's garage in Douala and I was very impressed. The mechanics were kind and they explained everything to me. Moreover, I often repair bicycles for my friends. My father works as a driver and he says that mechanics is a good job. I want to study mechanics at the technical school. In conclusion, mechanics is more than a job for me: it is a passion I want to turn into my future profession.}
\end{exoresolu}

\courssub{La relecture : la méthode C.O.P.S.}

\begin{methode}[Relire sa composition avec la checklist C.O.P.S.]
Avant de rendre votre copie, relisez-la quatre fois, en cherchant chaque fois un seul type de problème :
\begin{enumerate}
\item \textbf{C — Capital letters} : majuscule en début de phrase, au pronom \eng{I}, aux noms propres (\eng{Douala, Cameroon}), aux jours et aux mois (\eng{Monday, July}).
\item \textbf{O — Organisation} : introduction + 2 ou 3 paragraphes de développement + conclusion ; au moins 4 connecteurs différents ; chaque paragraphe parle de \emph{votre} intérêt personnel.
\item \textbf{P — Punctuation} : un point en fin de phrase, une virgule après les connecteurs en tête de phrase (\eng{Moreover, ...}), pas de point-virgule au hasard.
\item \textbf{S — Spelling} : orthographe des mots du vocabulaire (\eng{nurse, skilful, profession, experience}), verbes irréguliers corrects (\eng{learnt, taught, wore}).
\end{enumerate}
Ajoutez enfin deux vérifications spéciales anglais : les \textbf{temps verbaux} (un seul temps par paragraphe, sauf raison claire) et le \textbf{-s} de la 3e personne du singulier. Comptez vos mots : environ 180 mots sont attendus.
\end{methode}

\begin{exercice}[À vous de jouer]
Relisez ce début de composition avec la méthode C.O.P.S. et réécrivez-le correctement :

\eng{i want to be a teacher. because my mother is a teacher and she love her job and she help many children and i want to do the same.}

Vous devriez trouver : une majuscule, un connecteur mal placé en début de phrase, deux -s manquants et une répétition de \eng{and}.
\end{exercice}

\begin{corrige}[Exercice — relecture C.O.P.S.]
\eng{\textbf{I} want to be a teacher\textbf{ because} my mother is a teacher\textbf{.} She love\textbf{s} her job and she help\textbf{s} many children\textbf{; therefore,} I want to do the same.}

Corrections : majuscule à \eng{I} ; on ne commence pas une phrase par \eng{because} seul ; \eng{she loves}, \eng{she helps} ; le deuxième \eng{and} est remplacé par le connecteur de conséquence \eng{therefore}.
\end{corrige}

\begin{aretenir}[L'essentiel de la section]
Une bonne composition n'est pas un texte parfait du premier coup : c'est un texte \textbf{relu}. Retenez le trio gagnant : parlez de \emph{vous} (\eng{I, my}), reliez vos idées (\eng{moreover, however, therefore}), et relisez avec \cle{C.O.P.S.}. Vous êtes maintenant prêt pour la dernière étape : la grille d'évaluation et votre propre production guidée.
\end{aretenir}

\coursec{Grille d'évaluation et production guidée}

Dans la section précédente, nous avons commenté un modèle de composition et repéré les erreurs à éviter. Vous savez maintenant ce qu'est une \emph{bonne} composition. Mais comment le correcteur décide-t-il de votre note ? Et surtout, comment pouvez-vous vous corriger vous-même avant de rendre votre copie ? C'est l'objet de cette dernière section : comprendre la \cle{grille d'évaluation}, apprendre à s'auto-évaluer et à évaluer un camarade, puis produire votre propre composition guidée, étape par étape.

\courssub{La grille d'évaluation critériée}

\begin{definition}[Grille critériée]
Une \cle{grille d'évaluation critériée} (\eng{marking scheme} ou \eng{rubric}) est un tableau qui décompose la note finale en plusieurs \cle{critères} précis : contenu, organisation, vocabulaire, grammaire, orthographe. Elle rend la correction \emph{objective} : deux correcteurs différents, avec la même grille, attribuent presque la même note à la même copie.
\end{definition}

Observez d'abord comment un correcteur lit votre copie. Il ne lit pas seulement « si c'est beau » : il se pose cinq questions, dans cet ordre :

\begin{enumerate}
\item \eng{Did the candidate answer the question?} — Le candidat a-t-il répondu au sujet ?
\item \eng{Is the composition well organised?} — Y a-t-il une introduction, des paragraphes, une conclusion ?
\item \eng{Is the vocabulary rich and accurate?} — Le vocabulaire est-il varié et précis ?
\item \eng{Is the grammar correct?} — Les temps et les structures sont-ils corrects ?
\item \eng{What about spelling and punctuation?} — L'orthographe et la ponctuation sont-elles soignées ?
\end{enumerate}

\begin{propriete}[Grille d'évaluation — composition sur 20 points]
\begin{itemize}
\item \textbf{Contenu et respect du sujet} (\eng{content and task fulfilment}) : \textbf{/5}. Toutes les parties de la consigne sont traitées (le métier, l'origine de l'intérêt, le partage de cet intérêt).
\item \textbf{Organisation} (\eng{organisation}) : \textbf{/5}. Introduction, paragraphes distincts, conclusion, connecteurs logiques.
\item \textbf{Vocabulaire riche et précis} (\eng{vocabulary}) : \textbf{/4}. Champs lexicaux du métier et de la formation, pas de répétitions.
\item \textbf{Grammaire et temps corrects} (\eng{grammar and tenses}) : \textbf{/4}. Prétérit pour le passé, présent pour l'intérêt actuel, futur pour les projets ; accords sujet-verbe.
\item \textbf{Orthographe et ponctuation} (\eng{spelling and punctuation}) : \textbf{/2}.
\end{itemize}
\textbf{Total : /20.}
\end{propriete}

\begin{popfigure}
\begin{center}
\begin{tikzpicture}[x=3.1cm,y=0.72cm]
\fill[popblue] (0,5) rectangle (2.6,6);
\node[white,font=\bfseries,align=center] at (1.3,5.5) {Critère};
\fill[popgreen] (2.6,5) rectangle (3.9,6);
\node[white,font=\bfseries,align=center] at (3.25,5.5) {Excellent};
\fill[popblue!70] (3.9,5) rectangle (5.2,6);
\node[white,font=\bfseries,align=center] at (4.55,5.5) {Bon};
\fill[poporange] (5.2,5) rectangle (6.5,6);
\node[white,font=\bfseries,align=center] at (5.85,5.5) {Moyen};
\fill[poppink] (6.5,5) rectangle (7.8,6);
\node[white,font=\bfseries,align=center] at (7.15,5.5) {Faible};
\foreach \y/\crit/\e/\b/\m/\f in {
 4/{Contenu et sujet (/5)}/{5}/{4}/{3}/{0--2},
 3/{Organisation (/5)}/{5}/{4}/{3}/{0--2},
 2/{Vocabulaire (/4)}/{4}/{3}/{2}/{0--1},
 1/{Grammaire et temps (/4)}/{4}/{3}/{2}/{0--1},
 0/{Orthographe, ponct. (/2)}/{2}/{1.5}/{1}/{0}}{
 \fill[popblueL] (0,\y) rectangle (2.6,\y+1);
 \node[align=left,font=\small,anchor=west] at (0.08,\y+0.5) {\crit};
 \node[font=\small] at (3.25,\y+0.5) {\e\ pts};
 \node[font=\small] at (4.55,\y+0.5) {\b\ pts};
 \node[font=\small] at (5.85,\y+0.5) {\m\ pts};
 \node[font=\small] at (7.15,\y+0.5) {\f\ pts};
}
\draw[popdark,thick] (0,0) grid[xstep=1.3cm,ystep=0.72cm] (7.8,6);
\draw[popdark,very thick] (0,0) rectangle (7.8,6);
\end{tikzpicture}
\end{center}
\legende{Grille d'évaluation critériée : chaque critère est noté selon quatre niveaux de maîtrise.}
\end{popfigure}

\begin{exemplebox}[Lire la grille en pratique]
Une élève traite bien tout le sujet (contenu : 5/5), mais écrit sa composition en un seul bloc sans paragraphes (organisation : 2/5), avec un vocabulaire correct (3/4), une grammaire bonne (3/4) et une orthographe soignée (2/2). Total : \textbf{15/20}. Sans le problème de paragraphes, elle aurait eu 18/20. L'organisation coûte cher !
\end{exemplebox}

\begin{attention}[Le piège du « bloc unique »]
Une copie rédigée en \textbf{un seul paragraphe} perd des points d'organisation, \emph{même si les idées sont excellentes}. Le correcteur ne peut pas voir votre plan. \textbf{Sautez une ligne} entre l'introduction, chaque paragraphe de développement et la conclusion. En anglais, on peut aussi commencer chaque paragraphe par un léger retrait (\eng{indent}).
\end{attention}

\begin{savaistu}[Le WASSCE en Afrique anglophone]
Au Ghana et au Nigeria, les élèves passent le \emph{WASSCE} (\eng{West African Senior School Certificate Examination}). Les compositions y sont corrigées avec une grille presque identique à la nôtre : \eng{content} (contenu), \eng{organisation} (organisation), \eng{expression} (vocabulaire et grammaire) et \eng{mechanical accuracy} (orthographe, ponctuation). Les compétences que vous entraînez ici sont donc exactement celles exigées dans toute l'Afrique de l'Ouest anglophone — utile si vous étudiez ou travaillez un jour à Lagos ou à Accra !
\end{savaistu}

\courssub{Auto-évaluation et évaluation croisée}

Savoir se corriger soi-même est une compétence de champion. En anglais, on parle de \eng{self-assessment} (auto-évaluation) et de \eng{peer assessment} (évaluation entre pairs).

\begin{methode}[Auto-évaluation : corriger sa propre copie]
\begin{enumerate}
\item Relisez votre composition \textbf{une fois par critère} : une lecture pour le contenu, une pour l'organisation, une pour le vocabulaire, une pour la grammaire, une pour l'orthographe.
\item Pour chaque critère, cochez le niveau atteint (\emph{excellent}, \emph{bon}, \emph{moyen}, \emph{faible}) sur la grille.
\item Soulignez en rouge \textbf{une seule catégorie d'erreurs à la fois} : d'abord les verbes, puis les prépositions, puis l'orthographe.
\item Notez votre total sur 20, puis choisissez \textbf{un seul point à améliorer} pour la prochaine composition.
\end{enumerate}
\end{methode}

\begin{experience}[Évaluation croisée entre voisins]
Échangez votre copie avec votre voisin. Chacun note la copie de l'autre avec la grille, puis écrit \textbf{un commentaire bienveillant} en deux parties : d'abord un point fort (\eng{What I liked: ...}), ensuite une suggestion (\eng{To improve: ...}). Exemple de commentaire : \eng{What I liked: your introduction is clear and your past tenses are correct. To improve: add more linking words between paragraphs.} Rendez ensuite la copie à son auteur, qui compare la note du voisin avec sa propre auto-évaluation.
\end{experience}

\begin{exemplebox}[Consignes type Bac : comprendre ce qu'on vous demande]
\eng{Explain how you developed this interest.} « Expliquez comment vous avez développé cet intérêt. » → temps attendu : \textbf{prétérit}.\par
\eng{Share your plans for the future.} « Partagez vos projets d'avenir. » → temps attendu : \textbf{futur} (\eng{be going to}, \eng{will}, \eng{want to}).\par
\eng{Describe the job and its importance.} « Décrivez le métier et son importance. » → temps attendu : \textbf{présent simple}.
\end{exemplebox}

\begin{center}
\begin{tabularx}{\linewidth}{|X|X|X|}
\hline
\rowcolor{popblueL} Consigne anglaise & Traduction française & Temps attendu \\
\hline
\eng{Explain how you developed this interest.} & Expliquez comment vous avez développé cet intérêt. & Prétérit \\
\hline
\eng{Say why this job attracts you.} & Dites pourquoi ce métier vous attire. & Présent \\
\hline
\eng{Share your plans for the future.} & Partagez vos projets d'avenir. & Futur \\
\hline
\eng{Describe a typical day in this job.} & Décrivez une journée typique de ce métier. & Présent \\
\hline
\end{tabularx}
\end{center}

\courssub{Production guidée : à vous d'écrire !}

\begin{textebox}[Your task — Sujet de production]
\eng{Write a composition (150--200 words) about a trade, job or profession you are interested in. Explain how you developed this interest and how you share it or plan to share it.}\par
« Rédigez une composition (150 à 200 mots) sur un métier, un emploi ou une profession qui vous intéresse. Expliquez comment vous avez développé cet intérêt et comment vous le partagez — ou comptez le partager. »\par
\textbf{Glossaire :} \emph{trade} (nom) : métier manuel ; \emph{to develop an interest} : développer un intérêt ; \emph{to share} : partager ; \emph{to plan to} : avoir l'intention de.
\end{textebox}

\begin{methode}[Les quatre étapes imposées (40 minutes)]
\begin{enumerate}
\item \textbf{Brainstorming (5 min)} : notez au brouillon tous les mots qui viennent — le métier, les personnes qui vous ont inspiré, les lieux, les verbes d'action.
\item \textbf{Plan (5 min)} : organisez vos idées en trois parties — introduction (le métier choisi), développement (origine de l'intérêt au prétérit, partage au présent), conclusion (projets au futur).
\item \textbf{Rédaction (25 min)} : écrivez en suivant votre plan, en sautant une ligne entre les paragraphes et en utilisant au moins quatre connecteurs.
\item \textbf{Relecture avec la checklist (5 min)} : vérifiez chaque critère de la grille, un par un.
\end{enumerate}
\end{methode}

\begin{aretenir}[Aides autorisées pendant la rédaction]
\begin{itemize}
\item Le \textbf{tableau des connecteurs} : \eng{first, then, moreover, however, as a result, in conclusion}.
\item Les \textbf{champs lexicaux} des métiers et de la formation : \eng{apprenticeship, workshop, training, skills, to learn a trade, to run a business}.
\item Le \textbf{modèle de plan} : introduction → développement (2 paragraphes) → conclusion.
\end{itemize}
Ces aides sont vos outils d'entraînement : au Bac, vous les aurez mémorisés !
\end{aretenir}

\begin{methode}[Réussir le jour du Bac]
\begin{enumerate}
\item \textbf{Lisez le sujet deux fois} avant d'écrire quoi que ce soit.
\item \textbf{Soulignez les mots-clés} : le thème, le nombre de mots, les questions posées.
\item \textbf{Faites un plan au brouillon} : trois parties, une idée principale par paragraphe.
\item \textbf{Écrivez} sans chercher la perfection phrase par phrase — avancez.
\item \textbf{Gardez 5 minutes pour relire} : verbes, accords, orthographe, paragraphes.
\end{enumerate}
\end{methode}

\begin{exoresolu}[Appliquer la grille à une copie d'élève]
Énoncé — Voici un extrait de copie : \eng{I want to become a nurse. My aunt is a nurse in Bamenda and she inspired me. Last year I help her during the holidays. In the future, I will study nursing and I will share my knowledge with young people in my village.} Notez cet extrait sur le critère « grammaire et temps » (/4) et corrigez l'erreur.
\tcblower
Solution — L'extrait est clair et les temps sont globalement bien choisis (présent pour l'intérêt, futur pour les projets). Mais il contient une erreur de prétérit : \eng{Last year I \textbf{help} her} doit devenir \eng{Last year I \textbf{helped} her}, car \eng{last year} exige le prétérit. Note attribuée : \textbf{3/4} (une erreur de temps, sinon très bon). Correction complète : \eng{Last year I helped her during the holidays.} « L'année dernière, je l'ai aidée pendant les vacances. »
\end{exoresolu}

\begin{exercice}[À votre tour]
1. Notez avec la grille cette phrase d'ouverture : \eng{Since my childhood, I have always been fascinated by mechanics because of my father garage in Douala.} Identifiez l'erreur et corrigez-la.\par
2. Rédigez votre composition (150--200 mots) en suivant les quatre étapes imposées, puis auto-évaluez-vous avec la grille avant de la remettre.
\end{exercice}

\begin{corrige}[Exercice — question 1]
Erreur : il manque le génitif saxon — \eng{my father garage} doit devenir \eng{my \textbf{father's} garage} (« le garage de mon père »). Phrase corrigée : \eng{Since my childhood, I have always been fascinated by mechanics because of my father's garage in Douala.} Le \eng{present perfect} (\eng{have always been}) est correct ici, car l'intérêt a commencé dans le passé et continue aujourd'hui. Barème indicatif : contenu bon, grammaire 3/4 à cause du génitif.
\end{corrige}

Vous disposez désormais de tous les outils du rédacteur efficace : une grille pour comprendre ce qu'on attend de vous, une méthode en quatre étapes pour organiser votre temps, et une checklist pour relire. La prochaine fois que vous écrirez, vous ne rendrez plus une copie « au feeling » : vous rendrez une copie \emph{construite}. \eng{Now, pick up your pen and write!}

\newpage

\coursec{Activité d'intégration}

\courssub{Objectifs de l'activité}

À la fin de cette activité, tu seras capable de :
\begin{itemize}
\item rédiger, en conditions réelles et en temps limité (40 minutes), une composition complète de 150 à 200 mots intitulée \eng{The trade or profession I dream of} ;
\item planifier ta production : brainstorming, plan en trois parties (introduction, développement, conclusion), rédaction, relecture ;
\item utiliser correctement les connecteurs logiques (\eng{first of all, moreover, however, finally}), les temps verbaux adaptés (présent simple, futur avec \eng{will} et \eng{be going to}) et le vocabulaire des métiers et de la formation ;
\item évaluer la copie d'un camarade à l'aide d'une grille sur 20 points et formuler un commentaire constructif.
\end{itemize}

\courssub{Situation}

Ton école, \eng{Government High School Buea}, publie chaque trimestre un magazine scolaire en anglais, \eng{The Horizon}. Pour le prochain numéro spécial « métiers », la rédaction lance un appel à contributions à toutes les classes de Première. Voici l'annonce parue dans le magazine :

\begin{textebox}[The Horizon — Call for articles, Special Issue: Careers]
Dear students,

The Horizon is preparing its special issue on careers and we need YOUR voice! This month, our question is: “The trade or profession I dream of”.

Write a composition of 150 to 200 words. Tell us:
\begin{itemize}
\item what job or trade you dream of and why it attracts you;
\item how you discovered this interest (a person, an event, a visit, a book);
\item what training or studies you will need;
\item how this job can help your community or your country.
\end{itemize}

The best compositions will be published in the magazine and displayed on the school notice board. Deadline: Friday, 4 p.m. Give your work to your English teacher.

The Editorial Team
\end{textebox}

\begin{definition}[Glossaire utile]
\begin{itemize}
\item \cle{trade} (n.) : métier manuel ou artisanal (menuisier, mécanicien, tailleur).
\item \cle{profession} (n.) : profession intellectuelle exigeant de longues études (médecin, avocat, ingénieur).
\item \cle{deadline} (n.) : date limite. Prononciation : « déd-laïne ».
\item \cle{notice board} (n.) : panneau d'affichage.
\item \cle{career} (n.) : carrière, parcours professionnel.
\end{itemize}
\end{definition}

\courssub{Consignes}

\textbf{Tâche 1 — Compréhension de l'appel à contributions (5 min).}
Lis attentivement l'annonce du magazine et réponds par écrit en anglais, par des phrases complètes :
\begin{enumerate}
\item What is the topic of the special issue?
\item How many words must your composition contain?
\item What are the four points you must write about?
\item Where will the best compositions be displayed?
\end{enumerate}

\textbf{Tâche 2 — Brainstorming (5 min).}
Sur une feuille de brouillon, note au moins huit mots ou expressions liés au métier de tes rêves : nom du métier, qualités requises, formation, lieux de travail, verbes d'action. Exemple : \eng{nurse — caring — patient — hospital — look after — medical school — save lives}.

\textbf{Tâche 3 — Plan de rédaction (5 min).}
Organise tes idées selon la structure en trois parties :
\begin{enumerate}
\item \eng{Introduction} : annonce le métier choisi et ton intérêt général.
\item \eng{Body} (2 ou 3 paragraphes) : l'origine de ton intérêt, les raisons de ton choix, la formation nécessaire, l'utilité pour la communauté.
\item \eng{Conclusion} : réaffirme ton ambition et ton projet d'avenir.
\end{enumerate}

\textbf{Tâche 4 — Rédaction (25 min).}
Rédige ta composition de 150 à 200 mots au propre. Utilise au moins quatre connecteurs différents et varie les temps : présent simple pour les faits et les goûts, \eng{will} ou \eng{be going to} pour tes projets.

\textbf{Tâche 5 — Relecture avec la checklist C.O.P.S. (5 min).}
Vérifie ta copie :
\begin{itemize}
\item \cle{C}apitalization : majuscules en début de phrase, au pronom \eng{I}, aux noms propres.
\item \cle{O}rganization : trois parties visibles, connecteurs, paragraphes.
\item \cle{P}unctuation : points, virgules, pas de phrase interminable.
\item \cle{S}pelling : orthographe et grammaire (accord sujet-verbe, \eng{-s} à la 3\up{e} personne du singulier).
\end{itemize}

\textbf{Tâche 6 — Évaluation croisée (10 min).}
Échange ta copie avec ton voisin. Note sa composition sur 20 avec la grille ci-dessous, puis écris en bas de la copie \textbf{un commentaire positif} (\eng{What I liked: ...}) et \textbf{un conseil d'amélioration} (\eng{To improve: ...}).

\begin{center}
\begin{tabularx}{\linewidth}{|X|c|}
\hline
\rowcolor{popblueL} Critère & Points \\
\hline
Respect du sujet et de la longueur (150--200 mots) & /4 \\
\hline
Structure : introduction, développement, conclusion & /4 \\
\hline
Connecteurs logiques et cohérence des idées & /4 \\
\hline
Correction grammaticale et temps verbaux & /4 \\
\hline
Vocabulaire riche et varié sur les métiers & /4 \\
\hline
\end{tabularx}
\end{center}

\textbf{Tâche 7 — Publication.} Les trois meilleures compositions de la classe sont lues à voix haute par leurs auteurs, puis affichées au mur de la classe comme pages du numéro spécial de \eng{The Horizon}.

\courssub{Ressources à mobiliser}

\begin{aretenir}[Boîte à outils de la rédaction]
\textbf{Vocabulaire :} \eng{a carpenter, a nurse, an engineer, a tailor, a farmer, a journalist, a mechanic, a teacher, a pilot, an accountant} ; \eng{to be passionate about, to dream of, to train as, to apply for, to earn a living, to serve my community}.

\textbf{Connecteurs :} \eng{first of all} (tout d'abord), \eng{moreover} (de plus), \eng{however} (cependant), \eng{for example} (par exemple), \eng{that is why} (c'est pourquoi), \eng{to conclude} (pour conclure).

\textbf{Temps :} présent simple pour les faits et les goûts — \eng{I love helping people.} ; \eng{be going to} pour un projet déjà décidé — \eng{I am going to study nursing.} ; \eng{will} pour une prédiction ou une promesse — \eng{I will open my own workshop one day.}

\textbf{Structure type :} Introduction (2--3 phrases) / Développement (2--3 paragraphes) / Conclusion (2 phrases).
\end{aretenir}

\begin{propriete}[Choisir entre will et be going to]
\begin{itemize}
\item \eng{be going to} + base verbale : projet déjà décidé, intention ferme. \eng{I am going to attend a technical college.} « Je vais entrer dans un collège technique. »
\item \eng{will} + base verbale : prédiction, promesse ou décision spontanée. \eng{I will succeed.} « Je réussirai. »
\item Forme négative : \eng{will not} (contraction \eng{won't}) ; \eng{I am not going to give up.}
\end{itemize}
\end{propriete}

\begin{methode}[Rédiger une composition en 5 étapes]
\begin{enumerate}
\item \textbf{Comprendre} la consigne : souligne le sujet, la longueur et les points obligatoires.
\item \textbf{Brainstormer} : liste des mots-clés (métier, qualités, formation, utilité).
\item \textbf{Planifier} : répartis tes idées en introduction, développement, conclusion.
\item \textbf{Rédiger} : écris des phrases simples et correctes, reliées par des connecteurs.
\item \textbf{Relire} : applique la checklist C.O.P.S. et compte tes mots.
\end{enumerate}
\end{methode}

\begin{attention}[Erreurs à éviter]
\begin{itemize}
\item Ne dis pas \eng{I want to make this job} mais \eng{I want to \textbf{do} this job} ou \eng{to \textbf{take up} this profession}.
\item Attention au faux ami : « formation » se traduit par \eng{training} ou \eng{studies}, pas par \eng{formation}.
\item N'oublie pas le \eng{-s} à la troisième personne du singulier : \eng{My uncle \textbf{works} in a garage.}
\item \eng{Actually} signifie « en fait », pas « actuellement » ; dis \eng{nowadays} ou \eng{at present}.
\item \eng{Job} est dénombrable : \eng{a good job}, \eng{many jobs} ; ne dis pas \eng{a work} pour « un travail, un emploi ».
\end{itemize}
\end{attention}

\courssub{Proposition de production et corrigé commenté}

\begin{exoresolu}[Modèle de composition — The trade I dream of]
Rédige une composition de 150 à 200 mots sur le métier de tes rêves, en suivant le plan étudié.

\tcblower
\textbf{Production modèle :}

\eng{The trade I dream of is carpentry. Ever since I was a child, I have admired people who can transform a simple piece of wood into beautiful furniture.}

\eng{First of all, I discovered this trade thanks to my uncle, who owns a small workshop in Bamenda. During the holidays, I used to watch him make tables, chairs and cupboards. Moreover, carpentry attracts me because it is a creative and useful job: a good carpenter builds what families need every day. However, I know that this trade requires patience, precision and hard work.}

\eng{To achieve my dream, I am going to attend a technical college after my A-Levels, where I will learn modern techniques of woodwork. Then, I will open my own workshop in my town.}

\eng{To conclude, carpentry is not just a job for me; it is a passion. I am convinced that with determination and good training, I will succeed and create jobs for other young people in my community.}

(178 mots)

\textbf{Commentaire des choix de langue :}
\begin{itemize}
\item \textbf{Introduction} : le métier est annoncé dès la première phrase, avec le présent perfect \eng{I have admired} pour exprimer un intérêt qui dure depuis l'enfance.
\item \textbf{Connecteurs} : quatre connecteurs variés (\eng{first of all, moreover, however, to conclude}) structurent clairement les idées.
\item \textbf{Temps} : \eng{used to watch} (habitude passée), \eng{am going to attend} (projet décidé), \eng{will learn / will open / will succeed} (prédictions et promesses d'avenir).
\item \textbf{Vocabulaire riche} : \eng{workshop, furniture, precision, to achieve my dream, training}.
\item \textbf{Conclusion} : réaffirmation de l'ambition et ouverture sur l'utilité communautaire, comme le demande l'appel à contributions.
\end{itemize}
\end{exoresolu}

\begin{savaistu}[Le mot carpenter]
Le mot anglais \eng{carpenter} (menuisier) vient du latin \emph{carpentarius}, « artisan du char ». Au Cameroun anglophone, les collèges techniques (\eng{technical colleges}) préparent au \eng{City and Guilds} ou au \eng{GCE Technical}, diplômes très recherchés par les employeurs de Buea, Bamenda ou Douala.
\end{savaistu}

\courssub{Bilan de l'activité}

Dans cette activité d'intégration, tu as mobilisé l'ensemble des savoir-faire de la leçon : comprendre une consigne authentique en anglais, activer le lexique des métiers et de la formation, planifier une composition en trois parties, rédiger en temps limité avec des connecteurs et des temps verbaux variés, puis relire et corriger ta production grâce à la méthode C.O.P.S. L'évaluation croisée t'a en outre appris à lire un texte avec un regard critique et bienveillant : repérer les réussites d'un camarade et formuler un conseil constructif fait aussi partie des compétences d'expression écrite. Ces compétences te serviront directement à l'examen du Probatoire, où la production écrite est notée sur des critères voisins : respect du sujet, organisation, correction grammaticale et richesse du vocabulaire. Garde ta composition corrigée dans ton portfolio : elle pourra servir de base à une version améliorée pour le prochain numéro du magazine scolaire.

\newpage

\coursec{Méthodes et exercices résolus}

\coursec{Lesson 25 — Writing: Compositions about developing and sharing interests in a trade/job/profession}

\courssub{Découverte : un modèle de composition}

\begin{textebox}[Modèle de composition — “The profession I dream of and what I am doing to achieve my dream”]
Have you ever met someone whose dedication changed your view of a profession? For me, that person was Mrs. Awah, the civil engineer who supervised the construction of the bridge over the Menoua River near Dschang. Watching her coordinate workers, read complex plans, and solve problems on the site sparked a fascination that has never left me.

First of all, I discovered this profession during a school trip in Form Four. Before that day, I thought engineering was only about mathematics in a classroom. Mrs. Awah showed me that it is also about improving people’s lives: that bridge allows farmers to reach the market in Bafoussam even during the rainy season. Since then, I have dreamt of becoming a civil engineer myself.

Moreover, this profession attracts me because it combines technical rigour and social impact. I love physics and mathematics, but I also want a career that serves my community. Cameroon needs better roads, resilient buildings, and clean water systems. As a civil engineer, I could contribute directly to the development of our country.

Furthermore, I am already taking concrete steps to reach my goal. At present, I study science subjects seriously, especially mathematics and physics. I am a member of the school’s STEM club where we build model bridges. During the last holidays, I did an internship at a construction firm in Yaoundé where I learnt to read simple architectural drawings.

To conclude, civil engineering is not merely a job for me; it is a vocation. I am determined to succeed in my GCE Advanced Level, enter the National Advanced School of Public Works in Yaoundé, and one day build bridges that connect people. I hope more young Cameroonians will embrace technical professions to build our future.
\end{textebox}

\courssub{Le vocabulaire des métiers et de la formation}

\begin{aretenir}[Vocabulaire thématique : métiers, formation et projet professionnel]
\begin{center}
\begin{tabularx}{\linewidth}{|X|X|X|}
\hline
\rowcolor{popblueL} \textbf{Français} & \textbf{English} & \textbf{Remarque / Contexte} \\
\hline
un métier (manuel) & \eng{a trade} & Ex. : \eng{carpentry, plumbing, welding}. Souvent apprentissage. \\
\hline
un emploi, un poste & \eng{a job} & Terme général. \eng{to apply for a job}, \eng{to get a job}. \\
\hline
une profession (libérale, intellectuelle) & \eng{a profession} & Ex. : \eng{teacher, doctor, lawyer, engineer}. Exige un diplôme. \\
\hline
la formation professionnelle & \eng{vocational training} / \eng{apprenticeship} & \textbf{Attention :} pas de \eng{a formation} (faux ami). \\
\hline
un stage & \eng{an internship} / \eng{a work placement} & \eng{to do an internship}. \\
\hline
un apprentissage & \eng{an apprenticeship} & Contrat travail + formation. \\
\hline
un diplôme & \eng{a diploma} / \eng{a degree} & \eng{Bachelor's degree, Master's degree, PhD}. \\
\hline
une qualification & \eng{a qualification} & \eng{to gain qualifications}. \\
\hline
les compétences & \eng{skills} & \eng{hard skills} (techniques), \eng{soft skills} (relationnelles). \\
\hline
les qualités & \eng{qualities} / \eng{attributes} & \eng{hardworking, reliable, creative, team player}. \\
\hline
la passion, la vocation & \eng{a passion} / \eng{a vocation} & \eng{to have a passion for}, \eng{to feel a vocation for}. \\
\hline
le rêve, l'ambition & \eng{a dream} / \eng{an ambition} & \eng{to pursue one's dream}, \eng{to achieve one's ambition}. \\
\hline
\end{tabularx}
\end{center}

\begin{center}
\begin{tabularx}{\linewidth}{|X|X|X|}
\hline
\rowcolor{poporangeL} \textbf{Verbes d'action} & \textbf{English} & \textbf{Construction / Exemple} \\
\hline
découvrir & \eng{to discover} & \eng{I discovered this trade during a visit.} \\
\hline
s'intéresser à & \eng{to be interested in} & \textbf{Pas} \eng{interested by}. \eng{I am interested in mechanics.} \\
\hline
rêver de & \eng{to dream of} / \eng{to dream about} & \eng{I dream of becoming a pilot.} \\
\hline
choisir & \eng{to choose} (chose, chosen) & Irrégulier. \eng{I chose this path.} \\
\hline
apprendre (qch à qqn) & \eng{to teach} (taught, taught) & \textbf{Faux ami :} \eng{learn} = apprendre (pour soi). \eng{My boss teaches me.} \\
\hline
se former & \eng{to train} / \eng{to train for} & \eng{He is training to be a nurse.} \\
\hline
postuler & \eng{to apply for} & \eng{She applied for the position.} \\
\hline
réussir & \eng{to succeed in} / \eng{to pass} (examen) & \eng{I hope to succeed in my exams.} \\
\hline
atteindre / réaliser & \eng{to achieve} / \eng{to reach} & \eng{to achieve a goal}, \eng{to reach a target}. \\
\hline
partager & \eng{to share} & \eng{I want to share my passion with others.} \\
\hline
\end{tabularx}
\end{center}

\end{aretenir}

\begin{attention}[Faux amis fréquents — Monde du travail]
\begin{enumerate}
\item \eng{Actually} = « en fait, en réalité » (pas « actuellement » → \eng{currently, at present}).
\item \eng{To assist} = « aider, assister qqn » (pas « assister à» → \eng{to attend a meeting/course}).
\item \eng{A formation} n'existe pas pour « une formation » → \eng{training, a course, a programme, vocational training}.
\item \eng{To learn} = « apprendre » (soi-même) ; « apprendre à quelqu'un » = \eng{to teach somebody}.
\item \eng{Salary} = « salaire » (mensuel) ; \eng{wage} = « salaire » (horaire/hebdo, souvent ouvrier) ; \eng{pay} = « paie, rémunération » (général).
\item \eng{Career} = « carrière » (parcours pro) ; \eng{carrier} = « transporteur, porteur ».
\end{enumerate}
\end{attention}

\courssub{Connecteurs logiques et temps verbaux}

\begin{methode}[Les connecteurs logiques — Donner du rythme à sa composition]
Les \eng{linking words} structurent la pensée. Choisis-les selon la relation logique entre les idées.

\begin{center}
\begin{tabularx}{\linewidth}{|l|X|X|}
\hline
\rowcolor{popblueL} \textbf{Relation} & \textbf{Connecteurs (variés)} & \textbf{Exemple en contexte} \\
\hline
Addition & \eng{moreover, furthermore, also, in addition, besides} & \eng{Moreover, this trade offers stability.} \\
\hline
Cause & \eng{because, since, as} & \eng{Since I love children, I chose teaching.} \\
\hline
Conséquence & \eng{therefore, consequently, as a result, that is why} & \eng{The training is hard; therefore, I work daily.} \\
\hline
Opposition / Concession & \eng{however, nevertheless, although, even though, whereas, but} & \eng{Although the pay is low, I love this job.} \\
\hline
Ordre chronologique / Énumération & \eng{first of all, to begin with, then, next, after that, finally} & \eng{First of all, I need my A Levels.} \\
\hline
Illustration / Exemple & \eng{for example, for instance, such as} & \eng{Trades such as welding are in demand.} \\
\hline
Conclusion & \eng{to conclude, in conclusion, all in all, ultimately} & \eng{To conclude, I am ready for the challenge.} \\
\hline
\end{tabularx}
\end{center}

\end{methode}

\begin{propriete}[Règles d'or des connecteurs]
\begin{enumerate}
\item \textbf{Virgule obligatoire} après un connecteur en début de phrase : \eng{However, I disagree.}
\item \textbf{Pas de virgule} devant \eng{because, since, as} au milieu de phrase : \eng{I stay because I like it.}
\item \textbf{Varier} : ne jamais démarrer deux paragraphes par le même connecteur.
\item \textbf{But} : \eng{in order to} + BV / \eng{so as to} + BV (formel) : \eng{I study hard in order to succeed.}
\end{enumerate}
\end{propriete}

\begin{exoresolu}[Compléter un paragraphe avec les connecteurs]
\textbf{Énoncé.} Complete the paragraph with the following linking words: \emph{however — first of all — because — moreover — to conclude}.

\eng{\ldots\ldots\ldots, I must say that becoming a mechanic was not my first choice. I wanted to be a pilot \ldots\ldots\ldots I loved aeroplanes. \ldots\ldots\ldots, when my uncle opened his garage in Douala, I discovered the beauty of engines. \ldots\ldots\ldots, I realised that mechanics earn a good living here. \ldots\ldots\ldots, I am now proud of my choice and I work hard every day.}

\tcblower
\textbf{Solution détaillée.}

\begin{enumerate}
\item \eng{\textbf{First of all}, I must say that becoming a mechanic was not my first choice.} — Ouverture du paragraphe : connecteur d'ordre initial, suivi d'une virgule.
\item \eng{I wanted to be a pilot \textbf{because} I loved aeroplanes.} — Cause. Pas de virgule devant \eng{because} en milieu de phrase.
\item \eng{\textbf{However}, when my uncle opened his garage in Douala, I discovered the beauty of engines.} — Opposition (premier choix vs découverte). Connecteur de contraste + virgule.
\item \eng{\textbf{Moreover}, I realised that mechanics earn a good living here.} — Argument additionnel. Connecteur d'addition + virgule.
\item \eng{\textbf{To conclude}, I am now proud of my choice and I work hard every day.} — Conclusion du paragraphe. Connecteur de conclusion + virgule.
\end{enumerate}

\textbf{Conclusion.} Vérifie toujours la \emph{logique} avant la forme : la phrase ajoute-t-elle, oppose-t-elle, explique-t-elle ou conclut-elle ?
\end{exoresolu}

\begin{methode}[Les temps verbaux — La frise temporelle de la composition]
Une composition sur un projet professionnel navigue entre passé, présent et futur. Identifie le marqueur temporel pour choisir le temps.

\begin{center}
\begin{tabular}{|l|l|l|l|}
\hline
\rowcolor{popgreenL} \textbf{Fonction} & \textbf{Marqueurs typiques} & \textbf{Temps} & \textbf{Exemple} \\
\hline
Événement passé, daté, terminé & \eng{yesterday, last year, two years ago, in 2020, when I was...} & \textbf{Prétérit simple} & \eng{I \textbf{met} a carpenter in Buea three years ago.} \\
\hline
Expérience vécue (passé $\rightarrow$ présent) & \eng{since 2020, since then, for three years, recently, ever} & \textbf{Present Perfect} & \eng{I \textbf{have learnt} a lot \textbf{since then}.} \\
\hline
Goûts, vérités, situation actuelle & \eng{now, at present, currently, usually, always} & \textbf{Présent simple} & \eng{I \textbf{love} repairing engines.} (verbes d'état : \eng{love, want, know, prefer} $\neq$ \emph{-ing}) \\
\hline
Action en cours maintenant & \eng{at the moment, currently, this year} & \textbf{Present Continuous} & \eng{I \textbf{am studying} hard \textbf{this term}.} \\
\hline
Projet décidé, intention & \eng{soon, next year, after my exams, tomorrow} & \textbf{Be going to} + BV & \eng{I \textbf{am going to} apply to ENSTP.} \\
\hline
Prédiction, promesse, décision sur l'instant & \eng{probably, certainly, I hope, I promise} & \textbf{Will} + BV & \eng{I \textbf{will} never give up.} \\
\hline
Hypothèse (conditionnel) & \eng{if + prétérit} & \textbf{Would} + BV & \eng{If I \textbf{had} money, I \textbf{would buy} tools.} \\
\hline
\end{tabular}
\end{center}

\end{methode}

\begin{propriete}[Verbes irréguliers clés — À connaître par cœur]
\begin{center}
\begin{tabular}{|l|l|l|l|}
\hline
\rowcolor{poppurpleL} \textbf{Base} & \textbf{Prétérit} & \textbf{Participe passé} & \textbf{Traduction} \\
\hline
\eng{begin} & \eng{began} & \eng{begun} & commencer \\
\hline
\eng{choose} & \eng{chose} & \eng{chosen} & choisir \\
\hline
\eng{do} & \eng{did} & \eng{done} & faire \\
\hline
\eng{learn} & \eng{learnt/learned} & \eng{learnt/learned} & apprendre \\
\hline
\eng{meet} & \eng{met} & \eng{met} & rencontrer \\
\hline
\eng{teach} & \eng{taught} & \eng{taught} & enseigner, apprendre à \\
\hline
\eng{think} & \eng{thought} & \eng{thought} & penser, croire \\
\hline
\eng{become} & \eng{became} & \eng{become} & devenir \\
\hline
\eng{get} & \eng{got} & \eng{got/gotten} & obtenir, devenir \\
\hline
\end{tabular}
\end{center}
\end{propriete}

\begin{exoresolu}[Mettre les verbes au bon temps]
\textbf{Énoncé.} Put the verbs in brackets into the correct tense.

\eng{Two years ago, my father (take) \ldots\ldots\ldots me to his cocoa farm in the Centre region. I (discover) \ldots\ldots\ldots how chocolate begins its journey. Since that day, I (always / want) \ldots\ldots\ldots to become an agricultural engineer. At present, I (study) \ldots\ldots\ldots very hard in science subjects. After my baccalauréat, I (go) \ldots\ldots\ldots to university in Yaoundé.}

\tcblower
\textbf{Solution détaillée.}

\begin{enumerate}
\item \eng{Two years ago, my father \textbf{took} me to his cocoa farm.} — \eng{Two years ago} = repère de passé terminé $\rightarrow$ prétérit simple. Verbe irrégulier : \eng{take/took/taken}.
\item \eng{I \textbf{discovered} how chocolate begins its journey.} — Même contexte passé $\rightarrow$ prétérit. Verbe régulier : \eng{discover} + \emph{-ed}.
\item \eng{Since that day, I \textbf{have always wanted} to become an agricultural engineer.} — \eng{Since that day} impose le present perfect : \eng{have} + participe passé \eng{wanted}. \textbf{Piège :} on ne dit pas \eng{I want since that day} (calque du français).
\item \eng{At present, I \textbf{study} very hard in science subjects.} — Habitude actuelle $\rightarrow$ présent simple. (Le present continuous \eng{am studying} est accepté pour insister sur la période en cours.)
\item \eng{After my baccalauréat, I \textbf{will go} / \textbf{am going to go} to university in Yaoundé.} — Projet d'avenir : \eng{going to} (projet décidé) est le plus naturel ; \eng{will} (prédiction) est accepté.
\end{enumerate}

\textbf{Conclusion.} La cohérence des temps est un critère majeur de la grille d'évaluation : un seul marqueur temporel suffit à déterminer le temps du verbe.
\end{exoresolu}

\courssub{La structure de la composition et le plan de rédaction}

\begin{methode}[Planifier sa composition avant d'écrire — 5 étapes]
Une composition réussie se prépare \textbf{avant} d'écrire la première phrase.

\begin{enumerate}
\item \textbf{Analyser le sujet} : souligne les mots-clés (\eng{trade}, \eng{job}, \eng{profession}, \eng{interest}, \eng{share}, \eng{develop}) et identifie la double tâche : \emph{décrire} (le métier, l'intérêt) + \emph{expliquer/justifier} (actions, raisons, partage).
\item \textbf{Faire un remue-méninges} (\eng{brainstorming}) : note en 3 minutes toutes les idées et le vocabulaire en anglais, sans trier.
\item \textbf{Choisir 3 idées principales} : une idée = un paragraphe de développement. Ex. : (1) Découverte du métier, (2) Raisons de l'intérêt / adéquation avec valeurs, (3) Actions concrètes pour y parvenir / partage.
\item \textbf{Rédiger un plan détaillé} :
\begin{itemize}
\item Introduction (accroche + sujet + annonce du plan).
\item Développement (2 ou 3 paragraphes, 1 idée principale chacun, connecteur en tête).
\item Conclusion (bilan + ouverture).
\end{itemize}
\item \textbf{Prévoir les connecteurs} : associe à chaque paragraphe son connecteur de départ (\eng{First of all}, \eng{Moreover}, \eng{Finally}).
\end{enumerate}

\textbf{Réflexe d'examen} : consacre 10 minutes au plan sur 2 heures. Une composition bien planifiée s'écrit deux fois plus vite et sans ratures.
\end{methode}

\begin{exoresolu}[Analyser un sujet et construire un plan]
\textbf{Énoncé.} Your English teacher has asked you to write a composition (about 200 words) entitled: \emph{“The profession I dream of and what I am doing to achieve my dream.”}

\begin{enumerate}
\item Identify the two tasks contained in the topic.
\item Propose a three-paragraph plan for the body of your composition.
\item Suggest one sentence to open your introduction.
\end{enumerate}

\tcblower
\textbf{Solution détaillée.}

\textbf{1. Les deux tâches.} Le sujet contient deux verbes d'action : \eng{dream of} (le métier dont je rêve $\rightarrow$ décrire et justifier) et \eng{am doing} (ce que je fais concrètement $\rightarrow$ actions présentes). Oublier l'une des deux parties fait perdre la moitié des points de contenu.

\textbf{2. Plan du développement.}
\begin{itemize}
\item Paragraph 1: How I discovered the profession (a nurse who helped my grandmother in Bamenda).
\item Paragraph 2: Why nursing attracts me (helping people, job opportunities, respect in the community).
\item Paragraph 3: What I am doing now (working hard in biology, volunteering at the health centre, learning first aid).
\end{itemize}

\textbf{3. Phrase d'ouverture possible.}

\eng{Ever since I was a child, I have always admired people who wear a white coat and take care of the sick.}

« Depuis mon enfance, j'ai toujours admiré les personnes qui portent une blouse blanche et s'occupent des malades. »

\textbf{Conclusion méthodologique.} Remarque le temps utilisé : \eng{have always admired} (present perfect) car l'admiration a commencé dans le passé et continue aujourd'hui. Une bonne phrase d'accroche est personnelle, courte et annonce le thème sans le copier mot à mot.
\end{exoresolu}

\begin{methode}[Construire une introduction et une conclusion efficaces]
\textbf{L'introduction} (3 à 4 phrases) suit le schéma : \textbf{Accroche} $\rightarrow$ \textbf{Sujet} $\rightarrow$ \textbf{Annonce du plan}.

\begin{enumerate}
\item \textbf{Accroche} : une question, un souvenir personnel ou une vérité générale.
\eng{Have you ever watched a carpenter transform a simple piece of wood into a beautiful table?}
\item \textbf{Sujet} : phrase de transition qui personnalise le thème.
\eng{For me, this is more than a job: it is a passion I wish to share.}
\item \textbf{Annonce du plan} : utilise \eng{will} + verbes d'action + marqueurs de séquence.
\eng{I will first explain how I discovered this trade, then I will show why it matters to our community.}
\end{enumerate}

\textbf{La conclusion} (2 à 3 phrases) : \textbf{Bilan} + \textbf{Ouverture}, sans idée nouvelle.
\begin{enumerate}
\item \textbf{Bilan} : résume l'idée force.
\eng{To conclude, carpentry is not only my dream job; it is also a way of serving my country.}
\item \textbf{Ouverture} : perspective future, vœu, appel.
\eng{I hope that more young Cameroonians will dare to choose technical trades.}
\end{enumerate}

\end{methode}

\begin{aretenir}[À retenir pour l'intro et la conclusion]
\begin{itemize}
\item Ne recopie \textbf{jamais} le sujet comme première phrase.
\item N'écris \textbf{jamais} « the end » à la fin.
\item La conclusion ne doit \textbf{pas} introduire d'argument nouveau.
\item L'annonce du plan utilise \eng{I will} (ou \eng{This composition will}) + \eng{first / then / finally}.
\end{itemize}
\end{aretenir}

\begin{exoresolu}[Rédiger une introduction complète]
\textbf{Énoncé.} Write the introduction (3–4 sentences) of a composition entitled: \emph{“Why I want to become a teacher of English.”} Use a question as a hook and announce your plan.

\tcblower
\textbf{Solution détaillée.}

\textbf{Réflexion.} Il faut trois éléments : (1) une question d'accroche, (2) une phrase qui pose le sujet, (3) une annonce du plan en deux temps. On utilise le présent simple pour les vérités générales et \eng{will} pour annoncer le plan.

\textbf{Introduction rédigée.}

\eng{Have you ever met a person who changed your life with simple words? For me, that person was my English teacher in Form Three. Since then, I have dreamt of becoming a teacher of English myself. In this composition, I will first explain why this profession attracts me, and then I will describe the steps I am taking to reach my goal.}

« Avez-vous déjà rencontré une personne qui a changé votre vie avec de simples mots ? Pour moi, cette personne était mon professeur d'anglais en classe de troisième. Depuis, je rêve de devenir moi-même professeur d'anglais. Dans cette composition, j'expliquerai d'abord pourquoi ce métier m'attire, puis je décrirai les étapes que je franchis pour atteindre mon but. »

\textbf{Justifications grammaticales.} \eng{Have you ever met} : present perfect pour une expérience de vie indatée. \eng{have dreamt} : present perfect avec \eng{since then}. \eng{I will first explain\ldots{} and then I will describe} : futur pour l'annonce du plan, avec les marqueurs \eng{first} / \eng{then}.

\textbf{Conclusion.} Une introduction de ce type prend moins de 5 minutes à rédiger quand le plan est prêt, et elle donne immédiatement au correcteur une impression de maîtrise.
\end{exoresolu}

\courssub{Modèle commenté et erreurs à éviter}

\begin{exemplebox}[Analyse commentée du modèle (texte de la Découverte)]
\textbf{Structure :}
\begin{itemize}
\item \textbf{Intro (3 phrases) :} Question (\eng{Have you ever met...}) $\rightarrow$ Souvenir précis (\eng{Mrs. Awah... Menoua River}) $\rightarrow$ Annonce plan (\eng{First of all... Moreover... Furthermore... To conclude}).
\item \textbf{§1 (Discovery) :} \eng{First of all} + Prétérit (\eng{discovered, thought, showed, allows}). Contexte daté (\eng{during a school trip in Form Four}).
\item \textbf{§2 (Reasons) :} \eng{Moreover} + Présent simple (\eng{attracts, combines, love, want, needs}). Verbes d'état.
\item \textbf{§3 (Actions) :} \eng{Furthermore} + Present Continuous (\eng{am taking, am studying, am a member}) + Prétérit (\eng{did, learnt}) pour stage passé.
\item \textbf{Conclusion (2 phrases) :} \eng{To conclude} + Présent (\eng{is, am determined}) + \eng{will} (\eng{will embrace, will build}). Ouverture sur la jeunesse camerounaise.
\end{itemize}

\textbf{Points de langue remarquables :}
\begin{itemize}
\item \eng{sparked a fascination that has never left me} : prétérit + relative au present perfect (lien passé-présent).
\item \eng{allows farmers to reach} : structure \eng{allow + objet + to + BV}.
\item \eng{dreamt of becoming} : \eng{dream of + V-ing}.
\item \eng{contribute directly to} : \eng{contribute to + nom/V-ing}.
\end{itemize}
\end{exemplebox}

\begin{attention}[Les 5 erreurs qui coûtent des points au Bac]
\begin{enumerate}
\item \textbf{Faux amis du monde du travail.}
\eng{Actually} = « en fait » (pas « actuellement » $\rightarrow$ \eng{currently, at present}).
\eng{To assist} = « aider » (« assister à un cours » = \eng{to attend a course}).
\eng{A formation} n'existe pas : dis \eng{training} ou \eng{vocational training}.
\item \textbf{Present perfect avec un repère de passé daté.}
\eng{I have visited} \textbf{✗} avec \eng{last year / two years ago / in 2020} $\rightarrow$ prétérit obligatoire : \eng{I visited a workshop last year.}
\item \textbf{Prépositions calquées du français.}
\eng{interested \textbf{in}} (pas \emph{by}), \eng{to dream \textbf{of}}, \eng{to depend \textbf{on}}, \eng{to be good \textbf{at}}, \eng{to look \textbf{for} a job} (chercher un emploi), \eng{to work \textbf{as} a nurse} (travailler comme).
\item \textbf{Accord sujet-verbe à la 3\textsuperscript{e} personne du singulier.}
\eng{My brother work} \textbf{✗} $\rightarrow$ \eng{My brother work\textbf{s} as a mechanic.} Le \emph{s} du présent simple est la faute la plus fréquente et la plus sanctionnée.
\item \textbf{Ordre des mots et orthographe.}
Pas d'inversion du sujet et du verbe après un connecteur : \eng{Therefore, I decided} (jamais \eng{Therefore decided I}).
Orthographe à revoir : \eng{profession} (un \emph{f}), \eng{business}, \eng{successful} (un \emph{l} final, deux \emph{c}), \eng{career} (pas \emph{carreer}).
\end{enumerate}
\end{attention}

\courssub{Grille d'évaluation et production guidée}

\begin{methode}[La relecture en 4 passes — Réserve 10 minutes]
Effectue \textbf{4 passes séparées} pour ne rien rater.

\begin{enumerate}
\item \textbf{Pass 1 — Contenu et structure} : Ai-je répondu à \textbf{tout} le sujet ? Y a-t-il une intro, 2–3 paragraphes de développement, une conclusion ? Chaque paragraphe a-t-il une idée principale claire ?
\item \textbf{Pass 2 — Grammaire} : Accord sujet-verbe (\eng{he work\textbf{s}}), cohérence des temps verbaux, articles (\eng{a/an/the}), pluriels (\eng{two job\textbf{s}}), formation du passé (irréguliers).
\item \textbf{Pass 3 — Vocabulaire et orthographe} : Faux amis (\eng{actually}, \eng{assist}, \eng{training}), mots de métier corrects (\eng{a trade}, \eng{a profession}), orthographe (\eng{profession}, \eng{business}, \eng{successful}, \eng{career}).
\item \textbf{Pass 4 — Présentation et cohésion} : Connecteurs variés et bien ponctués (virgule), majuscules en début de phrase, longueur respectée (environ 200 mots $\pm$ 10\%).
\end{enumerate}

\textbf{Grille d'auto-évaluation rapide (sur 20) :}
\begin{center}
\begin{tabular}{|l|c|}
\hline
\rowcolor{popgoldL} \textbf{Critère} & \textbf{Points} \\
\hline
Contenu et respect du sujet & /5 \\
\hline
Organisation et connecteurs & /5 \\
\hline
Grammaire et temps verbaux & /5 \\
\hline
Vocabulaire et orthographe & /5 \\
\hline
\textbf{Total} & \textbf{/20} \\
\hline
\end{tabular}
\end{center}
\end{methode}

\begin{exoresolu}[Corriger un extrait d'élève — Entraînement à la relecture]
\textbf{Énoncé.} The following paragraph contains \textbf{six mistakes}. Find them and rewrite the paragraph correctly.

\eng{Last year, I have choosed to become a electrician. Actualy, I am very interesting by this trade because I like repair things. My father, who is a electrician too, learn me a lot.}

\tcblower
\textbf{Solution détaillée.}

\begin{enumerate}
\item \eng{I have choosed} $\rightarrow$ \eng{I \textbf{chose}}. Double erreur : avec \eng{last year} (passé daté), on utilise le prétérit, pas le present perfect ; et le participe passé de \eng{choose} est \eng{chosen} (jamais \eng{choosed}) : \eng{choose/chose/chosen}.
\item \eng{a electrician} $\rightarrow$ \eng{\textbf{an} electrician}. Voyelle initiale (\emph{e}) $\rightarrow$ article \eng{an}.
\item \eng{Actualy} $\rightarrow$ \eng{\textbf{Actually}}. Orthographe : deux \emph{l}. Attention au sens : \eng{actually} signifie « en fait » ; pour « actuellement », dis \eng{at present} ou \eng{currently}.
\item \eng{I am very interesting by} $\rightarrow$ \eng{I am very \textbf{interested in}}. Double correction : \eng{interested} (celui qui ressent l'intérêt) et non \eng{interesting} (ce qui provoque l'intérêt) ; préposition \eng{in}, pas \eng{by} (calque du français « intéressé par »).
\item \eng{I like repair things} $\rightarrow$ \eng{I like \textbf{repairing} things} (ou \eng{to repair}). Après \eng{like}, gérondif en \emph{-ing} ou infinitif, jamais la base verbale seule.
\item \eng{learn me a lot} $\rightarrow$ \eng{\textbf{teaches} me a lot}. Faux ami : \eng{to learn} = « apprendre » ; « apprendre à quelqu'un » = \eng{to teach somebody}. Et accord 3\textsuperscript{e} personne : \eng{he teaches}.
\end{enumerate}

\textbf{Paragraphe corrigé.}

\eng{Last year, I chose to become an electrician. Actually, I am very interested in this trade because I like repairing things. My father, who is an electrician too, teaches me a lot.}

\textbf{Conclusion.} Six fautes de ce type dans une composition font chuter la note de langue. La relecture ciblée (passes 2 et 3) permet d'en éliminer la quasi-totalité.
\end{exoresolu}

\begin{exercice}[Production guidée — À vous de jouer !]
\textbf{Consigne.} Write a composition of about \textbf{200 words} on the following topic:

\emph{“The trade I have chosen and why I want to share this passion with others.”}

\textbf{Étapes de travail (chronométrées sur 2 h) :}
\begin{enumerate}
\item \textbf{10 min — Planification :} Analyse le sujet (double tâche : \eng{chosen} + \eng{share}). Fais ton brainstorming. Choisis 3 idées pour le développement. Rédige ton plan avec connecteurs.
\item \textbf{15 min — Rédaction de l'introduction :} Accroche (question/souvenir), Sujet, Annonce du plan.
\item \textbf{50 min — Rédaction du développement :} 3 paragraphes. 1 idée = 1 paragraphe. Connecteur en tête. Vocabulaire précis (métier, formation, qualités). Temps verbaux cohérents.
\item \textbf{15 min — Rédaction de la conclusion :} Bilan + Ouverture. Pas d'idée nouvelle.
\item \textbf{20 min — Relecture (4 passes) :} Contenu $\rightarrow$ Grammaire $\rightarrow$ Vocab/Ortho $\rightarrow$ Présentation. Compte les mots.
\item \textbf{10 min — Copie propre :} Recopie lisiblement sur la copie d'examen.
\end{enumerate}

\textbf{Pistes pour le brainstorming (exemples) :}
\begin{itemize}
\item \textbf{Métiers manuels :} \eng{carpentry, plumbing, welding, hairdressing, tailoring, mechanics}.
\item \textbf{Métiers de service :} \eng{nursing, teaching, catering, tourism guide, IT technician}.
\item \textbf{Partage :} \eng{mentoring young apprentices, teaching in a vocational centre, opening a workshop to train others, promoting the trade on social media, encouraging girls in technical trades}.
\end{itemize}
\end{exercice}


\begin{corrige}[Exercice Production guidée — Grille de correction attendue]
\textbf{Contenu (5 pts) :}
\begin{itemize}
\item Le métier est clairement identifié et justifié (découverte, passion) : 2 pts.
\item Les actions concrètes pour le partager (formation d'autres, promotion, mentorat) sont décrites : 2 pts.
\item Le ton est personnel, engagé et adapté au contexte camerounais : 1 pt.
\end{itemize}

\textbf{Organisation (5 pts) :}
\begin{itemize}
\item Introduction complète (accroche, sujet, annonce de plan) : 1 pt.
\item 3 paragraphes de développement distincts, 1 idée principale chacun : 2 pts.
\item Connecteurs variés et correctement ponctués (min. 4 différents) : 1 pt.
\item Conclusion (bilan + ouverture) sans idée nouvelle : 1 pt.
\end{itemize}

\textbf{Langue — Grammaire (5 pts) :}
\begin{itemize}
\item Temps verbaux cohérents avec marqueurs (prétérit, present perfect, présent, futur) : 2 pts.
\item Accords sujet-verbe (surtout 3\textsuperscript{e} pers. sg.) et articles corrects : 1,5 pt.
\item Structures complexes tentées (relatives, gérondifs, \eng{in order to}, conditionnel) : 1,5 pt.
\end{itemize}

\textbf{Langue — Vocabulaire / Orthographe (5 pts) :}
\begin{itemize}
\item Vocabulaire précis du thème (métiers, formation, qualités) : 2 pts.
\item Pas de faux amis (\eng{actually, assist, training}), prépositions correctes (\eng{interested in, good at, work as}) : 2 pts.
\item Orthographe soignée (mots-clés : \eng{profession, successful, career, business}) : 1 pt.
\end{itemize}

\textbf{Note finale :} Somme des 4 critères / 20. \textbf{Seuil de réussite : 10/20.}
\end{corrige}

\newpage

\coursec{Exercices}

\courssub{Je vérifie mes connaissances}

\begin{exercice}[QCM : le bon temps verbal]
Pour chaque phrase, choisis la forme verbale correcte (A, B ou C).

\begin{enumerate}
\item I \ldots\ to become a nurse since I was twelve.\\
A. want \quad B. have wanted \quad C. wanted
\item Last year, my cousin \ldots\ an apprenticeship in a garage in Douala.\\
A. starts \quad B. has started \quad C. started
\item After my A Levels, I \ldots\ a computer science course.\\
A. am going to follow \quad B. follow \quad C. followed
\item A good tailor \ldots\ patience and precision.\\
A. require \quad B. requires \quad C. is requiring
\item My father \ldots\ as a mechanic for twenty years now.\\
A. works \quad B. has worked \quad C. worked
\end{enumerate}
\end{exercice}

\begin{exercice}[Vrai ou faux ?]
Indique si chaque affirmation est vraie (V) ou fausse (F) et justifie ta réponse en une phrase.

\begin{enumerate}
\item The introduction of a composition presents the topic and announces the plan.
\item \eng{However} is used to add a similar idea.
\item \eng{In conclusion} is normally placed at the beginning of the first paragraph.
\item The present simple is used to describe habits and general truths about a job.
\item A composition about a trade should mix ideas without paragraphs to save time.
\end{enumerate}
\end{exercice}

\begin{exercice}[Vocabulaire : métiers et formation]
Complète chaque phrase avec le mot anglais qui convient, choisi dans la liste : \emph{apprenticeship, skills, workshop, qualification, trade, career}.

\begin{enumerate}
\item My brother is learning carpentry in a small \ldots\ near the Mokolo market.
\item Nursing is a noble \ldots\ that demands love for others.
\item She obtained her \ldots\ as an electrician after two years of training.
\item During his \ldots\ , the young plumber worked with an experienced master.
\item Good communication \ldots\ are necessary in every profession.
\item He wants to make a \ldots\ in journalism.
\end{enumerate}
\end{exercice}

\begin{exercice}[Conjugaison à compléter]
Mets le verbe entre parenthèses au temps qui convient (present simple, past simple, present perfect ou be going to).

\begin{enumerate}
\item When I was a child, I often \ldots\ (watch) my uncle repair radios.
\item Since then, I \ldots\ (develop) a real passion for electronics.
\item Every Saturday, I \ldots\ (help) him in his shop in Bamenda.
\item Next year, I \ldots\ (enter) a technical college.
\item So far, this experience \ldots\ (teach) me discipline and hard work.
\end{enumerate}
\end{exercice}

\courssub{Je m'entraîne}

\begin{exercice}[Traduction]
Traduis les phrases suivantes en anglais, en utilisant le vocabulaire et les connecteurs étudiés dans la leçon.

\begin{enumerate}
\item Depuis mon enfance, je m'intéresse au métier d'infirmière.
\item D'abord, ce métier demande de la patience ; ensuite, il exige une bonne formation.
\item Mon frère a commencé un apprentissage chez un menuisier il y a deux ans.
\item Cependant, beaucoup de jeunes préfèrent les métiers de l'informatique.
\item En conclusion, je suis déterminé à partager ma passion avec mes camarades.
\end{enumerate}
\end{exercice}

\begin{exercice}[Relier les phrases avec un connecteur]
Relie chaque paire de phrases en une seule phrase en utilisant le connecteur indiqué entre parenthèses : \emph{however, therefore, moreover, for example, in conclusion}.

\begin{enumerate}
\item Training as a mechanic is long. It is very rewarding. (however)
\item Many trades require practical abilities. A welder must master his tools perfectly. (for example)
\item I love helping sick people. I have decided to become a nurse. (therefore)
\item Carpentry develops creativity. It offers good job opportunities in our towns. (moreover)
\item I have explained my project and my motivations. I am ready to start my training. (in conclusion)
\end{enumerate}
\end{exercice}

\begin{exercice}[Composition à trous]
Complète le texte suivant avec les mots de la liste : \emph{apprenticeship, skills, workshop, qualification, trade, interested, training, share}.

\begin{textebox}[My interest in carpentry]
Since my childhood, I have been \ldots\ in carpentry. My uncle owns a small \ldots\ in Bafoussam, and I often spent my holidays there watching him work. Two years ago, I started an \ldots\ with him. During this \ldots\ , I learnt many useful \ldots\ , such as measuring wood and using machines safely. Carpentry is a \ldots\ that I respect deeply. After my \ldots\ , I hope to open my own business and \ldots\ my passion with young people of my quarter.
\end{textebox}
\end{exercice}

\begin{exercice}[Corrige le texte d'un élève]
Le texte suivant contient \textbf{8 erreurs} (temps verbaux, 3\textsuperscript{e} personne, calques du français, connecteurs). Souligne chaque erreur et écris la correction.

\begin{textebox}[A student's draft]
I am interesting in the job of teacher since five years. My mother teach in a primary school in Yaoundé and I admire she very much. Firstable, teachers form the future of the nation. Morever, they share their knowledges with children. Last year, I have helped my little brother with his homework every evening. I am going to pass my A Level next year, and after I will enter in a teacher training college. In conclusion, the teaching is my dream job.
\end{textebox}
\end{exercice}

\courssub{Je me prépare au Bac}

\begin{exercice}[Compréhension et langue — type Bac]
Lis le texte suivant, puis réponds aux questions.

\begin{textebox}[A passion for tailoring]
Amina is a seventeen-year-old student at a high school in Douala. Since she was ten, she has been fascinated by tailoring. Her grandmother, who owned a small sewing shop in Deido, taught her how to use a needle and thread. At first, Amina only repaired her own clothes; gradually, she started designing dresses for her sisters and neighbours.

Last year, during the long holidays, Amina followed a three-month training course in a fashion workshop. There, she learnt how to take measurements, cut fabric and use an electric sewing machine. Moreover, she discovered that tailoring is not only a practical skill but also an art.

However, some of her classmates laughed at her dream. They believed that only office jobs were respectable. Amina did not listen to them. Therefore, she worked harder and even won a prize at a local fashion show.

After her A Levels, Amina is going to attend a professional school of fashion design. In conclusion, she wants to open her own workshop one day and train other young girls, so that they too can develop and share their interest in this beautiful trade.
\end{textebox}

\textbf{A. Comprehension questions} — Answer in complete sentences.
\begin{enumerate}
\item Who taught Amina the basics of tailoring?
\item Mention two things she learnt during her training course.
\item Why did some classmates laugh at her dream?
\item What are Amina's plans for the future?
\end{enumerate}

\textbf{B. Language work}
\begin{enumerate}
\item Find in the text two connectors of addition and one connector of opposition.
\item Identify the tense used in “she has been fascinated” and justify its use.
\item Transform into the passive voice: “Her grandmother taught her how to use a needle.”
\end{enumerate}
\end{exercice}

\begin{exercice}[Remettre la composition dans l'ordre — type Bac]
Les cinq paragraphes (A–E) d'une composition intitulée \emph{“My interest in the nursing profession”} ont été mélangés. Remets-les dans l'ordre logique (introduction, développement, conclusion) et justifie ton ordre en relevant les connecteurs et les indices de sens.

\begin{enumerate}
\item[\textbf{A.}] In conclusion, nursing is the profession of my heart, and I am determined to share this interest with other students of my school.
\item[\textbf{B.}] Moreover, nurses are highly respected in our communities because they comfort families in difficult moments.
\item[\textbf{C.}] Many young people dream of becoming doctors or engineers, but my greatest wish is to become a nurse. I will explain the origin of this interest and my plans for the future.
\item[\textbf{D.}] First of all, I discovered this profession when my grandmother fell ill and I spent two weeks with her in a hospital in Buea. The kindness of the nurses deeply touched me.
\item[\textbf{E.}] Therefore, after my A Levels, I am going to apply to a nursing school in order to obtain my qualification.
\end{enumerate}
\end{exercice}

\begin{exercice}[Sujet de composition — type Bac]
\textbf{Topic:} Write a composition of \textbf{150–200 words} on the following subject:

\begin{center}
\emph{“Developing and sharing my interest in a trade, job or profession.”}
\end{center}

\textbf{Instructions:}
\begin{enumerate}
\item Your composition must contain an \textbf{introduction} (presentation of the trade and origin of your interest), a \textbf{development} (two or three paragraphs: qualities required, training, advantages) and a \textbf{conclusion} (your plans and how you will share your interest).
\item Use at least \textbf{four different link words} (e.g. \emph{first of all, moreover, however, therefore, in conclusion}).
\item Use the correct tenses: past simple for the origin of your interest, present simple for general truths, \emph{be going to} for your future plans.
\item Write neatly and respect the number of words. Indicate the word count at the end.
\end{enumerate}
\end{exercice}

\newpage

\coursec{Corrigés des exercices}

\begin{corrige}[Exercice 1 — QCM : le bon temps verbal]
\begin{enumerate}
\item \textbf{B. have wanted} — \eng{I have wanted to become a nurse since I was twelve.} La présence de \eng{since I was twelve} impose le \cle{present perfect} (action commencée dans le passé et qui continue). « Depuis mes douze ans, je veux devenir infirmière. »
\item \textbf{C. started} — \eng{Last year, my cousin started an apprenticeship in a garage in Douala.} Le repère \eng{last year} (moment passé terminé) impose le \cle{past simple}.
\item \textbf{A. am going to follow} — \eng{After my A Levels, I am going to follow a computer science course.} Il s'agit d'un \cle{projet d'avenir} déjà décidé : \eng{be going to + base verbale}.
\item \textbf{B. requires} — \eng{A good tailor requires patience and precision.} Vérité générale au \cle{present simple} ; attention à la 3\textsuperscript{e} personne du singulier : \eng{require\textbf{s}}.
\item \textbf{B. has worked} — \eng{My father has worked as a mechanic for twenty years now.} \eng{For twenty years now} indique une durée qui continue : \cle{present perfect}.
\end{enumerate}
\end{corrige}

\begin{corrige}[Exercice 2 — Vrai ou faux ?]
\begin{enumerate}
\item \textbf{V} — \eng{The introduction presents the topic and announces the plan.} C'est la définition même de l'introduction : présenter le sujet et annoncer le plan.
\item \textbf{F} — \eng{However} exprime l'\cle{opposition}, pas l'addition. Pour ajouter une idée semblable, on utilise \eng{moreover}, \eng{furthermore} ou \eng{also}.
\item \textbf{F} — \eng{In conclusion} se place au début du \cle{dernier} paragraphe (la conclusion), jamais au début du premier.
\item \textbf{V} — Le \cle{present simple} décrit les habitudes et les vérités générales : \eng{A nurse works in a hospital.}
\item \textbf{F} — Une composition doit être organisée en \cle{paragraphes} (introduction, développement, conclusion) ; mélanger les idées rend le texte incompréhensible et fait perdre des points.
\end{enumerate}
\end{corrige}

\begin{corrige}[Exercice 3 — Vocabulaire : métiers et formation]
\begin{enumerate}
\item \textbf{workshop} — \eng{My brother is learning carpentry in a small workshop near the Mokolo market.} (atelier)
\item \textbf{trade} — \eng{Nursing is a noble trade that demands love for others.} (métier)
\item \textbf{qualification} — \eng{She obtained her qualification as an electrician after two years of training.} (diplôme, qualification)
\item \textbf{apprenticeship} — \eng{During his apprenticeship, the young plumber worked with an experienced master.} (apprentissage)
\item \textbf{skills} — \eng{Good communication skills are necessary in every profession.} (compétences)
\item \textbf{career} — \eng{He wants to make a career in journalism.} (carrière)
\end{enumerate}
\end{corrige}

\begin{corrige}[Exercice 4 — Conjugaison à compléter]
\begin{enumerate}
\item \textbf{watched} — \eng{When I was a child, I often watched my uncle repair radios.} Past simple : habitude passée révolue (\eng{when I was a child}).
\item \textbf{have developed} — \eng{Since then, I have developed a real passion for electronics.} Present perfect avec \eng{since then}.
\item \textbf{help} — \eng{Every Saturday, I help him in his shop in Bamenda.} Present simple : habitude actuelle (\eng{every Saturday}).
\item \textbf{am going to enter} — \eng{Next year, I am going to enter a technical college.} Projet d'avenir : \eng{be going to}.
\item \textbf{has taught} — \eng{So far, this experience has taught me discipline and hard work.} Present perfect avec \eng{so far} ; participe passé irrégulier : \eng{teach -- taught -- taught}.
\end{enumerate}
\end{corrige}

\begin{corrige}[Exercice 5 — Traduction]
\begin{enumerate}
\item \eng{Since my childhood, I have been interested in the nursing profession.} — Attention : \eng{since} + present perfect ; on dit \eng{to be intereste\textbf{d} in} (pas \eng{interesting in}).
\item \eng{First of all, this job requires patience; then, it demands good training.} — Connecteurs d'énumération : \eng{first of all}, \eng{then}.
\item \eng{My brother started an apprenticeship with a carpenter two years ago.} — \eng{Two years ago} impose le past simple ; « menuisier » = \eng{carpenter}.
\item \eng{However, many young people prefer computer jobs.} — \eng{However} suivi d'une virgule ; « les métiers de l'informatique » = \eng{computer jobs} ou \eng{jobs in computing}.
\item \eng{In conclusion, I am determined to share my passion with my classmates.} — « déterminé à » = \eng{determined to + base verbale}.
\end{enumerate}
\end{corrige}

\begin{corrige}[Exercice 6 — Relier les phrases avec un connecteur]
\begin{enumerate}
\item \eng{Training as a mechanic is long; however, it is very rewarding.} (opposition)
\item \eng{Many trades require practical abilities; for example, a welder must master his tools perfectly.} (exemple)
\item \eng{I love helping sick people; therefore, I have decided to become a nurse.} (conséquence)
\item \eng{Carpentry develops creativity; moreover, it offers good job opportunities in our towns.} (addition)
\item \eng{In conclusion, I have explained my project and my motivations, and I am ready to start my training.} (conclusion)
\end{enumerate}
\textbf{Point de méthode :} les connecteurs \eng{however}, \eng{therefore} et \eng{moreover} peuvent aussi être placés en début de la seconde phrase, suivis d'une virgule : \eng{Training as a mechanic is long. However, it is very rewarding.}
\end{corrige}

\begin{corrige}[Exercice 7 — Composition à trous]
\begin{textebox}[Corrigé — My interest in carpentry]
Since my childhood, I have been \textbf{interested} in carpentry. My uncle owns a small \textbf{workshop} in Bafoussam, and I often spent my holidays there watching him work. Two years ago, I started an \textbf{apprenticeship} with him. During this \textbf{training}, I learnt many useful \textbf{skills}, such as measuring wood and using machines safely. Carpentry is a \textbf{trade} that I respect deeply. After my \textbf{qualification}, I hope to open my own business and \textbf{share} my passion with young people of my quarter.
\end{textebox}
\textbf{Remarques :} \eng{to be interested in} (et non \eng{interesting}) ; \eng{an apprenticeship} (voyelle initiale, donc \eng{an}) ; \eng{to share something with somebody}.
\end{corrige}

\begin{corrige}[Exercice 8 — Corrige le texte d'un élève]
\begin{center}
\begin{tabularx}{\linewidth}{|X|X|}\hline
\rowcolor{popblueL} \textbf{Erreur} & \textbf{Correction} \\ \hline
I am \emph{interesting} in & I am \textbf{interested} in (on est « intéressé » : participe passé) \\ \hline
\emph{since five years} & \textbf{for five years} (\eng{since} + point de départ ; \eng{for} + durée) \\ \hline
I am interested\ldots\ since five years & \textbf{I have been interested} in\ldots\ for five years (present perfect avec durée) \\ \hline
My mother \emph{teach} & My mother \textbf{teaches} (3\textsuperscript{e} personne du singulier) \\ \hline
I admire \emph{she} & I admire \textbf{her} (pronom complément) \\ \hline
\emph{Firstable} & \textbf{First of all} (calque du français « d'abord » ; \eng{firstable} n'existe pas) \\ \hline
\emph{Morever} & \textbf{Moreover} (orthographe) \\ \hline
their \emph{knowledges} & their \textbf{knowledge} (\eng{knowledge} est indénombrable) \\ \hline
Last year, I \emph{have helped} & Last year, I \textbf{helped} (past simple avec \eng{last year}) \\ \hline
I will enter \emph{in} a college & I will \textbf{enter} a college (\eng{to enter} se construit sans préposition) \\ \hline
\emph{the teaching} is my dream job & \textbf{teaching} is my dream job (pas d'article devant un nom général) \\ \hline
\end{tabularx}
\end{center}
\textbf{Texte corrigé :} \eng{I have been interested in the job of teacher for five years. My mother teaches in a primary school in Yaoundé and I admire her very much. First of all, teachers form the future of the nation. Moreover, they share their knowledge with children. Last year, I helped my little brother with his homework every evening. I am going to pass my A Level next year, and after I will enter a teacher training college. In conclusion, teaching is my dream job.}
\end{corrige}

\begin{corrige}[Exercice 9 — Compréhension et langue — type Bac]
\textbf{A. Comprehension questions}
\begin{enumerate}
\item \eng{Her grandmother, who owned a small sewing shop in Deido, taught her the basics of tailoring.}
\item \eng{During her training course, she learnt how to take measurements, cut fabric and use an electric sewing machine.} (deux de ces trois éléments suffisent)
\item \eng{Some classmates laughed at her dream because they believed that only office jobs were respectable.}
\item \eng{After her A Levels, she is going to attend a professional school of fashion design; then she wants to open her own workshop and train other young girls.}
\end{enumerate}
\textbf{B. Language work}
\begin{enumerate}
\item Connecteurs d'addition : \eng{moreover}, \eng{not only\ldots\ but also} ; connecteur d'opposition : \eng{however}.
\item \eng{She has been fascinated} est au \cle{present perfect} (passif) : il exprime un état commencé dans le passé (\eng{since she was ten}) et qui continue aujourd'hui.
\item Voix passive : \eng{She was taught how to use a needle by her grandmother.}
\end{enumerate}
\textbf{Barème indicatif (20 pts) :} A : 4 $\times$ 2 pts = 8 pts (phrase complète et correcte) ; B : 6 pts (2 pts par item) ; qualité de la langue : 6 pts.
\textbf{Points de vigilance :} répondre par des phrases complètes, pas par des mots isolés ; reprendre le temps de la question ; à la voix passive, ne pas oublier l'auxiliaire \eng{be} conjugué + participe passé (\eng{was taught}).
\end{corrige}

\begin{corrige}[Exercice 10 — Remettre la composition dans l'ordre — type Bac]
\textbf{Ordre logique : C -- D -- B -- E -- A.}
\begin{itemize}
\item \textbf{C} = introduction : présentation du sujet (\eng{my greatest wish is to become a nurse}) et annonce du plan (\eng{I will explain the origin of this interest and my plans}).
\item \textbf{D} = premier paragraphe du développement : \eng{First of all} marque le début de l'argumentation (l'origine de l'intérêt).
\item \textbf{B} = deuxième paragraphe : \eng{Moreover} ajoute un argument.
\item \textbf{E} = troisième paragraphe : \eng{Therefore} tire la conséquence (le projet de formation).
\item \textbf{A} = conclusion : \eng{In conclusion} clôt la composition.
\end{itemize}
\textbf{Barème indicatif (10 pts) :} ordre correct 5 pts ; justification par les connecteurs et indices de sens 5 pts.
\textbf{Point de vigilance :} repérer d'abord les connecteurs d'ouverture (\eng{first of all}) et de clôture (\eng{in conclusion}), puis la chaîne logique addition $\rightarrow$ conséquence.
\end{corrige}

\begin{corrige}[Exercice 11 — Sujet de composition — type Bac]
\textbf{Proposition de rédaction modèle (environ 180 mots) :}
\begin{textebox}[Model composition]
Many students of my school dream of becoming lawyers or engineers, but my greatest wish is to become an electrician. I will explain the origin of this interest and my plans for the future.

First of all, I discovered this trade when I was thirteen. Our house in Bamenda had a serious electrical problem, and a young technician repaired everything in one afternoon. I was fascinated by his skill and his confidence. Since that day, I have been interested in electricity.

Moreover, this profession requires precision, patience and responsibility, because a small mistake can be dangerous. However, it is also very rewarding: an electrician brings light and comfort to families. For example, during the holidays, I helped my uncle install lamps in our quarter, and I felt very proud.

Therefore, after my A Levels, I am going to follow a two-year training course in a technical college in order to obtain my qualification.

In conclusion, electricity is the profession of my heart, and I am determined to share this interest with my classmates by creating a science club in my school. (180 words)
\end{textebox}
\textbf{Barème indicatif (20 pts) :} respect du sujet et du nombre de mots 4 pts ; structure en trois parties 4 pts ; connecteurs (au moins quatre) 4 pts ; correction des temps verbaux 4 pts ; vocabulaire et orthographe 4 pts.
\textbf{Points de vigilance :}
\begin{itemize}
\item Past simple pour l'origine (\eng{I discovered}, \eng{I was fascinated}) ; present perfect avec \eng{since} ; \eng{be going to} pour les projets.
\item Ne pas écrire \eng{I am interesting in} mais \eng{I am interested in}.
\item Indiquer le nombre de mots à la fin et soigner la présentation en paragraphes.
\end{itemize}
\end{corrige}

\newpage

\coursec{Fiche bilan}

\begin{popfigure}
\begin{center}\tcbox[colback=popblue,colframe=popblue,coltext=white,arc=9pt,boxsep=4pt,left=6pt,right=6pt]{\bfseries\small Writing: a trade, job or profession}\end{center}
\vspace{-2pt}
\begin{tcbraster}[raster columns=3,raster equal height=rows,raster column skip=6pt,raster row skip=6pt,enhanced,blankest]
\begin{tcolorbox}[enhanced,colback=white,colframe=poporange,boxrule=0.9pt,arc=3pt,title={Vocabulary trades \& jobs},fonttitle=\bfseries\scriptsize,coltitle=white,colbacktitle=poporange,left=4pt,right=4pt,top=2pt,bottom=2pt,boxsep=1pt,toptitle=1.5pt,bottomtitle=1.5pt,halign=left]
\begin{itemize}[nosep,leftmargin=*,itemsep=1pt,font=\scriptsize]
\item occupation, trade, craft
\item apprentice, trainee
\end{itemize}
\end{tcolorbox}
\begin{tcolorbox}[enhanced,colback=white,colframe=popgreen,boxrule=0.9pt,arc=3pt,title={Training \& skills},fonttitle=\bfseries\scriptsize,coltitle=white,colbacktitle=popgreen,left=4pt,right=4pt,top=2pt,bottom=2pt,boxsep=1pt,toptitle=1.5pt,bottomtitle=1.5pt,halign=left]
\begin{itemize}[nosep,leftmargin=*,itemsep=1pt,font=\scriptsize]
\item vocational school
\item to learn a trade
\end{itemize}
\end{tcolorbox}
\begin{tcolorbox}[enhanced,colback=white,colframe=poppurple,boxrule=0.9pt,arc=3pt,title={Structure 3 parts},fonttitle=\bfseries\scriptsize,coltitle=white,colbacktitle=poppurple,left=4pt,right=4pt,top=2pt,bottom=2pt,boxsep=1pt,toptitle=1.5pt,bottomtitle=1.5pt,halign=left]
\begin{itemize}[nosep,leftmargin=*,itemsep=1pt,font=\scriptsize]
\item introduction
\item body (2--3 §)
\item conclusion
\end{itemize}
\end{tcolorbox}
\begin{tcolorbox}[enhanced,colback=white,colframe=poppink,boxrule=0.9pt,arc=3pt,title={Link words},fonttitle=\bfseries\scriptsize,coltitle=white,colbacktitle=poppink,left=4pt,right=4pt,top=2pt,bottom=2pt,boxsep=1pt,toptitle=1.5pt,bottomtitle=1.5pt,halign=left]
\begin{itemize}[nosep,leftmargin=*,itemsep=1pt,font=\scriptsize]
\item first, then, finally
\item because, so that
\end{itemize}
\end{tcolorbox}
\begin{tcolorbox}[enhanced,colback=white,colframe=popgold,boxrule=0.9pt,arc=3pt,title={Verb tenses},fonttitle=\bfseries\scriptsize,coltitle=white,colbacktitle=popgold,left=4pt,right=4pt,top=2pt,bottom=2pt,boxsep=1pt,toptitle=1.5pt,bottomtitle=1.5pt,halign=left]
\begin{itemize}[nosep,leftmargin=*,itemsep=1pt,font=\scriptsize]
\item present simple
\item present perfect
\item be going to / will
\end{itemize}
\end{tcolorbox}
\begin{tcolorbox}[enhanced,colback=white,colframe=popblue!70,boxrule=0.9pt,arc=3pt,title={Method 5 steps},fonttitle=\bfseries\scriptsize,coltitle=white,colbacktitle=popblue!70,left=4pt,right=4pt,top=2pt,bottom=2pt,boxsep=1pt,toptitle=1.5pt,bottomtitle=1.5pt,halign=left]
\begin{itemize}[nosep,leftmargin=*,itemsep=1pt,font=\scriptsize]
\item plan $\to$ draft
\item check \& correct
\end{itemize}
\end{tcolorbox}
\begin{tcolorbox}[enhanced,colback=white,colframe=popdark,boxrule=0.9pt,arc=3pt,title={Pitfalls to avoid},fonttitle=\bfseries\scriptsize,coltitle=white,colbacktitle=popdark,left=4pt,right=4pt,top=2pt,bottom=2pt,boxsep=1pt,toptitle=1.5pt,bottomtitle=1.5pt,halign=left]
\begin{itemize}[nosep,leftmargin=*,itemsep=1pt,font=\scriptsize]
\item faux amis
\item word order
\end{itemize}
\end{tcolorbox}
\end{tcbraster}

\legende{Carte mentale de la leçon : rédiger une composition sur un métier ou une profession.}
\end{popfigure}

\begin{bilan}
\textbf{1. Lexique clé (à connaître par cœur)}
\begin{itemize}
  \item \eng{a trade / a craft} : un métier manuel ; \eng{an occupation} : une profession ; \eng{a job} : un emploi.
  \item \eng{a tailor, a carpenter, a mechanic, a hairdresser, a nurse, an engineer} : tailleur, menuisier, mécanicien, coiffeur, infirmier, ingénieur.
  \item \eng{an apprentice / an apprenticeship} : un apprenti / un apprentissage ; \eng{a trainee} : un stagiaire.
  \item \eng{to learn a trade} : apprendre un métier ; \eng{to be trained as} : être formé comme ; \eng{to earn a living} : gagner sa vie.
  \item \eng{vocational training} : formation professionnelle ; \eng{a workshop} : un atelier ; \eng{a skill} : une compétence.
  \item \eng{to be interested in / to be keen on / to be passionate about} : s'intéresser à / être passionné par.
\end{itemize}

\textbf{2. Grammaire et connecteurs}
\begin{center}
\begin{tabularx}{\linewidth}{|X|X|X|}\hline
\rowcolor{popblueL} Règle & Exemple anglais & Traduction \\ \hline
Présent simple (habitudes, goûts) & \eng{I work in my uncle's workshop every day.} & Je travaille chaque jour dans l'atelier de mon oncle. \\ \hline
Present perfect (expérience) & \eng{I have learnt a lot since last year.} & J'ai beaucoup appris depuis l'année dernière. \\ \hline
\eng{be going to} (projet) & \eng{I am going to open my own shop.} & Je vais ouvrir ma propre boutique. \\ \hline
\eng{want / hope / plan + to + BV} & \eng{She hopes to become a nurse.} & Elle espère devenir infirmière. \\ \hline
\eng{since / for} (durée) & \eng{for two years / since 2023} & depuis deux ans / depuis 2023 \\ \hline
\end{tabularx}
\end{center}
Connecteurs : \eng{first of all, then, moreover, however, because, so that, as a result, to conclude}.

\textbf{3. Méthode express (5 étapes)}
\begin{enumerate}
  \item \textbf{Analyser} le sujet et noter les mots-clés.
  \item \textbf{Brainstormer} : lister idées et vocabulaire (métier, raisons, formation, projets).
  \item \textbf{Planifier} : introduction (présenter le métier) $\to$ développement (raisons, formation, qualités) $\to$ conclusion (projet d'avenir).
  \item \textbf{Rédiger} avec connecteurs et temps corrects.
  \item \textbf{Relire} : accords, 3\textsuperscript{e} personne (\eng{-s}), orthographe, ponctuation.
\end{enumerate}
\end{bilan}

\begin{aretenir}[Checklist avant le Bac]
\begin{itemize}
  \item $\square$ Je connais le vocabulaire des métiers et de la formation professionnelle.
  \item $\square$ Je sais distinguer \eng{job}, \eng{trade}, \eng{occupation} et \eng{profession}.
  \item $\square$ J'évite le faux ami \eng{formation} $\neq$ \eng{training}.
  \item $\square$ Ma composition a une introduction, un développement (2--3 paragraphes) et une conclusion.
  \item $\square$ J'utilise au moins quatre connecteurs logiques variés.
  \item $\square$ J'emploie le présent simple pour les habitudes et \eng{be going to} pour mes projets.
  \item $\square$ Je n'oublie pas le \eng{-s} à la 3\textsuperscript{e} personne du singulier (\eng{he works}).
  \item $\square$ J'utilise \eng{since} + point de départ et \eng{for} + durée correctement.
  \item $\square$ Je construis \eng{to be interested in + V-ing} (\eng{I am interested in sewing}).
  \item $\square$ Je relis ma copie : orthographe, ponctuation, majuscules.
  \item $\square$ Je respecte le nombre de mots demandé et je soigne ma présentation.
\end{itemize}
\end{aretenir}

\findecours

\end{document}
